% =====================================================================
% Kapitel: Verteilte Ausfuehrung auf dem Rechencluster
% Autor: Mier
%
% EINBINDEN: % =====================================================================
% Kapitel: Verteilte Ausfuehrung auf dem Rechencluster
% Autor: Mier
%
% EINBINDEN: % =====================================================================
% Kapitel: Verteilte Ausfuehrung auf dem Rechencluster
% Autor: Mier
%
% EINBINDEN: % =====================================================================
% Kapitel: Verteilte Ausfuehrung auf dem Rechencluster
% Autor: Mier
%
% EINBINDEN: \input{Kapitel_Cluster_Mier} im Master.
% Benoetigte Pakete: amsmath, listings, graphicx, booktabs.
% Fuer saubere Optik im Master: fontenc, babel[ngerman], microtype.
% Abbildungen: latexDocs/figures/cluster_arch.png, mc_convergence.png
%
% STANDALONE-VORSCHAU:
%   \documentclass[12pt]{article}
%   \usepackage[utf8]{inputenc}\usepackage[T1]{fontenc}
%   \usepackage[ngerman]{babel}\usepackage{microtype}
%   \usepackage{amsmath}\usepackage{listings}
%   \usepackage{graphicx}\usepackage{booktabs}
%   \begin{document}
%   \input{Kapitel_Cluster_Mier}
%   \end{document}
% =====================================================================

\lstdefinestyle{clusterlst}{
    basicstyle=\ttfamily\small,
    showstringspaces=false,
    breaklines=true,
    numbers=left,
    numberstyle=\tiny,
    frame=single,
    tabsize=2
}

%======================================================================
\section{Verteilte Ausführung auf dem Rechencluster}
\label{sec:cluster}

\subsection{Einordnung und Aufgabenstellung}
\label{subsec:cluster-einordnung}

Die Strahlverfolgung ist der teuerste Schritt der gesamten Pipeline. Für eine
brauchbare Impulsantwort müssen Hunderttausende bis Millionen einzelner
Schallstrahlen durch die Raumgeometrie verfolgt werden, und jeder dieser Strahlen
verlangt mehrere Schnitt- und Reflexionsberechnungen. Der entscheidende Vorteil
dabei: Kein Strahl hängt vom Ergebnis eines anderen ab. Das Problem ist damit
\emph{embarrassingly parallel} und lässt sich ohne Kommunikation zwischen den
Teilrechnungen auf beliebig viele Rechner verteilen.

Dieses Kapitel dokumentiert die Verteilungsschicht des Systems. Sie startet die
Rust-Simulation gleichzeitig auf allen acht Knoten des Labor-Clusters, wartet auf
deren Abschluss, sammelt die Teilergebnisse ein und fügt sie zu einer einzigen,
statistisch dichteren Impulsantwort zusammen. Aus Sicht der übrigen Pipeline
ändert sich dabei nichts: Am Ende steht wieder eine Datei im gewohnten Schema,
die die Signalverarbeitung unverändert weiterverarbeiten kann.

Zur Abgrenzung von den Nachbarkapiteln gehören zwei Punkte. Erstens ist die
\emph{prozessinterne} Nebenläufigkeit der Rust-Engine, also die Verteilung auf
die Kerne eines einzelnen Rechners mittels \texttt{rayon}, an anderer Stelle
beschrieben und bleibt hier außen vor. Sie parallelisiert innerhalb eines
Prozesses; dieses Kapitel parallelisiert über Prozess- und Rechnergrenzen hinweg.
Zweitens geht es nicht um den einzelnen Subprozess-Aufruf aus der Python-Pipeline,
sondern um die Ebene darüber: die gleichzeitige Ausführung auf mehreren physischen
Knoten und die anschließende Aggregation ihrer Ergebnisse.

\subsection{Theoretischer Hintergrund}
\label{subsec:cluster-theorie}

\subsubsection{Monte-Carlo-Strahlverfolgung}
\label{subsubsec:mc}

Die Engine schätzt die Impulsantwort eines Raumes stochastisch. Sie verschickt
zufällig gerichtete Strahlen von der Schallquelle, verfolgt ihre Reflexionen und
registriert, wann und mit welchem Restdruck sie den Empfänger erreichen. Jeder
einzelne Strahl ist eine Zufallsstichprobe; erst die Summe vieler Strahlen ergibt
ein aussagekräftiges Bild der Raumakustik. Solche Monte-Carlo-Verfahren haben eine
charakteristische Eigenschaft: Ihr Schätzfehler sinkt mit der Wurzel der
Stichprobengröße. Verdoppelt man die Zahl der Strahlen, halbiert sich der Fehler
also nicht, sondern sinkt nur um den Faktor $\sqrt{2}$. Genau diese langsame
Konvergenz macht große Strahlzahlen nötig und motiviert die Verteilung auf mehrere
Rechner.

\subsubsection{Zwei Wege der Verteilung}
\label{subsubsec:strategien}

Ein Monte-Carlo-Verfahren lässt sich auf zwei grundsätzlich verschiedene Arten
verteilen. Tabelle~\ref{tab:strategien} stellt sie gegenüber.

\begin{table}[htbp]
    \centering
    \caption{Zwei Strategien zur Verteilung der Strahlverfolgung.}
    \label{tab:strategien}
    \begin{tabular}{p{3.2cm} p{5.2cm} p{4.2cm}}
        \toprule
        & \textbf{Work-Partitioning} & \textbf{Replikat-Parallelismus} \\
        \midrule
        Prinzip & Gesamtzahl der Strahlen wird in disjunkte Anteile zerlegt & Jeder Knoten rechnet denselben Auftrag mit eigenem Zufallsstrom \\
        Koordination & Zuteilung der Anteile nötig & keine, jeder Knoten ist autark \\
        Zufallszahlen & ein globaler Strom, aufgeteilt & unabhängige Ströme pro Knoten \\
        Robustheit & Ausfall trifft einen festen Anteil & Ausfall kostet nur Stichprobengröße \\
        \bottomrule
    \end{tabular}
\end{table}

Das umgesetzte System verwendet den zweiten Weg. Alle acht Knoten starten die
Simulation mit derselben Konfigurationsdatei \texttt{input\_config.json} und
werfen dennoch jeweils andere Strahlen, weil jeder Prozess seinen eigenen
Zufallsstrom besitzt. Der Ansatz ist bewusst schlicht gewählt: Er braucht keine
zentrale Zuteilung und keinen Austausch zwischen den Knoten, sondern nur, dass am
Ende alle Ergebnisse eingesammelt werden.

\subsubsection{Unabhängigkeit der Zufallsströme}
\label{subsubsec:rng}

Die Korrektheit des Replikat-Ansatzes steht und fällt mit einer Bedingung: Die
Zufallszahlen auf den Knoten müssen unabhängig sein. Würde jeder Knoten dieselbe
Folge zufälliger Richtungen erzeugen, so lieferten alle acht exakt dasselbe
Ergebnis, und das Zusammenführen brächte keinen Erkenntnisgewinn, sondern nur
achtfach dieselben Strahlen. Die Rust-Engine schließt das aus, indem sie ihren
Zufallsstrom pro Thread bzw. Prozess frisch aus der Systementropie zieht:

\begin{lstlisting}[style=clusterlst, language=Rust, caption={Initialisierung des Zufallsstroms pro Prozess (\texttt{simulation.rs}).}]
let mut rng = rand::rng();                       // Omnidirectional, random rays
let random_angle: f32 = rng.random_range(0.0..(2.0 * PI));
let direction = Vec2::new(random_angle.cos(), random_angle.sin());
\end{lstlisting}

Weil jeder der acht Prozesse auf seinem Knoten einen anderen Startzustand erhält,
erzeugt jeder eine andere Stichprobe. Acht Knoten mit je einer Million Strahlen
liefern damit acht Millionen \emph{unabhängige} Stichproben derselben Raumakustik.

\subsubsection{Statistische Konvergenz}
\label{subsubsec:konvergenz}

Statistisch verhält sich der verteilte Lauf wie eine einzige, achtfach größere
Simulation. Für einen Monte-Carlo-Schätzer aus $N$ unabhängigen Teilstichproben
gleicher Größe sinkt die Standardabweichung des Mittelwerts mit
\[
    \sigma_N \;=\; \frac{\sigma_1}{\sqrt{N}} .
\]
Acht Knoten reduzieren die Streuung der geschätzten Impulsantwort demnach um den
Faktor $\sqrt{8} \approx 2{,}83$ gegenüber einem einzelnen Knoten.
Abbildung~\ref{fig:mc-convergence} bestätigt dieses Verhalten in einem numerischen
Experiment: Aufgetragen ist der mittlere quadratische Fehler der geschätzten
Impulsantwort über der Knotenzahl. Die Messpunkte folgen sauber der
$1/\sqrt{N}$-Referenzkurve; von einem auf acht Knoten sinkt der Fehler um den
Faktor $2{,}78$, praktisch exakt der theoretisch erwartete Wert.

\begin{figure}[htbp]
    \centering
    \includegraphics[width=0.8\textwidth]{figures/mc_convergence}
    \caption{Mittlerer quadratischer Fehler der geschätzten Impulsantwort über der
    Knotenzahl $N$ (je $5\cdot10^4$ Strahlen, doppelt logarithmisch). Die
    Messpunkte folgen der $1/\sqrt{N}$-Referenz; acht Knoten senken den Fehler um
    den Faktor $2{,}78 \approx \sqrt{8}$.}
    \label{fig:mc-convergence}
\end{figure}

Bemerkenswert ist die Kehrseite dieser Kurve: Um den Fehler zu halbieren, braucht
es die vierfache Knotenzahl. Der Nutzen zusätzlicher Knoten nimmt also ab, was bei
der Wahl der Clustergröße zu berücksichtigen ist. Für die hier angestrebte
Genauigkeit sind acht Knoten ein sinnvoller Kompromiss zwischen Aufwand und
Qualität.

\subsection{Architektur des Clusters}
\label{subsec:cluster-arch}

Das System folgt einem \emph{Manager-Worker}-Muster
(Abbildung~\ref{fig:cluster-arch}). Ein Manager-Knoten koordiniert den Ablauf,
acht Worker-Knoten führen die eigentliche Simulation aus. Der Manager selbst
rechnet nicht mit, sondern startet die Worker, wartet auf sie und führt am Ende
die Ergebnisse zusammen.

\begin{figure}[htbp]
    \centering
    \includegraphics[width=0.82\textwidth]{figures/cluster_arch}
    \caption{Manager-Worker-Architektur. Der Manager startet die acht Knoten
    parallel per \texttt{ssh}, holt nach Abschluss jede Ergebnisdatei per
    \texttt{scp} zurück und führt sie mit \texttt{merge\_results.py} zusammen.}
    \label{fig:cluster-arch}
\end{figure}

Die Kommunikation läuft vollständig über Standard-Werkzeuge des Betriebssystems.
\texttt{ssh} startet die Prozesse aus der Ferne, \texttt{scp} überträgt die
Ergebnisdateien zurück. Ein eigenes Netzwerkprotokoll oder ein
Message-Passing-Framework kommt nicht zum Einsatz. Für einen Stapelbetrieb, bei
dem jeder Knoten genau einmal startet, eine Weile rechnet und ein einziges
Ergebnis zurückliefert, wäre solcher Aufwand nicht gerechtfertigt: Es gibt keine
laufende Zustandssynchronisation und keinen Nachrichtenaustausch während der
Rechnung, den ein Framework verwalten müsste.

Damit das funktioniert, müssen drei Voraussetzungen erfüllt sein. Auf jedem Knoten
liegt dieselbe vorkompilierte Binärdatei \texttt{rust\_service} samt der
Konfiguration \texttt{input\_config.json}. Der Manager kann sich per SSH ohne
interaktive Passworteingabe anmelden, was eine Schlüssel-basierte
Authentifizierung voraussetzt. Und die Knoten sind über feste Adressen im
Labornetz erreichbar.

\subsection{Ablaufsteuerung: \texttt{run\_cluster.sh}}
\label{subsec:cluster-run}

Das Skript \texttt{run\_cluster.sh} steuert den verteilten Lauf in drei Phasen:
paralleles Starten, Warten und Einsammeln.

\begin{lstlisting}[style=clusterlst, language=bash, caption={Verteilter Start und Einsammeln der Ergebnisse (gekürzt).}]
# Phase 1: alle Knoten gleichzeitig starten (Hintergrundprozesse)
ssh node1@139.6.65.29 "./rust_service --config-file input_config.json" &
ssh node2@139.6.65.18 "./rust_service --config-file input_config.json" &
# ... node3 bis node8 analog ...

wait                       # Phase 2: auf Abschluss ALLER Knoten warten

# Phase 3: Ergebnis jedes Knotens zum Manager kopieren
scp node1@139.6.65.29:~/ir_output.json ~/ir_node1.json
scp node2@139.6.65.18:~/ir_output.json ~/ir_node2.json
# ... node3 bis node8 analog ...
\end{lstlisting}

Zwei Eigenschaften der Shell nutzt der Ablauf gezielt. Das kaufmännische Und
(\texttt{\&}) am Zeilenende schickt jeden \texttt{ssh}-Aufruf in den Hintergrund.
Die acht Simulationen laufen dadurch nahezu gleichzeitig an, statt nacheinander
abgearbeitet zu werden. Das darauffolgende \texttt{wait} blockiert, bis
\emph{alle} Hintergrundprozesse beendet sind, und bildet so eine
Synchronisationsbarriere zwischen der Rechen- und der Einsammel-Phase. Erst wenn
der letzte Knoten fertig ist, beginnt \texttt{scp}, von jedem Knoten dessen
\texttt{ir\_output.json} zu holen und unter einem knotenspezifischen Namen auf dem
Manager abzulegen.

Diese Barriere hat eine wichtige Konsequenz für die Laufzeit: Der Gesamtlauf ist
so langsam wie der \emph{langsamste} Knoten. Solange die Hardware homogen ist,
fällt das kaum ins Gewicht; auf ungleichen Knoten würde ein einzelner langsamer
Rechner den ganzen Lauf ausbremsen, während die übrigen sieben bereits fertig
warten. Dieses Straggler-Problem ist der zentrale Grund, warum ein statischer
Replikat-Ansatz mit gleich großen Aufträgen homogene Hardware voraussetzt.

\subsection{Aggregation der Teilergebnisse: \texttt{merge\_results.py}}
\label{subsec:cluster-merge}

Sobald alle Teilergebnisse auf dem Manager liegen, führt \texttt{merge\_results.py}
sie zu einer einzigen Impulsantwort zusammen. Jeder Knoten liefert dasselbe
Ausgabeschema, nämlich Metadaten sowie zwei positionsgleiche Listen aus
Verzögerungen und Drücken. Die Aggregation besteht deshalb im Kern aus einer
schlichten Verkettung dieser Listen:

\begin{lstlisting}[style=clusterlst, language=Python, caption={Zusammenführen der Knoten-Ergebnisse, gekürzt (\texttt{merge\_results.py}).}]
all_delays, all_pressures = [], []
for f in sorted(glob.glob("/home/manager/ir_node*.json")):
    with open(f) as file:
        data = json.load(file)
        all_delays    += data["hits"]["delays_seconds"]
        all_pressures += data["hits"]["pressures"]

merged = {
    "metadata": {"sample_rate": 44100, "rays_cast": 1000000 * 8,
                 "rays_received": len(all_delays), "room_name": "MVP_Test_Room_1"},
    "hits": {"delays_seconds": all_delays, "pressures": all_pressures},
}
json.dump(merged, open("/home/manager/ir_merged.json", "w"), indent=2)
\end{lstlisting}

Dass die Verkettung hier die richtige Operation ist, folgt direkt aus dem
Replikat-Modell. Weil die einzelnen Läufe unabhängige Stichproben desselben Raums
sind, ist die Vereinigung aller Treffer statistisch identisch mit dem Ergebnis
eines einzigen Laufs mit acht Millionen Strahlen. Es findet keine Mittelung und
keine Gewichtung statt, sondern nur ein Aneinanderhängen. Das Feld
\texttt{rays\_cast} ist im Skript fest auf $8 \cdot 10^{6}$ gesetzt; solange der in
Abschnitt~\ref{subsec:cluster-grenzen} beschriebene Kopierfehler bei Knoten~6
besteht, fließen tatsächlich nur sieben Knoten in die Aggregation ein, und dieser
Metadatenwert überschätzt die reale Strahlzahl. \texttt{rays\_received} zählt
dagegen die tatsächlich empfangenen Treffer. Die
zusammengeführte Datei trägt exakt dasselbe Schema wie die Ausgabe eines einzelnen
Knotens und ist damit ohne weitere Anpassung für die DSP-Pipeline verwendbar.

\subsection{Datenfluss im Überblick}
\label{subsec:cluster-datenfluss}

Der gesamte verteilte Ablauf lässt sich in fünf Stufen zusammenfassen:

\begin{enumerate}
    \item Der Manager startet per SSH auf jedem der acht Knoten
    \texttt{rust\_service} mit identischer Konfiguration, alle parallel.
    \item Jeder Knoten simuliert unabhängig eine Million Strahlen und schreibt
    lokal seine \texttt{ir\_output.json}.
    \item \texttt{wait} synchronisiert: Die Einsammel-Phase beginnt erst, wenn der
    letzte Knoten fertig ist.
    \item \texttt{scp} überträgt jede \texttt{ir\_output.json} als
    \texttt{ir\_node$k$.json} zum Manager.
    \item \texttt{merge\_results.py} verkettet alle Trefferlisten zu einer
    \texttt{ir\_merged.json}, die anschließend die DSP-Pipeline verarbeitet.
\end{enumerate}

\subsection{Leistungsbetrachtung}
\label{subsec:cluster-leistung}

\subsubsection{Erwarteter Speedup}
\label{subsubsec:speedup}

Bei der Leistung sind zwei Sichtweisen zu trennen. In der umgesetzten
Replikat-Betriebsart rechnet jeder Knoten dieselbe Strahlzahl, sodass acht Knoten
in nahezu derselben Wanduhrzeit die achtfache Stichprobe liefern. Der Gewinn ist
hier keine kürzere Laufzeit, sondern die höhere Genauigkeit aus
Abbildung~\ref{fig:mc-convergence}; man spricht vom schwachen Skalieren im Sinne
von Gustafson, bei dem mit der Rechnerzahl die bearbeitete Problemgröße wächst,
nicht die Geschwindigkeit einer festen Aufgabe. Will man dagegen eine feste
Gesamtstrahlzahl in kürzerer Zeit erreichen, so verteilt man diese auf die Knoten,
sodass jeder nur $8\cdot10^{6}/k$ Strahlen rechnet. Erst in dieser
Partitionierungs-Betriebsart greift das Gesetz von Amdahl. Bezeichnet $p$ den parallelen
Anteil der Arbeit, hier die reine Strahlverfolgung, und $1-p$ den seriellen Rest
aus Prozessstart, Ergebnistransfer und Aggregation, so beträgt der theoretische
Speedup
\[
    S(k) \;=\; \frac{1}{(1-p) + \dfrac{p}{k}} .
\]
Da die Strahlverfolgung die Laufzeit klar dominiert, liegt $p$ nahe bei eins, und
$S(k)$ steigt bis $k=8$ nahezu linear. In der Praxis bleibt der reale Speedup
darunter, und zwar aus drei Gründen: dem bereits erwähnten Straggler-Effekt (der
Lauf ist so schnell wie der langsamste Knoten), dem Overhead für den
SSH-Verbindungsaufbau sowie der Übertragung und Zusammenführung der
Ergebnisdateien.

\subsubsection{Kommunikations- und I/O-Overhead}
\label{subsubsec:overhead}

Der Overhead dieses Ansatzes ist klein, aber nicht null. Pro Knoten fallen ein
SSH-Verbindungsaufbau und eine \texttt{scp}-Übertragung an. Übertragen wird jeweils
eine JSON-Datei, deren Größe mit der Zahl der empfangenen Treffer wächst; bei
einer Million geworfener Strahlen liegt sie im Bereich einiger Megabyte. Da diese
acht Übertragungen erst \emph{nach} der Rechenphase stattfinden und nicht mit ihr
überlappen, addieren sie sich voll zur Gesamtlaufzeit. Bei deutlich mehr Knoten
oder größeren Ergebnisdateien würde dieser sequenzielle Einsammel-Schritt zum
Flaschenhals; eine Überlappung von Rechnung und Transfer oder ein kompakteres
Ausgabeformat als das gut lesbare, aber platzintensive JSON wären dann die
naheliegenden Gegenmaßnahmen.

\subsubsection{Messung am Cluster}
\label{subsubsec:messung}

Die konkreten Laufzeiten hängen von der Laborhardware ab und sind am Cluster zu
messen. Um einen aussagekräftigen Speedup zu erhalten, wird in der
Partitionierungs-Betriebsart gemessen: In \texttt{input\_config.json} wird
\texttt{rays\_to\_cast} auf $8\cdot10^{6}/k$ gesetzt, sodass die Gesamtarbeit über
alle $k$ Knoten konstant bleibt. Dann wird \texttt{run\_cluster.sh} für
verschiedene Knotenzahlen mit der Shell-eigenen Zeitmessung ausgeführt:

\begin{lstlisting}[style=clusterlst, language=bash, caption={Messvorschrift für die Laufzeit am Cluster.}]
time ./run_cluster.sh        # Wanduhrzeit fuer den kompletten verteilten Lauf
\end{lstlisting}

Tabelle~\ref{tab:speedup} nimmt die Messwerte auf. Als Referenz dient der Lauf mit
einem einzelnen Knoten; der Speedup ergibt sich als Verhältnis der Wanduhrzeiten.

\begin{table}[htbp]
    \centering
    \caption{Wanduhrzeit und Speedup des verteilten Laufs (am Cluster zu messen).}
    \label{tab:speedup}
    \begin{tabular}{r r r}
        \toprule
        Knoten $k$ & Wanduhrzeit [s] & Speedup $S$ \\
        \midrule
        1 & \dots & $1{,}00$ \\
        2 & \dots & \dots \\
        4 & \dots & \dots \\
        8 & \dots & \dots \\
        \bottomrule
    \end{tabular}
\end{table}

\subsection{Fehlertoleranz und Sicherheit}
\label{subsec:cluster-fehler}

Zwei Querschnittsthemen verdienen eigene Betrachtung. Beim \textbf{Sicherheits}aspekt
ist die Schlüssel-basierte SSH-Authentifizierung zentral: Der Manager benötigt
passwortlosen Zugriff auf jeden Knoten, was über hinterlegte öffentliche Schlüssel
gelöst ist. Der private Schlüssel des Managers ist damit das entscheidende
Geheimnis; wer ihn besitzt, kann auf allen Knoten Code ausführen. In einem
abgeschlossenen Labornetz ist das vertretbar, für einen offeneren Betrieb wäre eine
engere Rechtevergabe angebracht.

Beim Thema \textbf{Fehlertoleranz} zeigt der aktuelle Stand seine deutlichste
Schwäche, die im folgenden Abschnitt zusammen mit den übrigen Grenzen diskutiert
wird. Der Kern des Problems: Ein ausgefallener Knoten wird nicht erkannt, sondern
schlicht übergangen, sodass die Stichprobengröße unbemerkt sinkt.

\subsection{Grenzen des aktuellen Standes und Ausblick}
\label{subsec:cluster-grenzen}

Die Verteilungsschicht ist als funktionsfähiger Prototyp umgesetzt. Für einen
robusten Betrieb bestehen mehrere klar benannte Schwachstellen, die zugleich die
nächsten Entwicklungsschritte vorgeben.

\begin{itemize}
    \item \textbf{Fest verdrahtete Konfiguration.} IP-Adressen, Knotennamen, die
    Knotenzahl und die Ablagepfade stehen direkt im Skript. Ein Umzug auf andere
    Hardware oder eine andere Knotenzahl erfordert Änderungen am Quelltext. Eine
    Auslagerung dieser Werte in eine Konfigurationsdatei wäre der naheliegende
    erste Schritt.

    \item \textbf{Kopierfehler bei Knoten~6.} Im Einsammel-Schritt wird das
    Ergebnis von Knoten~6 versehentlich nach \texttt{ir\_node5.json} kopiert und
    überschreibt damit das Ergebnis von Knoten~5. Effektiv fließen so nur sieben
    statt acht unabhängige Stichproben in die Aggregation ein. Der Fehler ist ein
    Tippfehler im Zielnamen und leicht zu beheben, führt aber vor Augen, wie
    fehleranfällig die von Hand ausgeschriebene Knotenliste ist -- ein weiteres
    Argument für die genannte Auslagerung in eine Schleife über eine
    Konfigurationsliste.

    \item \textbf{Keine Fehlererkennung.} Fällt ein Knoten aus oder liefert er kein
    Ergebnis, schlägt der zugehörige \texttt{scp}-Aufruf fehl, der Gesamtlauf setzt
    sich aber fort und aggregiert stillschweigend über die vorhandenen Dateien. Ein
    Knotenausfall bleibt damit unbemerkt und verfälscht die Stichprobengröße. Eine
    Prüfung der Rückgabecodes und eine Warnung bei fehlenden Ergebnissen wären hier
    angebracht.

    \item \textbf{Manuelle Kopplung an die DSP-Pipeline.} Die Aggregation erzeugt
    \texttt{ir\_merged.json}, die Python-DSP-Pipeline liest jedoch
    \texttt{ir\_output.json}. Der Übergang zwischen verteiltem Lauf und Faltung ist
    derzeit ein manueller Schritt. Eine saubere Integration würde den Cluster-Lauf
    direkt aus der Orchestrierung heraus anstoßen und die zusammengeführte Datei
    automatisch weiterreichen.

    \item \textbf{Statische Verteilung.} Der Replikat-Ansatz mit gleich großen
    Aufträgen ist einfach und robust, setzt aber homogene Hardware voraus und lastet
    unterschiedlich schnelle Knoten nicht unterschiedlich aus. Ein Modell mit
    dynamischer Arbeitsverteilung -- ein Manager vergibt Ray-Kontingente, sobald ein
    Knoten frei wird -- würde heterogene Hardware besser nutzen und die Knotenzahl
    flexibel machen.
\end{itemize}

Alle genannten Punkte betreffen Robustheit und Bedienbarkeit, nicht die
grundsätzliche Korrektheit des Verfahrens. Die verteilte Ausführung liefert eine
gültige, statistisch dichtere Impulsantwort, und das Manager-Worker-Muster über
SSH und \texttt{scp} erfüllt seinen Zweck im Laborbetrieb.

\subsection{Zusammenfassung}
\label{subsec:cluster-zusammenfassung}

Dieses Kapitel hat die verteilte Ausführung der Strahlverfolgung auf dem
Labor-Cluster beschrieben. Weil die Simulation embarrassingly parallel ist, genügt
ein Replikat-Ansatz: Alle acht Knoten rechnen denselben Auftrag mit unabhängigen
Zufallsströmen, was statistisch einer achtfach größeren Simulation entspricht und
den Schätzfehler der Impulsantwort um rund $\sqrt{8}$ senkt -- ein Verhalten, das
das numerische Experiment in Abbildung~\ref{fig:mc-convergence} bestätigt. Ein
Shell-Skript startet die Knoten per SSH parallel, synchronisiert über \texttt{wait}
und sammelt die Ergebnisse per \texttt{scp} ein; \texttt{merge\_results.py}
verkettet die Trefferlisten zu einer gemeinsamen Datei im gewohnten Schema. Der
Prototyp erfüllt seine Aufgabe, weist aber in Konfigurierbarkeit, Fehlererkennung
und der Anbindung an die DSP-Pipeline die diskutierten Schwächen auf, die den
Rahmen für die Weiterentwicklung abstecken.
 im Master.
% Benoetigte Pakete: amsmath, listings, graphicx, booktabs.
% Fuer saubere Optik im Master: fontenc, babel[ngerman], microtype.
% Abbildungen: latexDocs/figures/cluster_arch.png, mc_convergence.png
%
% STANDALONE-VORSCHAU:
%   \documentclass[12pt]{article}
%   \usepackage[utf8]{inputenc}\usepackage[T1]{fontenc}
%   \usepackage[ngerman]{babel}\usepackage{microtype}
%   \usepackage{amsmath}\usepackage{listings}
%   \usepackage{graphicx}\usepackage{booktabs}
%   \begin{document}
%   % =====================================================================
% Kapitel: Verteilte Ausfuehrung auf dem Rechencluster
% Autor: Mier
%
% EINBINDEN: \input{Kapitel_Cluster_Mier} im Master.
% Benoetigte Pakete: amsmath, listings, graphicx, booktabs.
% Fuer saubere Optik im Master: fontenc, babel[ngerman], microtype.
% Abbildungen: latexDocs/figures/cluster_arch.png, mc_convergence.png
%
% STANDALONE-VORSCHAU:
%   \documentclass[12pt]{article}
%   \usepackage[utf8]{inputenc}\usepackage[T1]{fontenc}
%   \usepackage[ngerman]{babel}\usepackage{microtype}
%   \usepackage{amsmath}\usepackage{listings}
%   \usepackage{graphicx}\usepackage{booktabs}
%   \begin{document}
%   \input{Kapitel_Cluster_Mier}
%   \end{document}
% =====================================================================

\lstdefinestyle{clusterlst}{
    basicstyle=\ttfamily\small,
    showstringspaces=false,
    breaklines=true,
    numbers=left,
    numberstyle=\tiny,
    frame=single,
    tabsize=2
}

%======================================================================
\section{Verteilte Ausführung auf dem Rechencluster}
\label{sec:cluster}

\subsection{Einordnung und Aufgabenstellung}
\label{subsec:cluster-einordnung}

Die Strahlverfolgung ist der teuerste Schritt der gesamten Pipeline. Für eine
brauchbare Impulsantwort müssen Hunderttausende bis Millionen einzelner
Schallstrahlen durch die Raumgeometrie verfolgt werden, und jeder dieser Strahlen
verlangt mehrere Schnitt- und Reflexionsberechnungen. Der entscheidende Vorteil
dabei: Kein Strahl hängt vom Ergebnis eines anderen ab. Das Problem ist damit
\emph{embarrassingly parallel} und lässt sich ohne Kommunikation zwischen den
Teilrechnungen auf beliebig viele Rechner verteilen.

Dieses Kapitel dokumentiert die Verteilungsschicht des Systems. Sie startet die
Rust-Simulation gleichzeitig auf allen acht Knoten des Labor-Clusters, wartet auf
deren Abschluss, sammelt die Teilergebnisse ein und fügt sie zu einer einzigen,
statistisch dichteren Impulsantwort zusammen. Aus Sicht der übrigen Pipeline
ändert sich dabei nichts: Am Ende steht wieder eine Datei im gewohnten Schema,
die die Signalverarbeitung unverändert weiterverarbeiten kann.

Zur Abgrenzung von den Nachbarkapiteln gehören zwei Punkte. Erstens ist die
\emph{prozessinterne} Nebenläufigkeit der Rust-Engine, also die Verteilung auf
die Kerne eines einzelnen Rechners mittels \texttt{rayon}, an anderer Stelle
beschrieben und bleibt hier außen vor. Sie parallelisiert innerhalb eines
Prozesses; dieses Kapitel parallelisiert über Prozess- und Rechnergrenzen hinweg.
Zweitens geht es nicht um den einzelnen Subprozess-Aufruf aus der Python-Pipeline,
sondern um die Ebene darüber: die gleichzeitige Ausführung auf mehreren physischen
Knoten und die anschließende Aggregation ihrer Ergebnisse.

\subsection{Theoretischer Hintergrund}
\label{subsec:cluster-theorie}

\subsubsection{Monte-Carlo-Strahlverfolgung}
\label{subsubsec:mc}

Die Engine schätzt die Impulsantwort eines Raumes stochastisch. Sie verschickt
zufällig gerichtete Strahlen von der Schallquelle, verfolgt ihre Reflexionen und
registriert, wann und mit welchem Restdruck sie den Empfänger erreichen. Jeder
einzelne Strahl ist eine Zufallsstichprobe; erst die Summe vieler Strahlen ergibt
ein aussagekräftiges Bild der Raumakustik. Solche Monte-Carlo-Verfahren haben eine
charakteristische Eigenschaft: Ihr Schätzfehler sinkt mit der Wurzel der
Stichprobengröße. Verdoppelt man die Zahl der Strahlen, halbiert sich der Fehler
also nicht, sondern sinkt nur um den Faktor $\sqrt{2}$. Genau diese langsame
Konvergenz macht große Strahlzahlen nötig und motiviert die Verteilung auf mehrere
Rechner.

\subsubsection{Zwei Wege der Verteilung}
\label{subsubsec:strategien}

Ein Monte-Carlo-Verfahren lässt sich auf zwei grundsätzlich verschiedene Arten
verteilen. Tabelle~\ref{tab:strategien} stellt sie gegenüber.

\begin{table}[htbp]
    \centering
    \caption{Zwei Strategien zur Verteilung der Strahlverfolgung.}
    \label{tab:strategien}
    \begin{tabular}{p{3.2cm} p{5.2cm} p{4.2cm}}
        \toprule
        & \textbf{Work-Partitioning} & \textbf{Replikat-Parallelismus} \\
        \midrule
        Prinzip & Gesamtzahl der Strahlen wird in disjunkte Anteile zerlegt & Jeder Knoten rechnet denselben Auftrag mit eigenem Zufallsstrom \\
        Koordination & Zuteilung der Anteile nötig & keine, jeder Knoten ist autark \\
        Zufallszahlen & ein globaler Strom, aufgeteilt & unabhängige Ströme pro Knoten \\
        Robustheit & Ausfall trifft einen festen Anteil & Ausfall kostet nur Stichprobengröße \\
        \bottomrule
    \end{tabular}
\end{table}

Das umgesetzte System verwendet den zweiten Weg. Alle acht Knoten starten die
Simulation mit derselben Konfigurationsdatei \texttt{input\_config.json} und
werfen dennoch jeweils andere Strahlen, weil jeder Prozess seinen eigenen
Zufallsstrom besitzt. Der Ansatz ist bewusst schlicht gewählt: Er braucht keine
zentrale Zuteilung und keinen Austausch zwischen den Knoten, sondern nur, dass am
Ende alle Ergebnisse eingesammelt werden.

\subsubsection{Unabhängigkeit der Zufallsströme}
\label{subsubsec:rng}

Die Korrektheit des Replikat-Ansatzes steht und fällt mit einer Bedingung: Die
Zufallszahlen auf den Knoten müssen unabhängig sein. Würde jeder Knoten dieselbe
Folge zufälliger Richtungen erzeugen, so lieferten alle acht exakt dasselbe
Ergebnis, und das Zusammenführen brächte keinen Erkenntnisgewinn, sondern nur
achtfach dieselben Strahlen. Die Rust-Engine schließt das aus, indem sie ihren
Zufallsstrom pro Thread bzw. Prozess frisch aus der Systementropie zieht:

\begin{lstlisting}[style=clusterlst, language=Rust, caption={Initialisierung des Zufallsstroms pro Prozess (\texttt{simulation.rs}).}]
let mut rng = rand::rng();                       // Omnidirectional, random rays
let random_angle: f32 = rng.random_range(0.0..(2.0 * PI));
let direction = Vec2::new(random_angle.cos(), random_angle.sin());
\end{lstlisting}

Weil jeder der acht Prozesse auf seinem Knoten einen anderen Startzustand erhält,
erzeugt jeder eine andere Stichprobe. Acht Knoten mit je einer Million Strahlen
liefern damit acht Millionen \emph{unabhängige} Stichproben derselben Raumakustik.

\subsubsection{Statistische Konvergenz}
\label{subsubsec:konvergenz}

Statistisch verhält sich der verteilte Lauf wie eine einzige, achtfach größere
Simulation. Für einen Monte-Carlo-Schätzer aus $N$ unabhängigen Teilstichproben
gleicher Größe sinkt die Standardabweichung des Mittelwerts mit
\[
    \sigma_N \;=\; \frac{\sigma_1}{\sqrt{N}} .
\]
Acht Knoten reduzieren die Streuung der geschätzten Impulsantwort demnach um den
Faktor $\sqrt{8} \approx 2{,}83$ gegenüber einem einzelnen Knoten.
Abbildung~\ref{fig:mc-convergence} bestätigt dieses Verhalten in einem numerischen
Experiment: Aufgetragen ist der mittlere quadratische Fehler der geschätzten
Impulsantwort über der Knotenzahl. Die Messpunkte folgen sauber der
$1/\sqrt{N}$-Referenzkurve; von einem auf acht Knoten sinkt der Fehler um den
Faktor $2{,}78$, praktisch exakt der theoretisch erwartete Wert.

\begin{figure}[htbp]
    \centering
    \includegraphics[width=0.8\textwidth]{figures/mc_convergence}
    \caption{Mittlerer quadratischer Fehler der geschätzten Impulsantwort über der
    Knotenzahl $N$ (je $5\cdot10^4$ Strahlen, doppelt logarithmisch). Die
    Messpunkte folgen der $1/\sqrt{N}$-Referenz; acht Knoten senken den Fehler um
    den Faktor $2{,}78 \approx \sqrt{8}$.}
    \label{fig:mc-convergence}
\end{figure}

Bemerkenswert ist die Kehrseite dieser Kurve: Um den Fehler zu halbieren, braucht
es die vierfache Knotenzahl. Der Nutzen zusätzlicher Knoten nimmt also ab, was bei
der Wahl der Clustergröße zu berücksichtigen ist. Für die hier angestrebte
Genauigkeit sind acht Knoten ein sinnvoller Kompromiss zwischen Aufwand und
Qualität.

\subsection{Architektur des Clusters}
\label{subsec:cluster-arch}

Das System folgt einem \emph{Manager-Worker}-Muster
(Abbildung~\ref{fig:cluster-arch}). Ein Manager-Knoten koordiniert den Ablauf,
acht Worker-Knoten führen die eigentliche Simulation aus. Der Manager selbst
rechnet nicht mit, sondern startet die Worker, wartet auf sie und führt am Ende
die Ergebnisse zusammen.

\begin{figure}[htbp]
    \centering
    \includegraphics[width=0.82\textwidth]{figures/cluster_arch}
    \caption{Manager-Worker-Architektur. Der Manager startet die acht Knoten
    parallel per \texttt{ssh}, holt nach Abschluss jede Ergebnisdatei per
    \texttt{scp} zurück und führt sie mit \texttt{merge\_results.py} zusammen.}
    \label{fig:cluster-arch}
\end{figure}

Die Kommunikation läuft vollständig über Standard-Werkzeuge des Betriebssystems.
\texttt{ssh} startet die Prozesse aus der Ferne, \texttt{scp} überträgt die
Ergebnisdateien zurück. Ein eigenes Netzwerkprotokoll oder ein
Message-Passing-Framework kommt nicht zum Einsatz. Für einen Stapelbetrieb, bei
dem jeder Knoten genau einmal startet, eine Weile rechnet und ein einziges
Ergebnis zurückliefert, wäre solcher Aufwand nicht gerechtfertigt: Es gibt keine
laufende Zustandssynchronisation und keinen Nachrichtenaustausch während der
Rechnung, den ein Framework verwalten müsste.

Damit das funktioniert, müssen drei Voraussetzungen erfüllt sein. Auf jedem Knoten
liegt dieselbe vorkompilierte Binärdatei \texttt{rust\_service} samt der
Konfiguration \texttt{input\_config.json}. Der Manager kann sich per SSH ohne
interaktive Passworteingabe anmelden, was eine Schlüssel-basierte
Authentifizierung voraussetzt. Und die Knoten sind über feste Adressen im
Labornetz erreichbar.

\subsection{Ablaufsteuerung: \texttt{run\_cluster.sh}}
\label{subsec:cluster-run}

Das Skript \texttt{run\_cluster.sh} steuert den verteilten Lauf in drei Phasen:
paralleles Starten, Warten und Einsammeln.

\begin{lstlisting}[style=clusterlst, language=bash, caption={Verteilter Start und Einsammeln der Ergebnisse (gekürzt).}]
# Phase 1: alle Knoten gleichzeitig starten (Hintergrundprozesse)
ssh node1@139.6.65.29 "./rust_service --config-file input_config.json" &
ssh node2@139.6.65.18 "./rust_service --config-file input_config.json" &
# ... node3 bis node8 analog ...

wait                       # Phase 2: auf Abschluss ALLER Knoten warten

# Phase 3: Ergebnis jedes Knotens zum Manager kopieren
scp node1@139.6.65.29:~/ir_output.json ~/ir_node1.json
scp node2@139.6.65.18:~/ir_output.json ~/ir_node2.json
# ... node3 bis node8 analog ...
\end{lstlisting}

Zwei Eigenschaften der Shell nutzt der Ablauf gezielt. Das kaufmännische Und
(\texttt{\&}) am Zeilenende schickt jeden \texttt{ssh}-Aufruf in den Hintergrund.
Die acht Simulationen laufen dadurch nahezu gleichzeitig an, statt nacheinander
abgearbeitet zu werden. Das darauffolgende \texttt{wait} blockiert, bis
\emph{alle} Hintergrundprozesse beendet sind, und bildet so eine
Synchronisationsbarriere zwischen der Rechen- und der Einsammel-Phase. Erst wenn
der letzte Knoten fertig ist, beginnt \texttt{scp}, von jedem Knoten dessen
\texttt{ir\_output.json} zu holen und unter einem knotenspezifischen Namen auf dem
Manager abzulegen.

Diese Barriere hat eine wichtige Konsequenz für die Laufzeit: Der Gesamtlauf ist
so langsam wie der \emph{langsamste} Knoten. Solange die Hardware homogen ist,
fällt das kaum ins Gewicht; auf ungleichen Knoten würde ein einzelner langsamer
Rechner den ganzen Lauf ausbremsen, während die übrigen sieben bereits fertig
warten. Dieses Straggler-Problem ist der zentrale Grund, warum ein statischer
Replikat-Ansatz mit gleich großen Aufträgen homogene Hardware voraussetzt.

\subsection{Aggregation der Teilergebnisse: \texttt{merge\_results.py}}
\label{subsec:cluster-merge}

Sobald alle Teilergebnisse auf dem Manager liegen, führt \texttt{merge\_results.py}
sie zu einer einzigen Impulsantwort zusammen. Jeder Knoten liefert dasselbe
Ausgabeschema, nämlich Metadaten sowie zwei positionsgleiche Listen aus
Verzögerungen und Drücken. Die Aggregation besteht deshalb im Kern aus einer
schlichten Verkettung dieser Listen:

\begin{lstlisting}[style=clusterlst, language=Python, caption={Zusammenführen der Knoten-Ergebnisse, gekürzt (\texttt{merge\_results.py}).}]
all_delays, all_pressures = [], []
for f in sorted(glob.glob("/home/manager/ir_node*.json")):
    with open(f) as file:
        data = json.load(file)
        all_delays    += data["hits"]["delays_seconds"]
        all_pressures += data["hits"]["pressures"]

merged = {
    "metadata": {"sample_rate": 44100, "rays_cast": 1000000 * 8,
                 "rays_received": len(all_delays), "room_name": "MVP_Test_Room_1"},
    "hits": {"delays_seconds": all_delays, "pressures": all_pressures},
}
json.dump(merged, open("/home/manager/ir_merged.json", "w"), indent=2)
\end{lstlisting}

Dass die Verkettung hier die richtige Operation ist, folgt direkt aus dem
Replikat-Modell. Weil die einzelnen Läufe unabhängige Stichproben desselben Raums
sind, ist die Vereinigung aller Treffer statistisch identisch mit dem Ergebnis
eines einzigen Laufs mit acht Millionen Strahlen. Es findet keine Mittelung und
keine Gewichtung statt, sondern nur ein Aneinanderhängen. Das Feld
\texttt{rays\_cast} ist im Skript fest auf $8 \cdot 10^{6}$ gesetzt; solange der in
Abschnitt~\ref{subsec:cluster-grenzen} beschriebene Kopierfehler bei Knoten~6
besteht, fließen tatsächlich nur sieben Knoten in die Aggregation ein, und dieser
Metadatenwert überschätzt die reale Strahlzahl. \texttt{rays\_received} zählt
dagegen die tatsächlich empfangenen Treffer. Die
zusammengeführte Datei trägt exakt dasselbe Schema wie die Ausgabe eines einzelnen
Knotens und ist damit ohne weitere Anpassung für die DSP-Pipeline verwendbar.

\subsection{Datenfluss im Überblick}
\label{subsec:cluster-datenfluss}

Der gesamte verteilte Ablauf lässt sich in fünf Stufen zusammenfassen:

\begin{enumerate}
    \item Der Manager startet per SSH auf jedem der acht Knoten
    \texttt{rust\_service} mit identischer Konfiguration, alle parallel.
    \item Jeder Knoten simuliert unabhängig eine Million Strahlen und schreibt
    lokal seine \texttt{ir\_output.json}.
    \item \texttt{wait} synchronisiert: Die Einsammel-Phase beginnt erst, wenn der
    letzte Knoten fertig ist.
    \item \texttt{scp} überträgt jede \texttt{ir\_output.json} als
    \texttt{ir\_node$k$.json} zum Manager.
    \item \texttt{merge\_results.py} verkettet alle Trefferlisten zu einer
    \texttt{ir\_merged.json}, die anschließend die DSP-Pipeline verarbeitet.
\end{enumerate}

\subsection{Leistungsbetrachtung}
\label{subsec:cluster-leistung}

\subsubsection{Erwarteter Speedup}
\label{subsubsec:speedup}

Bei der Leistung sind zwei Sichtweisen zu trennen. In der umgesetzten
Replikat-Betriebsart rechnet jeder Knoten dieselbe Strahlzahl, sodass acht Knoten
in nahezu derselben Wanduhrzeit die achtfache Stichprobe liefern. Der Gewinn ist
hier keine kürzere Laufzeit, sondern die höhere Genauigkeit aus
Abbildung~\ref{fig:mc-convergence}; man spricht vom schwachen Skalieren im Sinne
von Gustafson, bei dem mit der Rechnerzahl die bearbeitete Problemgröße wächst,
nicht die Geschwindigkeit einer festen Aufgabe. Will man dagegen eine feste
Gesamtstrahlzahl in kürzerer Zeit erreichen, so verteilt man diese auf die Knoten,
sodass jeder nur $8\cdot10^{6}/k$ Strahlen rechnet. Erst in dieser
Partitionierungs-Betriebsart greift das Gesetz von Amdahl. Bezeichnet $p$ den parallelen
Anteil der Arbeit, hier die reine Strahlverfolgung, und $1-p$ den seriellen Rest
aus Prozessstart, Ergebnistransfer und Aggregation, so beträgt der theoretische
Speedup
\[
    S(k) \;=\; \frac{1}{(1-p) + \dfrac{p}{k}} .
\]
Da die Strahlverfolgung die Laufzeit klar dominiert, liegt $p$ nahe bei eins, und
$S(k)$ steigt bis $k=8$ nahezu linear. In der Praxis bleibt der reale Speedup
darunter, und zwar aus drei Gründen: dem bereits erwähnten Straggler-Effekt (der
Lauf ist so schnell wie der langsamste Knoten), dem Overhead für den
SSH-Verbindungsaufbau sowie der Übertragung und Zusammenführung der
Ergebnisdateien.

\subsubsection{Kommunikations- und I/O-Overhead}
\label{subsubsec:overhead}

Der Overhead dieses Ansatzes ist klein, aber nicht null. Pro Knoten fallen ein
SSH-Verbindungsaufbau und eine \texttt{scp}-Übertragung an. Übertragen wird jeweils
eine JSON-Datei, deren Größe mit der Zahl der empfangenen Treffer wächst; bei
einer Million geworfener Strahlen liegt sie im Bereich einiger Megabyte. Da diese
acht Übertragungen erst \emph{nach} der Rechenphase stattfinden und nicht mit ihr
überlappen, addieren sie sich voll zur Gesamtlaufzeit. Bei deutlich mehr Knoten
oder größeren Ergebnisdateien würde dieser sequenzielle Einsammel-Schritt zum
Flaschenhals; eine Überlappung von Rechnung und Transfer oder ein kompakteres
Ausgabeformat als das gut lesbare, aber platzintensive JSON wären dann die
naheliegenden Gegenmaßnahmen.

\subsubsection{Messung am Cluster}
\label{subsubsec:messung}

Die konkreten Laufzeiten hängen von der Laborhardware ab und sind am Cluster zu
messen. Um einen aussagekräftigen Speedup zu erhalten, wird in der
Partitionierungs-Betriebsart gemessen: In \texttt{input\_config.json} wird
\texttt{rays\_to\_cast} auf $8\cdot10^{6}/k$ gesetzt, sodass die Gesamtarbeit über
alle $k$ Knoten konstant bleibt. Dann wird \texttt{run\_cluster.sh} für
verschiedene Knotenzahlen mit der Shell-eigenen Zeitmessung ausgeführt:

\begin{lstlisting}[style=clusterlst, language=bash, caption={Messvorschrift für die Laufzeit am Cluster.}]
time ./run_cluster.sh        # Wanduhrzeit fuer den kompletten verteilten Lauf
\end{lstlisting}

Tabelle~\ref{tab:speedup} nimmt die Messwerte auf. Als Referenz dient der Lauf mit
einem einzelnen Knoten; der Speedup ergibt sich als Verhältnis der Wanduhrzeiten.

\begin{table}[htbp]
    \centering
    \caption{Wanduhrzeit und Speedup des verteilten Laufs (am Cluster zu messen).}
    \label{tab:speedup}
    \begin{tabular}{r r r}
        \toprule
        Knoten $k$ & Wanduhrzeit [s] & Speedup $S$ \\
        \midrule
        1 & \dots & $1{,}00$ \\
        2 & \dots & \dots \\
        4 & \dots & \dots \\
        8 & \dots & \dots \\
        \bottomrule
    \end{tabular}
\end{table}

\subsection{Fehlertoleranz und Sicherheit}
\label{subsec:cluster-fehler}

Zwei Querschnittsthemen verdienen eigene Betrachtung. Beim \textbf{Sicherheits}aspekt
ist die Schlüssel-basierte SSH-Authentifizierung zentral: Der Manager benötigt
passwortlosen Zugriff auf jeden Knoten, was über hinterlegte öffentliche Schlüssel
gelöst ist. Der private Schlüssel des Managers ist damit das entscheidende
Geheimnis; wer ihn besitzt, kann auf allen Knoten Code ausführen. In einem
abgeschlossenen Labornetz ist das vertretbar, für einen offeneren Betrieb wäre eine
engere Rechtevergabe angebracht.

Beim Thema \textbf{Fehlertoleranz} zeigt der aktuelle Stand seine deutlichste
Schwäche, die im folgenden Abschnitt zusammen mit den übrigen Grenzen diskutiert
wird. Der Kern des Problems: Ein ausgefallener Knoten wird nicht erkannt, sondern
schlicht übergangen, sodass die Stichprobengröße unbemerkt sinkt.

\subsection{Grenzen des aktuellen Standes und Ausblick}
\label{subsec:cluster-grenzen}

Die Verteilungsschicht ist als funktionsfähiger Prototyp umgesetzt. Für einen
robusten Betrieb bestehen mehrere klar benannte Schwachstellen, die zugleich die
nächsten Entwicklungsschritte vorgeben.

\begin{itemize}
    \item \textbf{Fest verdrahtete Konfiguration.} IP-Adressen, Knotennamen, die
    Knotenzahl und die Ablagepfade stehen direkt im Skript. Ein Umzug auf andere
    Hardware oder eine andere Knotenzahl erfordert Änderungen am Quelltext. Eine
    Auslagerung dieser Werte in eine Konfigurationsdatei wäre der naheliegende
    erste Schritt.

    \item \textbf{Kopierfehler bei Knoten~6.} Im Einsammel-Schritt wird das
    Ergebnis von Knoten~6 versehentlich nach \texttt{ir\_node5.json} kopiert und
    überschreibt damit das Ergebnis von Knoten~5. Effektiv fließen so nur sieben
    statt acht unabhängige Stichproben in die Aggregation ein. Der Fehler ist ein
    Tippfehler im Zielnamen und leicht zu beheben, führt aber vor Augen, wie
    fehleranfällig die von Hand ausgeschriebene Knotenliste ist -- ein weiteres
    Argument für die genannte Auslagerung in eine Schleife über eine
    Konfigurationsliste.

    \item \textbf{Keine Fehlererkennung.} Fällt ein Knoten aus oder liefert er kein
    Ergebnis, schlägt der zugehörige \texttt{scp}-Aufruf fehl, der Gesamtlauf setzt
    sich aber fort und aggregiert stillschweigend über die vorhandenen Dateien. Ein
    Knotenausfall bleibt damit unbemerkt und verfälscht die Stichprobengröße. Eine
    Prüfung der Rückgabecodes und eine Warnung bei fehlenden Ergebnissen wären hier
    angebracht.

    \item \textbf{Manuelle Kopplung an die DSP-Pipeline.} Die Aggregation erzeugt
    \texttt{ir\_merged.json}, die Python-DSP-Pipeline liest jedoch
    \texttt{ir\_output.json}. Der Übergang zwischen verteiltem Lauf und Faltung ist
    derzeit ein manueller Schritt. Eine saubere Integration würde den Cluster-Lauf
    direkt aus der Orchestrierung heraus anstoßen und die zusammengeführte Datei
    automatisch weiterreichen.

    \item \textbf{Statische Verteilung.} Der Replikat-Ansatz mit gleich großen
    Aufträgen ist einfach und robust, setzt aber homogene Hardware voraus und lastet
    unterschiedlich schnelle Knoten nicht unterschiedlich aus. Ein Modell mit
    dynamischer Arbeitsverteilung -- ein Manager vergibt Ray-Kontingente, sobald ein
    Knoten frei wird -- würde heterogene Hardware besser nutzen und die Knotenzahl
    flexibel machen.
\end{itemize}

Alle genannten Punkte betreffen Robustheit und Bedienbarkeit, nicht die
grundsätzliche Korrektheit des Verfahrens. Die verteilte Ausführung liefert eine
gültige, statistisch dichtere Impulsantwort, und das Manager-Worker-Muster über
SSH und \texttt{scp} erfüllt seinen Zweck im Laborbetrieb.

\subsection{Zusammenfassung}
\label{subsec:cluster-zusammenfassung}

Dieses Kapitel hat die verteilte Ausführung der Strahlverfolgung auf dem
Labor-Cluster beschrieben. Weil die Simulation embarrassingly parallel ist, genügt
ein Replikat-Ansatz: Alle acht Knoten rechnen denselben Auftrag mit unabhängigen
Zufallsströmen, was statistisch einer achtfach größeren Simulation entspricht und
den Schätzfehler der Impulsantwort um rund $\sqrt{8}$ senkt -- ein Verhalten, das
das numerische Experiment in Abbildung~\ref{fig:mc-convergence} bestätigt. Ein
Shell-Skript startet die Knoten per SSH parallel, synchronisiert über \texttt{wait}
und sammelt die Ergebnisse per \texttt{scp} ein; \texttt{merge\_results.py}
verkettet die Trefferlisten zu einer gemeinsamen Datei im gewohnten Schema. Der
Prototyp erfüllt seine Aufgabe, weist aber in Konfigurierbarkeit, Fehlererkennung
und der Anbindung an die DSP-Pipeline die diskutierten Schwächen auf, die den
Rahmen für die Weiterentwicklung abstecken.

%   \end{document}
% =====================================================================

\lstdefinestyle{clusterlst}{
    basicstyle=\ttfamily\small,
    showstringspaces=false,
    breaklines=true,
    numbers=left,
    numberstyle=\tiny,
    frame=single,
    tabsize=2
}

%======================================================================
\section{Verteilte Ausführung auf dem Rechencluster}
\label{sec:cluster}

\subsection{Einordnung und Aufgabenstellung}
\label{subsec:cluster-einordnung}

Die Strahlverfolgung ist der teuerste Schritt der gesamten Pipeline. Für eine
brauchbare Impulsantwort müssen Hunderttausende bis Millionen einzelner
Schallstrahlen durch die Raumgeometrie verfolgt werden, und jeder dieser Strahlen
verlangt mehrere Schnitt- und Reflexionsberechnungen. Der entscheidende Vorteil
dabei: Kein Strahl hängt vom Ergebnis eines anderen ab. Das Problem ist damit
\emph{embarrassingly parallel} und lässt sich ohne Kommunikation zwischen den
Teilrechnungen auf beliebig viele Rechner verteilen.

Dieses Kapitel dokumentiert die Verteilungsschicht des Systems. Sie startet die
Rust-Simulation gleichzeitig auf allen acht Knoten des Labor-Clusters, wartet auf
deren Abschluss, sammelt die Teilergebnisse ein und fügt sie zu einer einzigen,
statistisch dichteren Impulsantwort zusammen. Aus Sicht der übrigen Pipeline
ändert sich dabei nichts: Am Ende steht wieder eine Datei im gewohnten Schema,
die die Signalverarbeitung unverändert weiterverarbeiten kann.

Zur Abgrenzung von den Nachbarkapiteln gehören zwei Punkte. Erstens ist die
\emph{prozessinterne} Nebenläufigkeit der Rust-Engine, also die Verteilung auf
die Kerne eines einzelnen Rechners mittels \texttt{rayon}, an anderer Stelle
beschrieben und bleibt hier außen vor. Sie parallelisiert innerhalb eines
Prozesses; dieses Kapitel parallelisiert über Prozess- und Rechnergrenzen hinweg.
Zweitens geht es nicht um den einzelnen Subprozess-Aufruf aus der Python-Pipeline,
sondern um die Ebene darüber: die gleichzeitige Ausführung auf mehreren physischen
Knoten und die anschließende Aggregation ihrer Ergebnisse.

\subsection{Theoretischer Hintergrund}
\label{subsec:cluster-theorie}

\subsubsection{Monte-Carlo-Strahlverfolgung}
\label{subsubsec:mc}

Die Engine schätzt die Impulsantwort eines Raumes stochastisch. Sie verschickt
zufällig gerichtete Strahlen von der Schallquelle, verfolgt ihre Reflexionen und
registriert, wann und mit welchem Restdruck sie den Empfänger erreichen. Jeder
einzelne Strahl ist eine Zufallsstichprobe; erst die Summe vieler Strahlen ergibt
ein aussagekräftiges Bild der Raumakustik. Solche Monte-Carlo-Verfahren haben eine
charakteristische Eigenschaft: Ihr Schätzfehler sinkt mit der Wurzel der
Stichprobengröße. Verdoppelt man die Zahl der Strahlen, halbiert sich der Fehler
also nicht, sondern sinkt nur um den Faktor $\sqrt{2}$. Genau diese langsame
Konvergenz macht große Strahlzahlen nötig und motiviert die Verteilung auf mehrere
Rechner.

\subsubsection{Zwei Wege der Verteilung}
\label{subsubsec:strategien}

Ein Monte-Carlo-Verfahren lässt sich auf zwei grundsätzlich verschiedene Arten
verteilen. Tabelle~\ref{tab:strategien} stellt sie gegenüber.

\begin{table}[htbp]
    \centering
    \caption{Zwei Strategien zur Verteilung der Strahlverfolgung.}
    \label{tab:strategien}
    \begin{tabular}{p{3.2cm} p{5.2cm} p{4.2cm}}
        \toprule
        & \textbf{Work-Partitioning} & \textbf{Replikat-Parallelismus} \\
        \midrule
        Prinzip & Gesamtzahl der Strahlen wird in disjunkte Anteile zerlegt & Jeder Knoten rechnet denselben Auftrag mit eigenem Zufallsstrom \\
        Koordination & Zuteilung der Anteile nötig & keine, jeder Knoten ist autark \\
        Zufallszahlen & ein globaler Strom, aufgeteilt & unabhängige Ströme pro Knoten \\
        Robustheit & Ausfall trifft einen festen Anteil & Ausfall kostet nur Stichprobengröße \\
        \bottomrule
    \end{tabular}
\end{table}

Das umgesetzte System verwendet den zweiten Weg. Alle acht Knoten starten die
Simulation mit derselben Konfigurationsdatei \texttt{input\_config.json} und
werfen dennoch jeweils andere Strahlen, weil jeder Prozess seinen eigenen
Zufallsstrom besitzt. Der Ansatz ist bewusst schlicht gewählt: Er braucht keine
zentrale Zuteilung und keinen Austausch zwischen den Knoten, sondern nur, dass am
Ende alle Ergebnisse eingesammelt werden.

\subsubsection{Unabhängigkeit der Zufallsströme}
\label{subsubsec:rng}

Die Korrektheit des Replikat-Ansatzes steht und fällt mit einer Bedingung: Die
Zufallszahlen auf den Knoten müssen unabhängig sein. Würde jeder Knoten dieselbe
Folge zufälliger Richtungen erzeugen, so lieferten alle acht exakt dasselbe
Ergebnis, und das Zusammenführen brächte keinen Erkenntnisgewinn, sondern nur
achtfach dieselben Strahlen. Die Rust-Engine schließt das aus, indem sie ihren
Zufallsstrom pro Thread bzw. Prozess frisch aus der Systementropie zieht:

\begin{lstlisting}[style=clusterlst, language=Rust, caption={Initialisierung des Zufallsstroms pro Prozess (\texttt{simulation.rs}).}]
let mut rng = rand::rng();                       // Omnidirectional, random rays
let random_angle: f32 = rng.random_range(0.0..(2.0 * PI));
let direction = Vec2::new(random_angle.cos(), random_angle.sin());
\end{lstlisting}

Weil jeder der acht Prozesse auf seinem Knoten einen anderen Startzustand erhält,
erzeugt jeder eine andere Stichprobe. Acht Knoten mit je einer Million Strahlen
liefern damit acht Millionen \emph{unabhängige} Stichproben derselben Raumakustik.

\subsubsection{Statistische Konvergenz}
\label{subsubsec:konvergenz}

Statistisch verhält sich der verteilte Lauf wie eine einzige, achtfach größere
Simulation. Für einen Monte-Carlo-Schätzer aus $N$ unabhängigen Teilstichproben
gleicher Größe sinkt die Standardabweichung des Mittelwerts mit
\[
    \sigma_N \;=\; \frac{\sigma_1}{\sqrt{N}} .
\]
Acht Knoten reduzieren die Streuung der geschätzten Impulsantwort demnach um den
Faktor $\sqrt{8} \approx 2{,}83$ gegenüber einem einzelnen Knoten.
Abbildung~\ref{fig:mc-convergence} bestätigt dieses Verhalten in einem numerischen
Experiment: Aufgetragen ist der mittlere quadratische Fehler der geschätzten
Impulsantwort über der Knotenzahl. Die Messpunkte folgen sauber der
$1/\sqrt{N}$-Referenzkurve; von einem auf acht Knoten sinkt der Fehler um den
Faktor $2{,}78$, praktisch exakt der theoretisch erwartete Wert.

\begin{figure}[htbp]
    \centering
    \includegraphics[width=0.8\textwidth]{figures/mc_convergence}
    \caption{Mittlerer quadratischer Fehler der geschätzten Impulsantwort über der
    Knotenzahl $N$ (je $5\cdot10^4$ Strahlen, doppelt logarithmisch). Die
    Messpunkte folgen der $1/\sqrt{N}$-Referenz; acht Knoten senken den Fehler um
    den Faktor $2{,}78 \approx \sqrt{8}$.}
    \label{fig:mc-convergence}
\end{figure}

Bemerkenswert ist die Kehrseite dieser Kurve: Um den Fehler zu halbieren, braucht
es die vierfache Knotenzahl. Der Nutzen zusätzlicher Knoten nimmt also ab, was bei
der Wahl der Clustergröße zu berücksichtigen ist. Für die hier angestrebte
Genauigkeit sind acht Knoten ein sinnvoller Kompromiss zwischen Aufwand und
Qualität.

\subsection{Architektur des Clusters}
\label{subsec:cluster-arch}

Das System folgt einem \emph{Manager-Worker}-Muster
(Abbildung~\ref{fig:cluster-arch}). Ein Manager-Knoten koordiniert den Ablauf,
acht Worker-Knoten führen die eigentliche Simulation aus. Der Manager selbst
rechnet nicht mit, sondern startet die Worker, wartet auf sie und führt am Ende
die Ergebnisse zusammen.

\begin{figure}[htbp]
    \centering
    \includegraphics[width=0.82\textwidth]{figures/cluster_arch}
    \caption{Manager-Worker-Architektur. Der Manager startet die acht Knoten
    parallel per \texttt{ssh}, holt nach Abschluss jede Ergebnisdatei per
    \texttt{scp} zurück und führt sie mit \texttt{merge\_results.py} zusammen.}
    \label{fig:cluster-arch}
\end{figure}

Die Kommunikation läuft vollständig über Standard-Werkzeuge des Betriebssystems.
\texttt{ssh} startet die Prozesse aus der Ferne, \texttt{scp} überträgt die
Ergebnisdateien zurück. Ein eigenes Netzwerkprotokoll oder ein
Message-Passing-Framework kommt nicht zum Einsatz. Für einen Stapelbetrieb, bei
dem jeder Knoten genau einmal startet, eine Weile rechnet und ein einziges
Ergebnis zurückliefert, wäre solcher Aufwand nicht gerechtfertigt: Es gibt keine
laufende Zustandssynchronisation und keinen Nachrichtenaustausch während der
Rechnung, den ein Framework verwalten müsste.

Damit das funktioniert, müssen drei Voraussetzungen erfüllt sein. Auf jedem Knoten
liegt dieselbe vorkompilierte Binärdatei \texttt{rust\_service} samt der
Konfiguration \texttt{input\_config.json}. Der Manager kann sich per SSH ohne
interaktive Passworteingabe anmelden, was eine Schlüssel-basierte
Authentifizierung voraussetzt. Und die Knoten sind über feste Adressen im
Labornetz erreichbar.

\subsection{Ablaufsteuerung: \texttt{run\_cluster.sh}}
\label{subsec:cluster-run}

Das Skript \texttt{run\_cluster.sh} steuert den verteilten Lauf in drei Phasen:
paralleles Starten, Warten und Einsammeln.

\begin{lstlisting}[style=clusterlst, language=bash, caption={Verteilter Start und Einsammeln der Ergebnisse (gekürzt).}]
# Phase 1: alle Knoten gleichzeitig starten (Hintergrundprozesse)
ssh node1@139.6.65.29 "./rust_service --config-file input_config.json" &
ssh node2@139.6.65.18 "./rust_service --config-file input_config.json" &
# ... node3 bis node8 analog ...

wait                       # Phase 2: auf Abschluss ALLER Knoten warten

# Phase 3: Ergebnis jedes Knotens zum Manager kopieren
scp node1@139.6.65.29:~/ir_output.json ~/ir_node1.json
scp node2@139.6.65.18:~/ir_output.json ~/ir_node2.json
# ... node3 bis node8 analog ...
\end{lstlisting}

Zwei Eigenschaften der Shell nutzt der Ablauf gezielt. Das kaufmännische Und
(\texttt{\&}) am Zeilenende schickt jeden \texttt{ssh}-Aufruf in den Hintergrund.
Die acht Simulationen laufen dadurch nahezu gleichzeitig an, statt nacheinander
abgearbeitet zu werden. Das darauffolgende \texttt{wait} blockiert, bis
\emph{alle} Hintergrundprozesse beendet sind, und bildet so eine
Synchronisationsbarriere zwischen der Rechen- und der Einsammel-Phase. Erst wenn
der letzte Knoten fertig ist, beginnt \texttt{scp}, von jedem Knoten dessen
\texttt{ir\_output.json} zu holen und unter einem knotenspezifischen Namen auf dem
Manager abzulegen.

Diese Barriere hat eine wichtige Konsequenz für die Laufzeit: Der Gesamtlauf ist
so langsam wie der \emph{langsamste} Knoten. Solange die Hardware homogen ist,
fällt das kaum ins Gewicht; auf ungleichen Knoten würde ein einzelner langsamer
Rechner den ganzen Lauf ausbremsen, während die übrigen sieben bereits fertig
warten. Dieses Straggler-Problem ist der zentrale Grund, warum ein statischer
Replikat-Ansatz mit gleich großen Aufträgen homogene Hardware voraussetzt.

\subsection{Aggregation der Teilergebnisse: \texttt{merge\_results.py}}
\label{subsec:cluster-merge}

Sobald alle Teilergebnisse auf dem Manager liegen, führt \texttt{merge\_results.py}
sie zu einer einzigen Impulsantwort zusammen. Jeder Knoten liefert dasselbe
Ausgabeschema, nämlich Metadaten sowie zwei positionsgleiche Listen aus
Verzögerungen und Drücken. Die Aggregation besteht deshalb im Kern aus einer
schlichten Verkettung dieser Listen:

\begin{lstlisting}[style=clusterlst, language=Python, caption={Zusammenführen der Knoten-Ergebnisse, gekürzt (\texttt{merge\_results.py}).}]
all_delays, all_pressures = [], []
for f in sorted(glob.glob("/home/manager/ir_node*.json")):
    with open(f) as file:
        data = json.load(file)
        all_delays    += data["hits"]["delays_seconds"]
        all_pressures += data["hits"]["pressures"]

merged = {
    "metadata": {"sample_rate": 44100, "rays_cast": 1000000 * 8,
                 "rays_received": len(all_delays), "room_name": "MVP_Test_Room_1"},
    "hits": {"delays_seconds": all_delays, "pressures": all_pressures},
}
json.dump(merged, open("/home/manager/ir_merged.json", "w"), indent=2)
\end{lstlisting}

Dass die Verkettung hier die richtige Operation ist, folgt direkt aus dem
Replikat-Modell. Weil die einzelnen Läufe unabhängige Stichproben desselben Raums
sind, ist die Vereinigung aller Treffer statistisch identisch mit dem Ergebnis
eines einzigen Laufs mit acht Millionen Strahlen. Es findet keine Mittelung und
keine Gewichtung statt, sondern nur ein Aneinanderhängen. Das Feld
\texttt{rays\_cast} ist im Skript fest auf $8 \cdot 10^{6}$ gesetzt; solange der in
Abschnitt~\ref{subsec:cluster-grenzen} beschriebene Kopierfehler bei Knoten~6
besteht, fließen tatsächlich nur sieben Knoten in die Aggregation ein, und dieser
Metadatenwert überschätzt die reale Strahlzahl. \texttt{rays\_received} zählt
dagegen die tatsächlich empfangenen Treffer. Die
zusammengeführte Datei trägt exakt dasselbe Schema wie die Ausgabe eines einzelnen
Knotens und ist damit ohne weitere Anpassung für die DSP-Pipeline verwendbar.

\subsection{Datenfluss im Überblick}
\label{subsec:cluster-datenfluss}

Der gesamte verteilte Ablauf lässt sich in fünf Stufen zusammenfassen:

\begin{enumerate}
    \item Der Manager startet per SSH auf jedem der acht Knoten
    \texttt{rust\_service} mit identischer Konfiguration, alle parallel.
    \item Jeder Knoten simuliert unabhängig eine Million Strahlen und schreibt
    lokal seine \texttt{ir\_output.json}.
    \item \texttt{wait} synchronisiert: Die Einsammel-Phase beginnt erst, wenn der
    letzte Knoten fertig ist.
    \item \texttt{scp} überträgt jede \texttt{ir\_output.json} als
    \texttt{ir\_node$k$.json} zum Manager.
    \item \texttt{merge\_results.py} verkettet alle Trefferlisten zu einer
    \texttt{ir\_merged.json}, die anschließend die DSP-Pipeline verarbeitet.
\end{enumerate}

\subsection{Leistungsbetrachtung}
\label{subsec:cluster-leistung}

\subsubsection{Erwarteter Speedup}
\label{subsubsec:speedup}

Bei der Leistung sind zwei Sichtweisen zu trennen. In der umgesetzten
Replikat-Betriebsart rechnet jeder Knoten dieselbe Strahlzahl, sodass acht Knoten
in nahezu derselben Wanduhrzeit die achtfache Stichprobe liefern. Der Gewinn ist
hier keine kürzere Laufzeit, sondern die höhere Genauigkeit aus
Abbildung~\ref{fig:mc-convergence}; man spricht vom schwachen Skalieren im Sinne
von Gustafson, bei dem mit der Rechnerzahl die bearbeitete Problemgröße wächst,
nicht die Geschwindigkeit einer festen Aufgabe. Will man dagegen eine feste
Gesamtstrahlzahl in kürzerer Zeit erreichen, so verteilt man diese auf die Knoten,
sodass jeder nur $8\cdot10^{6}/k$ Strahlen rechnet. Erst in dieser
Partitionierungs-Betriebsart greift das Gesetz von Amdahl. Bezeichnet $p$ den parallelen
Anteil der Arbeit, hier die reine Strahlverfolgung, und $1-p$ den seriellen Rest
aus Prozessstart, Ergebnistransfer und Aggregation, so beträgt der theoretische
Speedup
\[
    S(k) \;=\; \frac{1}{(1-p) + \dfrac{p}{k}} .
\]
Da die Strahlverfolgung die Laufzeit klar dominiert, liegt $p$ nahe bei eins, und
$S(k)$ steigt bis $k=8$ nahezu linear. In der Praxis bleibt der reale Speedup
darunter, und zwar aus drei Gründen: dem bereits erwähnten Straggler-Effekt (der
Lauf ist so schnell wie der langsamste Knoten), dem Overhead für den
SSH-Verbindungsaufbau sowie der Übertragung und Zusammenführung der
Ergebnisdateien.

\subsubsection{Kommunikations- und I/O-Overhead}
\label{subsubsec:overhead}

Der Overhead dieses Ansatzes ist klein, aber nicht null. Pro Knoten fallen ein
SSH-Verbindungsaufbau und eine \texttt{scp}-Übertragung an. Übertragen wird jeweils
eine JSON-Datei, deren Größe mit der Zahl der empfangenen Treffer wächst; bei
einer Million geworfener Strahlen liegt sie im Bereich einiger Megabyte. Da diese
acht Übertragungen erst \emph{nach} der Rechenphase stattfinden und nicht mit ihr
überlappen, addieren sie sich voll zur Gesamtlaufzeit. Bei deutlich mehr Knoten
oder größeren Ergebnisdateien würde dieser sequenzielle Einsammel-Schritt zum
Flaschenhals; eine Überlappung von Rechnung und Transfer oder ein kompakteres
Ausgabeformat als das gut lesbare, aber platzintensive JSON wären dann die
naheliegenden Gegenmaßnahmen.

\subsubsection{Messung am Cluster}
\label{subsubsec:messung}

Die konkreten Laufzeiten hängen von der Laborhardware ab und sind am Cluster zu
messen. Um einen aussagekräftigen Speedup zu erhalten, wird in der
Partitionierungs-Betriebsart gemessen: In \texttt{input\_config.json} wird
\texttt{rays\_to\_cast} auf $8\cdot10^{6}/k$ gesetzt, sodass die Gesamtarbeit über
alle $k$ Knoten konstant bleibt. Dann wird \texttt{run\_cluster.sh} für
verschiedene Knotenzahlen mit der Shell-eigenen Zeitmessung ausgeführt:

\begin{lstlisting}[style=clusterlst, language=bash, caption={Messvorschrift für die Laufzeit am Cluster.}]
time ./run_cluster.sh        # Wanduhrzeit fuer den kompletten verteilten Lauf
\end{lstlisting}

Tabelle~\ref{tab:speedup} nimmt die Messwerte auf. Als Referenz dient der Lauf mit
einem einzelnen Knoten; der Speedup ergibt sich als Verhältnis der Wanduhrzeiten.

\begin{table}[htbp]
    \centering
    \caption{Wanduhrzeit und Speedup des verteilten Laufs (am Cluster zu messen).}
    \label{tab:speedup}
    \begin{tabular}{r r r}
        \toprule
        Knoten $k$ & Wanduhrzeit [s] & Speedup $S$ \\
        \midrule
        1 & \dots & $1{,}00$ \\
        2 & \dots & \dots \\
        4 & \dots & \dots \\
        8 & \dots & \dots \\
        \bottomrule
    \end{tabular}
\end{table}

\subsection{Fehlertoleranz und Sicherheit}
\label{subsec:cluster-fehler}

Zwei Querschnittsthemen verdienen eigene Betrachtung. Beim \textbf{Sicherheits}aspekt
ist die Schlüssel-basierte SSH-Authentifizierung zentral: Der Manager benötigt
passwortlosen Zugriff auf jeden Knoten, was über hinterlegte öffentliche Schlüssel
gelöst ist. Der private Schlüssel des Managers ist damit das entscheidende
Geheimnis; wer ihn besitzt, kann auf allen Knoten Code ausführen. In einem
abgeschlossenen Labornetz ist das vertretbar, für einen offeneren Betrieb wäre eine
engere Rechtevergabe angebracht.

Beim Thema \textbf{Fehlertoleranz} zeigt der aktuelle Stand seine deutlichste
Schwäche, die im folgenden Abschnitt zusammen mit den übrigen Grenzen diskutiert
wird. Der Kern des Problems: Ein ausgefallener Knoten wird nicht erkannt, sondern
schlicht übergangen, sodass die Stichprobengröße unbemerkt sinkt.

\subsection{Grenzen des aktuellen Standes und Ausblick}
\label{subsec:cluster-grenzen}

Die Verteilungsschicht ist als funktionsfähiger Prototyp umgesetzt. Für einen
robusten Betrieb bestehen mehrere klar benannte Schwachstellen, die zugleich die
nächsten Entwicklungsschritte vorgeben.

\begin{itemize}
    \item \textbf{Fest verdrahtete Konfiguration.} IP-Adressen, Knotennamen, die
    Knotenzahl und die Ablagepfade stehen direkt im Skript. Ein Umzug auf andere
    Hardware oder eine andere Knotenzahl erfordert Änderungen am Quelltext. Eine
    Auslagerung dieser Werte in eine Konfigurationsdatei wäre der naheliegende
    erste Schritt.

    \item \textbf{Kopierfehler bei Knoten~6.} Im Einsammel-Schritt wird das
    Ergebnis von Knoten~6 versehentlich nach \texttt{ir\_node5.json} kopiert und
    überschreibt damit das Ergebnis von Knoten~5. Effektiv fließen so nur sieben
    statt acht unabhängige Stichproben in die Aggregation ein. Der Fehler ist ein
    Tippfehler im Zielnamen und leicht zu beheben, führt aber vor Augen, wie
    fehleranfällig die von Hand ausgeschriebene Knotenliste ist -- ein weiteres
    Argument für die genannte Auslagerung in eine Schleife über eine
    Konfigurationsliste.

    \item \textbf{Keine Fehlererkennung.} Fällt ein Knoten aus oder liefert er kein
    Ergebnis, schlägt der zugehörige \texttt{scp}-Aufruf fehl, der Gesamtlauf setzt
    sich aber fort und aggregiert stillschweigend über die vorhandenen Dateien. Ein
    Knotenausfall bleibt damit unbemerkt und verfälscht die Stichprobengröße. Eine
    Prüfung der Rückgabecodes und eine Warnung bei fehlenden Ergebnissen wären hier
    angebracht.

    \item \textbf{Manuelle Kopplung an die DSP-Pipeline.} Die Aggregation erzeugt
    \texttt{ir\_merged.json}, die Python-DSP-Pipeline liest jedoch
    \texttt{ir\_output.json}. Der Übergang zwischen verteiltem Lauf und Faltung ist
    derzeit ein manueller Schritt. Eine saubere Integration würde den Cluster-Lauf
    direkt aus der Orchestrierung heraus anstoßen und die zusammengeführte Datei
    automatisch weiterreichen.

    \item \textbf{Statische Verteilung.} Der Replikat-Ansatz mit gleich großen
    Aufträgen ist einfach und robust, setzt aber homogene Hardware voraus und lastet
    unterschiedlich schnelle Knoten nicht unterschiedlich aus. Ein Modell mit
    dynamischer Arbeitsverteilung -- ein Manager vergibt Ray-Kontingente, sobald ein
    Knoten frei wird -- würde heterogene Hardware besser nutzen und die Knotenzahl
    flexibel machen.
\end{itemize}

Alle genannten Punkte betreffen Robustheit und Bedienbarkeit, nicht die
grundsätzliche Korrektheit des Verfahrens. Die verteilte Ausführung liefert eine
gültige, statistisch dichtere Impulsantwort, und das Manager-Worker-Muster über
SSH und \texttt{scp} erfüllt seinen Zweck im Laborbetrieb.

\subsection{Zusammenfassung}
\label{subsec:cluster-zusammenfassung}

Dieses Kapitel hat die verteilte Ausführung der Strahlverfolgung auf dem
Labor-Cluster beschrieben. Weil die Simulation embarrassingly parallel ist, genügt
ein Replikat-Ansatz: Alle acht Knoten rechnen denselben Auftrag mit unabhängigen
Zufallsströmen, was statistisch einer achtfach größeren Simulation entspricht und
den Schätzfehler der Impulsantwort um rund $\sqrt{8}$ senkt -- ein Verhalten, das
das numerische Experiment in Abbildung~\ref{fig:mc-convergence} bestätigt. Ein
Shell-Skript startet die Knoten per SSH parallel, synchronisiert über \texttt{wait}
und sammelt die Ergebnisse per \texttt{scp} ein; \texttt{merge\_results.py}
verkettet die Trefferlisten zu einer gemeinsamen Datei im gewohnten Schema. Der
Prototyp erfüllt seine Aufgabe, weist aber in Konfigurierbarkeit, Fehlererkennung
und der Anbindung an die DSP-Pipeline die diskutierten Schwächen auf, die den
Rahmen für die Weiterentwicklung abstecken.
 im Master.
% Benoetigte Pakete: amsmath, listings, graphicx, booktabs.
% Fuer saubere Optik im Master: fontenc, babel[ngerman], microtype.
% Abbildungen: latexDocs/figures/cluster_arch.png, mc_convergence.png
%
% STANDALONE-VORSCHAU:
%   \documentclass[12pt]{article}
%   \usepackage[utf8]{inputenc}\usepackage[T1]{fontenc}
%   \usepackage[ngerman]{babel}\usepackage{microtype}
%   \usepackage{amsmath}\usepackage{listings}
%   \usepackage{graphicx}\usepackage{booktabs}
%   \begin{document}
%   % =====================================================================
% Kapitel: Verteilte Ausfuehrung auf dem Rechencluster
% Autor: Mier
%
% EINBINDEN: % =====================================================================
% Kapitel: Verteilte Ausfuehrung auf dem Rechencluster
% Autor: Mier
%
% EINBINDEN: \input{Kapitel_Cluster_Mier} im Master.
% Benoetigte Pakete: amsmath, listings, graphicx, booktabs.
% Fuer saubere Optik im Master: fontenc, babel[ngerman], microtype.
% Abbildungen: latexDocs/figures/cluster_arch.png, mc_convergence.png
%
% STANDALONE-VORSCHAU:
%   \documentclass[12pt]{article}
%   \usepackage[utf8]{inputenc}\usepackage[T1]{fontenc}
%   \usepackage[ngerman]{babel}\usepackage{microtype}
%   \usepackage{amsmath}\usepackage{listings}
%   \usepackage{graphicx}\usepackage{booktabs}
%   \begin{document}
%   \input{Kapitel_Cluster_Mier}
%   \end{document}
% =====================================================================

\lstdefinestyle{clusterlst}{
    basicstyle=\ttfamily\small,
    showstringspaces=false,
    breaklines=true,
    numbers=left,
    numberstyle=\tiny,
    frame=single,
    tabsize=2
}

%======================================================================
\section{Verteilte Ausführung auf dem Rechencluster}
\label{sec:cluster}

\subsection{Einordnung und Aufgabenstellung}
\label{subsec:cluster-einordnung}

Die Strahlverfolgung ist der teuerste Schritt der gesamten Pipeline. Für eine
brauchbare Impulsantwort müssen Hunderttausende bis Millionen einzelner
Schallstrahlen durch die Raumgeometrie verfolgt werden, und jeder dieser Strahlen
verlangt mehrere Schnitt- und Reflexionsberechnungen. Der entscheidende Vorteil
dabei: Kein Strahl hängt vom Ergebnis eines anderen ab. Das Problem ist damit
\emph{embarrassingly parallel} und lässt sich ohne Kommunikation zwischen den
Teilrechnungen auf beliebig viele Rechner verteilen.

Dieses Kapitel dokumentiert die Verteilungsschicht des Systems. Sie startet die
Rust-Simulation gleichzeitig auf allen acht Knoten des Labor-Clusters, wartet auf
deren Abschluss, sammelt die Teilergebnisse ein und fügt sie zu einer einzigen,
statistisch dichteren Impulsantwort zusammen. Aus Sicht der übrigen Pipeline
ändert sich dabei nichts: Am Ende steht wieder eine Datei im gewohnten Schema,
die die Signalverarbeitung unverändert weiterverarbeiten kann.

Zur Abgrenzung von den Nachbarkapiteln gehören zwei Punkte. Erstens ist die
\emph{prozessinterne} Nebenläufigkeit der Rust-Engine, also die Verteilung auf
die Kerne eines einzelnen Rechners mittels \texttt{rayon}, an anderer Stelle
beschrieben und bleibt hier außen vor. Sie parallelisiert innerhalb eines
Prozesses; dieses Kapitel parallelisiert über Prozess- und Rechnergrenzen hinweg.
Zweitens geht es nicht um den einzelnen Subprozess-Aufruf aus der Python-Pipeline,
sondern um die Ebene darüber: die gleichzeitige Ausführung auf mehreren physischen
Knoten und die anschließende Aggregation ihrer Ergebnisse.

\subsection{Theoretischer Hintergrund}
\label{subsec:cluster-theorie}

\subsubsection{Monte-Carlo-Strahlverfolgung}
\label{subsubsec:mc}

Die Engine schätzt die Impulsantwort eines Raumes stochastisch. Sie verschickt
zufällig gerichtete Strahlen von der Schallquelle, verfolgt ihre Reflexionen und
registriert, wann und mit welchem Restdruck sie den Empfänger erreichen. Jeder
einzelne Strahl ist eine Zufallsstichprobe; erst die Summe vieler Strahlen ergibt
ein aussagekräftiges Bild der Raumakustik. Solche Monte-Carlo-Verfahren haben eine
charakteristische Eigenschaft: Ihr Schätzfehler sinkt mit der Wurzel der
Stichprobengröße. Verdoppelt man die Zahl der Strahlen, halbiert sich der Fehler
also nicht, sondern sinkt nur um den Faktor $\sqrt{2}$. Genau diese langsame
Konvergenz macht große Strahlzahlen nötig und motiviert die Verteilung auf mehrere
Rechner.

\subsubsection{Zwei Wege der Verteilung}
\label{subsubsec:strategien}

Ein Monte-Carlo-Verfahren lässt sich auf zwei grundsätzlich verschiedene Arten
verteilen. Tabelle~\ref{tab:strategien} stellt sie gegenüber.

\begin{table}[htbp]
    \centering
    \caption{Zwei Strategien zur Verteilung der Strahlverfolgung.}
    \label{tab:strategien}
    \begin{tabular}{p{3.2cm} p{5.2cm} p{4.2cm}}
        \toprule
        & \textbf{Work-Partitioning} & \textbf{Replikat-Parallelismus} \\
        \midrule
        Prinzip & Gesamtzahl der Strahlen wird in disjunkte Anteile zerlegt & Jeder Knoten rechnet denselben Auftrag mit eigenem Zufallsstrom \\
        Koordination & Zuteilung der Anteile nötig & keine, jeder Knoten ist autark \\
        Zufallszahlen & ein globaler Strom, aufgeteilt & unabhängige Ströme pro Knoten \\
        Robustheit & Ausfall trifft einen festen Anteil & Ausfall kostet nur Stichprobengröße \\
        \bottomrule
    \end{tabular}
\end{table}

Das umgesetzte System verwendet den zweiten Weg. Alle acht Knoten starten die
Simulation mit derselben Konfigurationsdatei \texttt{input\_config.json} und
werfen dennoch jeweils andere Strahlen, weil jeder Prozess seinen eigenen
Zufallsstrom besitzt. Der Ansatz ist bewusst schlicht gewählt: Er braucht keine
zentrale Zuteilung und keinen Austausch zwischen den Knoten, sondern nur, dass am
Ende alle Ergebnisse eingesammelt werden.

\subsubsection{Unabhängigkeit der Zufallsströme}
\label{subsubsec:rng}

Die Korrektheit des Replikat-Ansatzes steht und fällt mit einer Bedingung: Die
Zufallszahlen auf den Knoten müssen unabhängig sein. Würde jeder Knoten dieselbe
Folge zufälliger Richtungen erzeugen, so lieferten alle acht exakt dasselbe
Ergebnis, und das Zusammenführen brächte keinen Erkenntnisgewinn, sondern nur
achtfach dieselben Strahlen. Die Rust-Engine schließt das aus, indem sie ihren
Zufallsstrom pro Thread bzw. Prozess frisch aus der Systementropie zieht:

\begin{lstlisting}[style=clusterlst, language=Rust, caption={Initialisierung des Zufallsstroms pro Prozess (\texttt{simulation.rs}).}]
let mut rng = rand::rng();                       // Omnidirectional, random rays
let random_angle: f32 = rng.random_range(0.0..(2.0 * PI));
let direction = Vec2::new(random_angle.cos(), random_angle.sin());
\end{lstlisting}

Weil jeder der acht Prozesse auf seinem Knoten einen anderen Startzustand erhält,
erzeugt jeder eine andere Stichprobe. Acht Knoten mit je einer Million Strahlen
liefern damit acht Millionen \emph{unabhängige} Stichproben derselben Raumakustik.

\subsubsection{Statistische Konvergenz}
\label{subsubsec:konvergenz}

Statistisch verhält sich der verteilte Lauf wie eine einzige, achtfach größere
Simulation. Für einen Monte-Carlo-Schätzer aus $N$ unabhängigen Teilstichproben
gleicher Größe sinkt die Standardabweichung des Mittelwerts mit
\[
    \sigma_N \;=\; \frac{\sigma_1}{\sqrt{N}} .
\]
Acht Knoten reduzieren die Streuung der geschätzten Impulsantwort demnach um den
Faktor $\sqrt{8} \approx 2{,}83$ gegenüber einem einzelnen Knoten.
Abbildung~\ref{fig:mc-convergence} bestätigt dieses Verhalten in einem numerischen
Experiment: Aufgetragen ist der mittlere quadratische Fehler der geschätzten
Impulsantwort über der Knotenzahl. Die Messpunkte folgen sauber der
$1/\sqrt{N}$-Referenzkurve; von einem auf acht Knoten sinkt der Fehler um den
Faktor $2{,}78$, praktisch exakt der theoretisch erwartete Wert.

\begin{figure}[htbp]
    \centering
    \includegraphics[width=0.8\textwidth]{figures/mc_convergence}
    \caption{Mittlerer quadratischer Fehler der geschätzten Impulsantwort über der
    Knotenzahl $N$ (je $5\cdot10^4$ Strahlen, doppelt logarithmisch). Die
    Messpunkte folgen der $1/\sqrt{N}$-Referenz; acht Knoten senken den Fehler um
    den Faktor $2{,}78 \approx \sqrt{8}$.}
    \label{fig:mc-convergence}
\end{figure}

Bemerkenswert ist die Kehrseite dieser Kurve: Um den Fehler zu halbieren, braucht
es die vierfache Knotenzahl. Der Nutzen zusätzlicher Knoten nimmt also ab, was bei
der Wahl der Clustergröße zu berücksichtigen ist. Für die hier angestrebte
Genauigkeit sind acht Knoten ein sinnvoller Kompromiss zwischen Aufwand und
Qualität.

\subsection{Architektur des Clusters}
\label{subsec:cluster-arch}

Das System folgt einem \emph{Manager-Worker}-Muster
(Abbildung~\ref{fig:cluster-arch}). Ein Manager-Knoten koordiniert den Ablauf,
acht Worker-Knoten führen die eigentliche Simulation aus. Der Manager selbst
rechnet nicht mit, sondern startet die Worker, wartet auf sie und führt am Ende
die Ergebnisse zusammen.

\begin{figure}[htbp]
    \centering
    \includegraphics[width=0.82\textwidth]{figures/cluster_arch}
    \caption{Manager-Worker-Architektur. Der Manager startet die acht Knoten
    parallel per \texttt{ssh}, holt nach Abschluss jede Ergebnisdatei per
    \texttt{scp} zurück und führt sie mit \texttt{merge\_results.py} zusammen.}
    \label{fig:cluster-arch}
\end{figure}

Die Kommunikation läuft vollständig über Standard-Werkzeuge des Betriebssystems.
\texttt{ssh} startet die Prozesse aus der Ferne, \texttt{scp} überträgt die
Ergebnisdateien zurück. Ein eigenes Netzwerkprotokoll oder ein
Message-Passing-Framework kommt nicht zum Einsatz. Für einen Stapelbetrieb, bei
dem jeder Knoten genau einmal startet, eine Weile rechnet und ein einziges
Ergebnis zurückliefert, wäre solcher Aufwand nicht gerechtfertigt: Es gibt keine
laufende Zustandssynchronisation und keinen Nachrichtenaustausch während der
Rechnung, den ein Framework verwalten müsste.

Damit das funktioniert, müssen drei Voraussetzungen erfüllt sein. Auf jedem Knoten
liegt dieselbe vorkompilierte Binärdatei \texttt{rust\_service} samt der
Konfiguration \texttt{input\_config.json}. Der Manager kann sich per SSH ohne
interaktive Passworteingabe anmelden, was eine Schlüssel-basierte
Authentifizierung voraussetzt. Und die Knoten sind über feste Adressen im
Labornetz erreichbar.

\subsection{Ablaufsteuerung: \texttt{run\_cluster.sh}}
\label{subsec:cluster-run}

Das Skript \texttt{run\_cluster.sh} steuert den verteilten Lauf in drei Phasen:
paralleles Starten, Warten und Einsammeln.

\begin{lstlisting}[style=clusterlst, language=bash, caption={Verteilter Start und Einsammeln der Ergebnisse (gekürzt).}]
# Phase 1: alle Knoten gleichzeitig starten (Hintergrundprozesse)
ssh node1@139.6.65.29 "./rust_service --config-file input_config.json" &
ssh node2@139.6.65.18 "./rust_service --config-file input_config.json" &
# ... node3 bis node8 analog ...

wait                       # Phase 2: auf Abschluss ALLER Knoten warten

# Phase 3: Ergebnis jedes Knotens zum Manager kopieren
scp node1@139.6.65.29:~/ir_output.json ~/ir_node1.json
scp node2@139.6.65.18:~/ir_output.json ~/ir_node2.json
# ... node3 bis node8 analog ...
\end{lstlisting}

Zwei Eigenschaften der Shell nutzt der Ablauf gezielt. Das kaufmännische Und
(\texttt{\&}) am Zeilenende schickt jeden \texttt{ssh}-Aufruf in den Hintergrund.
Die acht Simulationen laufen dadurch nahezu gleichzeitig an, statt nacheinander
abgearbeitet zu werden. Das darauffolgende \texttt{wait} blockiert, bis
\emph{alle} Hintergrundprozesse beendet sind, und bildet so eine
Synchronisationsbarriere zwischen der Rechen- und der Einsammel-Phase. Erst wenn
der letzte Knoten fertig ist, beginnt \texttt{scp}, von jedem Knoten dessen
\texttt{ir\_output.json} zu holen und unter einem knotenspezifischen Namen auf dem
Manager abzulegen.

Diese Barriere hat eine wichtige Konsequenz für die Laufzeit: Der Gesamtlauf ist
so langsam wie der \emph{langsamste} Knoten. Solange die Hardware homogen ist,
fällt das kaum ins Gewicht; auf ungleichen Knoten würde ein einzelner langsamer
Rechner den ganzen Lauf ausbremsen, während die übrigen sieben bereits fertig
warten. Dieses Straggler-Problem ist der zentrale Grund, warum ein statischer
Replikat-Ansatz mit gleich großen Aufträgen homogene Hardware voraussetzt.

\subsection{Aggregation der Teilergebnisse: \texttt{merge\_results.py}}
\label{subsec:cluster-merge}

Sobald alle Teilergebnisse auf dem Manager liegen, führt \texttt{merge\_results.py}
sie zu einer einzigen Impulsantwort zusammen. Jeder Knoten liefert dasselbe
Ausgabeschema, nämlich Metadaten sowie zwei positionsgleiche Listen aus
Verzögerungen und Drücken. Die Aggregation besteht deshalb im Kern aus einer
schlichten Verkettung dieser Listen:

\begin{lstlisting}[style=clusterlst, language=Python, caption={Zusammenführen der Knoten-Ergebnisse, gekürzt (\texttt{merge\_results.py}).}]
all_delays, all_pressures = [], []
for f in sorted(glob.glob("/home/manager/ir_node*.json")):
    with open(f) as file:
        data = json.load(file)
        all_delays    += data["hits"]["delays_seconds"]
        all_pressures += data["hits"]["pressures"]

merged = {
    "metadata": {"sample_rate": 44100, "rays_cast": 1000000 * 8,
                 "rays_received": len(all_delays), "room_name": "MVP_Test_Room_1"},
    "hits": {"delays_seconds": all_delays, "pressures": all_pressures},
}
json.dump(merged, open("/home/manager/ir_merged.json", "w"), indent=2)
\end{lstlisting}

Dass die Verkettung hier die richtige Operation ist, folgt direkt aus dem
Replikat-Modell. Weil die einzelnen Läufe unabhängige Stichproben desselben Raums
sind, ist die Vereinigung aller Treffer statistisch identisch mit dem Ergebnis
eines einzigen Laufs mit acht Millionen Strahlen. Es findet keine Mittelung und
keine Gewichtung statt, sondern nur ein Aneinanderhängen. Das Feld
\texttt{rays\_cast} ist im Skript fest auf $8 \cdot 10^{6}$ gesetzt; solange der in
Abschnitt~\ref{subsec:cluster-grenzen} beschriebene Kopierfehler bei Knoten~6
besteht, fließen tatsächlich nur sieben Knoten in die Aggregation ein, und dieser
Metadatenwert überschätzt die reale Strahlzahl. \texttt{rays\_received} zählt
dagegen die tatsächlich empfangenen Treffer. Die
zusammengeführte Datei trägt exakt dasselbe Schema wie die Ausgabe eines einzelnen
Knotens und ist damit ohne weitere Anpassung für die DSP-Pipeline verwendbar.

\subsection{Datenfluss im Überblick}
\label{subsec:cluster-datenfluss}

Der gesamte verteilte Ablauf lässt sich in fünf Stufen zusammenfassen:

\begin{enumerate}
    \item Der Manager startet per SSH auf jedem der acht Knoten
    \texttt{rust\_service} mit identischer Konfiguration, alle parallel.
    \item Jeder Knoten simuliert unabhängig eine Million Strahlen und schreibt
    lokal seine \texttt{ir\_output.json}.
    \item \texttt{wait} synchronisiert: Die Einsammel-Phase beginnt erst, wenn der
    letzte Knoten fertig ist.
    \item \texttt{scp} überträgt jede \texttt{ir\_output.json} als
    \texttt{ir\_node$k$.json} zum Manager.
    \item \texttt{merge\_results.py} verkettet alle Trefferlisten zu einer
    \texttt{ir\_merged.json}, die anschließend die DSP-Pipeline verarbeitet.
\end{enumerate}

\subsection{Leistungsbetrachtung}
\label{subsec:cluster-leistung}

\subsubsection{Erwarteter Speedup}
\label{subsubsec:speedup}

Bei der Leistung sind zwei Sichtweisen zu trennen. In der umgesetzten
Replikat-Betriebsart rechnet jeder Knoten dieselbe Strahlzahl, sodass acht Knoten
in nahezu derselben Wanduhrzeit die achtfache Stichprobe liefern. Der Gewinn ist
hier keine kürzere Laufzeit, sondern die höhere Genauigkeit aus
Abbildung~\ref{fig:mc-convergence}; man spricht vom schwachen Skalieren im Sinne
von Gustafson, bei dem mit der Rechnerzahl die bearbeitete Problemgröße wächst,
nicht die Geschwindigkeit einer festen Aufgabe. Will man dagegen eine feste
Gesamtstrahlzahl in kürzerer Zeit erreichen, so verteilt man diese auf die Knoten,
sodass jeder nur $8\cdot10^{6}/k$ Strahlen rechnet. Erst in dieser
Partitionierungs-Betriebsart greift das Gesetz von Amdahl. Bezeichnet $p$ den parallelen
Anteil der Arbeit, hier die reine Strahlverfolgung, und $1-p$ den seriellen Rest
aus Prozessstart, Ergebnistransfer und Aggregation, so beträgt der theoretische
Speedup
\[
    S(k) \;=\; \frac{1}{(1-p) + \dfrac{p}{k}} .
\]
Da die Strahlverfolgung die Laufzeit klar dominiert, liegt $p$ nahe bei eins, und
$S(k)$ steigt bis $k=8$ nahezu linear. In der Praxis bleibt der reale Speedup
darunter, und zwar aus drei Gründen: dem bereits erwähnten Straggler-Effekt (der
Lauf ist so schnell wie der langsamste Knoten), dem Overhead für den
SSH-Verbindungsaufbau sowie der Übertragung und Zusammenführung der
Ergebnisdateien.

\subsubsection{Kommunikations- und I/O-Overhead}
\label{subsubsec:overhead}

Der Overhead dieses Ansatzes ist klein, aber nicht null. Pro Knoten fallen ein
SSH-Verbindungsaufbau und eine \texttt{scp}-Übertragung an. Übertragen wird jeweils
eine JSON-Datei, deren Größe mit der Zahl der empfangenen Treffer wächst; bei
einer Million geworfener Strahlen liegt sie im Bereich einiger Megabyte. Da diese
acht Übertragungen erst \emph{nach} der Rechenphase stattfinden und nicht mit ihr
überlappen, addieren sie sich voll zur Gesamtlaufzeit. Bei deutlich mehr Knoten
oder größeren Ergebnisdateien würde dieser sequenzielle Einsammel-Schritt zum
Flaschenhals; eine Überlappung von Rechnung und Transfer oder ein kompakteres
Ausgabeformat als das gut lesbare, aber platzintensive JSON wären dann die
naheliegenden Gegenmaßnahmen.

\subsubsection{Messung am Cluster}
\label{subsubsec:messung}

Die konkreten Laufzeiten hängen von der Laborhardware ab und sind am Cluster zu
messen. Um einen aussagekräftigen Speedup zu erhalten, wird in der
Partitionierungs-Betriebsart gemessen: In \texttt{input\_config.json} wird
\texttt{rays\_to\_cast} auf $8\cdot10^{6}/k$ gesetzt, sodass die Gesamtarbeit über
alle $k$ Knoten konstant bleibt. Dann wird \texttt{run\_cluster.sh} für
verschiedene Knotenzahlen mit der Shell-eigenen Zeitmessung ausgeführt:

\begin{lstlisting}[style=clusterlst, language=bash, caption={Messvorschrift für die Laufzeit am Cluster.}]
time ./run_cluster.sh        # Wanduhrzeit fuer den kompletten verteilten Lauf
\end{lstlisting}

Tabelle~\ref{tab:speedup} nimmt die Messwerte auf. Als Referenz dient der Lauf mit
einem einzelnen Knoten; der Speedup ergibt sich als Verhältnis der Wanduhrzeiten.

\begin{table}[htbp]
    \centering
    \caption{Wanduhrzeit und Speedup des verteilten Laufs (am Cluster zu messen).}
    \label{tab:speedup}
    \begin{tabular}{r r r}
        \toprule
        Knoten $k$ & Wanduhrzeit [s] & Speedup $S$ \\
        \midrule
        1 & \dots & $1{,}00$ \\
        2 & \dots & \dots \\
        4 & \dots & \dots \\
        8 & \dots & \dots \\
        \bottomrule
    \end{tabular}
\end{table}

\subsection{Fehlertoleranz und Sicherheit}
\label{subsec:cluster-fehler}

Zwei Querschnittsthemen verdienen eigene Betrachtung. Beim \textbf{Sicherheits}aspekt
ist die Schlüssel-basierte SSH-Authentifizierung zentral: Der Manager benötigt
passwortlosen Zugriff auf jeden Knoten, was über hinterlegte öffentliche Schlüssel
gelöst ist. Der private Schlüssel des Managers ist damit das entscheidende
Geheimnis; wer ihn besitzt, kann auf allen Knoten Code ausführen. In einem
abgeschlossenen Labornetz ist das vertretbar, für einen offeneren Betrieb wäre eine
engere Rechtevergabe angebracht.

Beim Thema \textbf{Fehlertoleranz} zeigt der aktuelle Stand seine deutlichste
Schwäche, die im folgenden Abschnitt zusammen mit den übrigen Grenzen diskutiert
wird. Der Kern des Problems: Ein ausgefallener Knoten wird nicht erkannt, sondern
schlicht übergangen, sodass die Stichprobengröße unbemerkt sinkt.

\subsection{Grenzen des aktuellen Standes und Ausblick}
\label{subsec:cluster-grenzen}

Die Verteilungsschicht ist als funktionsfähiger Prototyp umgesetzt. Für einen
robusten Betrieb bestehen mehrere klar benannte Schwachstellen, die zugleich die
nächsten Entwicklungsschritte vorgeben.

\begin{itemize}
    \item \textbf{Fest verdrahtete Konfiguration.} IP-Adressen, Knotennamen, die
    Knotenzahl und die Ablagepfade stehen direkt im Skript. Ein Umzug auf andere
    Hardware oder eine andere Knotenzahl erfordert Änderungen am Quelltext. Eine
    Auslagerung dieser Werte in eine Konfigurationsdatei wäre der naheliegende
    erste Schritt.

    \item \textbf{Kopierfehler bei Knoten~6.} Im Einsammel-Schritt wird das
    Ergebnis von Knoten~6 versehentlich nach \texttt{ir\_node5.json} kopiert und
    überschreibt damit das Ergebnis von Knoten~5. Effektiv fließen so nur sieben
    statt acht unabhängige Stichproben in die Aggregation ein. Der Fehler ist ein
    Tippfehler im Zielnamen und leicht zu beheben, führt aber vor Augen, wie
    fehleranfällig die von Hand ausgeschriebene Knotenliste ist -- ein weiteres
    Argument für die genannte Auslagerung in eine Schleife über eine
    Konfigurationsliste.

    \item \textbf{Keine Fehlererkennung.} Fällt ein Knoten aus oder liefert er kein
    Ergebnis, schlägt der zugehörige \texttt{scp}-Aufruf fehl, der Gesamtlauf setzt
    sich aber fort und aggregiert stillschweigend über die vorhandenen Dateien. Ein
    Knotenausfall bleibt damit unbemerkt und verfälscht die Stichprobengröße. Eine
    Prüfung der Rückgabecodes und eine Warnung bei fehlenden Ergebnissen wären hier
    angebracht.

    \item \textbf{Manuelle Kopplung an die DSP-Pipeline.} Die Aggregation erzeugt
    \texttt{ir\_merged.json}, die Python-DSP-Pipeline liest jedoch
    \texttt{ir\_output.json}. Der Übergang zwischen verteiltem Lauf und Faltung ist
    derzeit ein manueller Schritt. Eine saubere Integration würde den Cluster-Lauf
    direkt aus der Orchestrierung heraus anstoßen und die zusammengeführte Datei
    automatisch weiterreichen.

    \item \textbf{Statische Verteilung.} Der Replikat-Ansatz mit gleich großen
    Aufträgen ist einfach und robust, setzt aber homogene Hardware voraus und lastet
    unterschiedlich schnelle Knoten nicht unterschiedlich aus. Ein Modell mit
    dynamischer Arbeitsverteilung -- ein Manager vergibt Ray-Kontingente, sobald ein
    Knoten frei wird -- würde heterogene Hardware besser nutzen und die Knotenzahl
    flexibel machen.
\end{itemize}

Alle genannten Punkte betreffen Robustheit und Bedienbarkeit, nicht die
grundsätzliche Korrektheit des Verfahrens. Die verteilte Ausführung liefert eine
gültige, statistisch dichtere Impulsantwort, und das Manager-Worker-Muster über
SSH und \texttt{scp} erfüllt seinen Zweck im Laborbetrieb.

\subsection{Zusammenfassung}
\label{subsec:cluster-zusammenfassung}

Dieses Kapitel hat die verteilte Ausführung der Strahlverfolgung auf dem
Labor-Cluster beschrieben. Weil die Simulation embarrassingly parallel ist, genügt
ein Replikat-Ansatz: Alle acht Knoten rechnen denselben Auftrag mit unabhängigen
Zufallsströmen, was statistisch einer achtfach größeren Simulation entspricht und
den Schätzfehler der Impulsantwort um rund $\sqrt{8}$ senkt -- ein Verhalten, das
das numerische Experiment in Abbildung~\ref{fig:mc-convergence} bestätigt. Ein
Shell-Skript startet die Knoten per SSH parallel, synchronisiert über \texttt{wait}
und sammelt die Ergebnisse per \texttt{scp} ein; \texttt{merge\_results.py}
verkettet die Trefferlisten zu einer gemeinsamen Datei im gewohnten Schema. Der
Prototyp erfüllt seine Aufgabe, weist aber in Konfigurierbarkeit, Fehlererkennung
und der Anbindung an die DSP-Pipeline die diskutierten Schwächen auf, die den
Rahmen für die Weiterentwicklung abstecken.
 im Master.
% Benoetigte Pakete: amsmath, listings, graphicx, booktabs.
% Fuer saubere Optik im Master: fontenc, babel[ngerman], microtype.
% Abbildungen: latexDocs/figures/cluster_arch.png, mc_convergence.png
%
% STANDALONE-VORSCHAU:
%   \documentclass[12pt]{article}
%   \usepackage[utf8]{inputenc}\usepackage[T1]{fontenc}
%   \usepackage[ngerman]{babel}\usepackage{microtype}
%   \usepackage{amsmath}\usepackage{listings}
%   \usepackage{graphicx}\usepackage{booktabs}
%   \begin{document}
%   % =====================================================================
% Kapitel: Verteilte Ausfuehrung auf dem Rechencluster
% Autor: Mier
%
% EINBINDEN: \input{Kapitel_Cluster_Mier} im Master.
% Benoetigte Pakete: amsmath, listings, graphicx, booktabs.
% Fuer saubere Optik im Master: fontenc, babel[ngerman], microtype.
% Abbildungen: latexDocs/figures/cluster_arch.png, mc_convergence.png
%
% STANDALONE-VORSCHAU:
%   \documentclass[12pt]{article}
%   \usepackage[utf8]{inputenc}\usepackage[T1]{fontenc}
%   \usepackage[ngerman]{babel}\usepackage{microtype}
%   \usepackage{amsmath}\usepackage{listings}
%   \usepackage{graphicx}\usepackage{booktabs}
%   \begin{document}
%   \input{Kapitel_Cluster_Mier}
%   \end{document}
% =====================================================================

\lstdefinestyle{clusterlst}{
    basicstyle=\ttfamily\small,
    showstringspaces=false,
    breaklines=true,
    numbers=left,
    numberstyle=\tiny,
    frame=single,
    tabsize=2
}

%======================================================================
\section{Verteilte Ausführung auf dem Rechencluster}
\label{sec:cluster}

\subsection{Einordnung und Aufgabenstellung}
\label{subsec:cluster-einordnung}

Die Strahlverfolgung ist der teuerste Schritt der gesamten Pipeline. Für eine
brauchbare Impulsantwort müssen Hunderttausende bis Millionen einzelner
Schallstrahlen durch die Raumgeometrie verfolgt werden, und jeder dieser Strahlen
verlangt mehrere Schnitt- und Reflexionsberechnungen. Der entscheidende Vorteil
dabei: Kein Strahl hängt vom Ergebnis eines anderen ab. Das Problem ist damit
\emph{embarrassingly parallel} und lässt sich ohne Kommunikation zwischen den
Teilrechnungen auf beliebig viele Rechner verteilen.

Dieses Kapitel dokumentiert die Verteilungsschicht des Systems. Sie startet die
Rust-Simulation gleichzeitig auf allen acht Knoten des Labor-Clusters, wartet auf
deren Abschluss, sammelt die Teilergebnisse ein und fügt sie zu einer einzigen,
statistisch dichteren Impulsantwort zusammen. Aus Sicht der übrigen Pipeline
ändert sich dabei nichts: Am Ende steht wieder eine Datei im gewohnten Schema,
die die Signalverarbeitung unverändert weiterverarbeiten kann.

Zur Abgrenzung von den Nachbarkapiteln gehören zwei Punkte. Erstens ist die
\emph{prozessinterne} Nebenläufigkeit der Rust-Engine, also die Verteilung auf
die Kerne eines einzelnen Rechners mittels \texttt{rayon}, an anderer Stelle
beschrieben und bleibt hier außen vor. Sie parallelisiert innerhalb eines
Prozesses; dieses Kapitel parallelisiert über Prozess- und Rechnergrenzen hinweg.
Zweitens geht es nicht um den einzelnen Subprozess-Aufruf aus der Python-Pipeline,
sondern um die Ebene darüber: die gleichzeitige Ausführung auf mehreren physischen
Knoten und die anschließende Aggregation ihrer Ergebnisse.

\subsection{Theoretischer Hintergrund}
\label{subsec:cluster-theorie}

\subsubsection{Monte-Carlo-Strahlverfolgung}
\label{subsubsec:mc}

Die Engine schätzt die Impulsantwort eines Raumes stochastisch. Sie verschickt
zufällig gerichtete Strahlen von der Schallquelle, verfolgt ihre Reflexionen und
registriert, wann und mit welchem Restdruck sie den Empfänger erreichen. Jeder
einzelne Strahl ist eine Zufallsstichprobe; erst die Summe vieler Strahlen ergibt
ein aussagekräftiges Bild der Raumakustik. Solche Monte-Carlo-Verfahren haben eine
charakteristische Eigenschaft: Ihr Schätzfehler sinkt mit der Wurzel der
Stichprobengröße. Verdoppelt man die Zahl der Strahlen, halbiert sich der Fehler
also nicht, sondern sinkt nur um den Faktor $\sqrt{2}$. Genau diese langsame
Konvergenz macht große Strahlzahlen nötig und motiviert die Verteilung auf mehrere
Rechner.

\subsubsection{Zwei Wege der Verteilung}
\label{subsubsec:strategien}

Ein Monte-Carlo-Verfahren lässt sich auf zwei grundsätzlich verschiedene Arten
verteilen. Tabelle~\ref{tab:strategien} stellt sie gegenüber.

\begin{table}[htbp]
    \centering
    \caption{Zwei Strategien zur Verteilung der Strahlverfolgung.}
    \label{tab:strategien}
    \begin{tabular}{p{3.2cm} p{5.2cm} p{4.2cm}}
        \toprule
        & \textbf{Work-Partitioning} & \textbf{Replikat-Parallelismus} \\
        \midrule
        Prinzip & Gesamtzahl der Strahlen wird in disjunkte Anteile zerlegt & Jeder Knoten rechnet denselben Auftrag mit eigenem Zufallsstrom \\
        Koordination & Zuteilung der Anteile nötig & keine, jeder Knoten ist autark \\
        Zufallszahlen & ein globaler Strom, aufgeteilt & unabhängige Ströme pro Knoten \\
        Robustheit & Ausfall trifft einen festen Anteil & Ausfall kostet nur Stichprobengröße \\
        \bottomrule
    \end{tabular}
\end{table}

Das umgesetzte System verwendet den zweiten Weg. Alle acht Knoten starten die
Simulation mit derselben Konfigurationsdatei \texttt{input\_config.json} und
werfen dennoch jeweils andere Strahlen, weil jeder Prozess seinen eigenen
Zufallsstrom besitzt. Der Ansatz ist bewusst schlicht gewählt: Er braucht keine
zentrale Zuteilung und keinen Austausch zwischen den Knoten, sondern nur, dass am
Ende alle Ergebnisse eingesammelt werden.

\subsubsection{Unabhängigkeit der Zufallsströme}
\label{subsubsec:rng}

Die Korrektheit des Replikat-Ansatzes steht und fällt mit einer Bedingung: Die
Zufallszahlen auf den Knoten müssen unabhängig sein. Würde jeder Knoten dieselbe
Folge zufälliger Richtungen erzeugen, so lieferten alle acht exakt dasselbe
Ergebnis, und das Zusammenführen brächte keinen Erkenntnisgewinn, sondern nur
achtfach dieselben Strahlen. Die Rust-Engine schließt das aus, indem sie ihren
Zufallsstrom pro Thread bzw. Prozess frisch aus der Systementropie zieht:

\begin{lstlisting}[style=clusterlst, language=Rust, caption={Initialisierung des Zufallsstroms pro Prozess (\texttt{simulation.rs}).}]
let mut rng = rand::rng();                       // Omnidirectional, random rays
let random_angle: f32 = rng.random_range(0.0..(2.0 * PI));
let direction = Vec2::new(random_angle.cos(), random_angle.sin());
\end{lstlisting}

Weil jeder der acht Prozesse auf seinem Knoten einen anderen Startzustand erhält,
erzeugt jeder eine andere Stichprobe. Acht Knoten mit je einer Million Strahlen
liefern damit acht Millionen \emph{unabhängige} Stichproben derselben Raumakustik.

\subsubsection{Statistische Konvergenz}
\label{subsubsec:konvergenz}

Statistisch verhält sich der verteilte Lauf wie eine einzige, achtfach größere
Simulation. Für einen Monte-Carlo-Schätzer aus $N$ unabhängigen Teilstichproben
gleicher Größe sinkt die Standardabweichung des Mittelwerts mit
\[
    \sigma_N \;=\; \frac{\sigma_1}{\sqrt{N}} .
\]
Acht Knoten reduzieren die Streuung der geschätzten Impulsantwort demnach um den
Faktor $\sqrt{8} \approx 2{,}83$ gegenüber einem einzelnen Knoten.
Abbildung~\ref{fig:mc-convergence} bestätigt dieses Verhalten in einem numerischen
Experiment: Aufgetragen ist der mittlere quadratische Fehler der geschätzten
Impulsantwort über der Knotenzahl. Die Messpunkte folgen sauber der
$1/\sqrt{N}$-Referenzkurve; von einem auf acht Knoten sinkt der Fehler um den
Faktor $2{,}78$, praktisch exakt der theoretisch erwartete Wert.

\begin{figure}[htbp]
    \centering
    \includegraphics[width=0.8\textwidth]{figures/mc_convergence}
    \caption{Mittlerer quadratischer Fehler der geschätzten Impulsantwort über der
    Knotenzahl $N$ (je $5\cdot10^4$ Strahlen, doppelt logarithmisch). Die
    Messpunkte folgen der $1/\sqrt{N}$-Referenz; acht Knoten senken den Fehler um
    den Faktor $2{,}78 \approx \sqrt{8}$.}
    \label{fig:mc-convergence}
\end{figure}

Bemerkenswert ist die Kehrseite dieser Kurve: Um den Fehler zu halbieren, braucht
es die vierfache Knotenzahl. Der Nutzen zusätzlicher Knoten nimmt also ab, was bei
der Wahl der Clustergröße zu berücksichtigen ist. Für die hier angestrebte
Genauigkeit sind acht Knoten ein sinnvoller Kompromiss zwischen Aufwand und
Qualität.

\subsection{Architektur des Clusters}
\label{subsec:cluster-arch}

Das System folgt einem \emph{Manager-Worker}-Muster
(Abbildung~\ref{fig:cluster-arch}). Ein Manager-Knoten koordiniert den Ablauf,
acht Worker-Knoten führen die eigentliche Simulation aus. Der Manager selbst
rechnet nicht mit, sondern startet die Worker, wartet auf sie und führt am Ende
die Ergebnisse zusammen.

\begin{figure}[htbp]
    \centering
    \includegraphics[width=0.82\textwidth]{figures/cluster_arch}
    \caption{Manager-Worker-Architektur. Der Manager startet die acht Knoten
    parallel per \texttt{ssh}, holt nach Abschluss jede Ergebnisdatei per
    \texttt{scp} zurück und führt sie mit \texttt{merge\_results.py} zusammen.}
    \label{fig:cluster-arch}
\end{figure}

Die Kommunikation läuft vollständig über Standard-Werkzeuge des Betriebssystems.
\texttt{ssh} startet die Prozesse aus der Ferne, \texttt{scp} überträgt die
Ergebnisdateien zurück. Ein eigenes Netzwerkprotokoll oder ein
Message-Passing-Framework kommt nicht zum Einsatz. Für einen Stapelbetrieb, bei
dem jeder Knoten genau einmal startet, eine Weile rechnet und ein einziges
Ergebnis zurückliefert, wäre solcher Aufwand nicht gerechtfertigt: Es gibt keine
laufende Zustandssynchronisation und keinen Nachrichtenaustausch während der
Rechnung, den ein Framework verwalten müsste.

Damit das funktioniert, müssen drei Voraussetzungen erfüllt sein. Auf jedem Knoten
liegt dieselbe vorkompilierte Binärdatei \texttt{rust\_service} samt der
Konfiguration \texttt{input\_config.json}. Der Manager kann sich per SSH ohne
interaktive Passworteingabe anmelden, was eine Schlüssel-basierte
Authentifizierung voraussetzt. Und die Knoten sind über feste Adressen im
Labornetz erreichbar.

\subsection{Ablaufsteuerung: \texttt{run\_cluster.sh}}
\label{subsec:cluster-run}

Das Skript \texttt{run\_cluster.sh} steuert den verteilten Lauf in drei Phasen:
paralleles Starten, Warten und Einsammeln.

\begin{lstlisting}[style=clusterlst, language=bash, caption={Verteilter Start und Einsammeln der Ergebnisse (gekürzt).}]
# Phase 1: alle Knoten gleichzeitig starten (Hintergrundprozesse)
ssh node1@139.6.65.29 "./rust_service --config-file input_config.json" &
ssh node2@139.6.65.18 "./rust_service --config-file input_config.json" &
# ... node3 bis node8 analog ...

wait                       # Phase 2: auf Abschluss ALLER Knoten warten

# Phase 3: Ergebnis jedes Knotens zum Manager kopieren
scp node1@139.6.65.29:~/ir_output.json ~/ir_node1.json
scp node2@139.6.65.18:~/ir_output.json ~/ir_node2.json
# ... node3 bis node8 analog ...
\end{lstlisting}

Zwei Eigenschaften der Shell nutzt der Ablauf gezielt. Das kaufmännische Und
(\texttt{\&}) am Zeilenende schickt jeden \texttt{ssh}-Aufruf in den Hintergrund.
Die acht Simulationen laufen dadurch nahezu gleichzeitig an, statt nacheinander
abgearbeitet zu werden. Das darauffolgende \texttt{wait} blockiert, bis
\emph{alle} Hintergrundprozesse beendet sind, und bildet so eine
Synchronisationsbarriere zwischen der Rechen- und der Einsammel-Phase. Erst wenn
der letzte Knoten fertig ist, beginnt \texttt{scp}, von jedem Knoten dessen
\texttt{ir\_output.json} zu holen und unter einem knotenspezifischen Namen auf dem
Manager abzulegen.

Diese Barriere hat eine wichtige Konsequenz für die Laufzeit: Der Gesamtlauf ist
so langsam wie der \emph{langsamste} Knoten. Solange die Hardware homogen ist,
fällt das kaum ins Gewicht; auf ungleichen Knoten würde ein einzelner langsamer
Rechner den ganzen Lauf ausbremsen, während die übrigen sieben bereits fertig
warten. Dieses Straggler-Problem ist der zentrale Grund, warum ein statischer
Replikat-Ansatz mit gleich großen Aufträgen homogene Hardware voraussetzt.

\subsection{Aggregation der Teilergebnisse: \texttt{merge\_results.py}}
\label{subsec:cluster-merge}

Sobald alle Teilergebnisse auf dem Manager liegen, führt \texttt{merge\_results.py}
sie zu einer einzigen Impulsantwort zusammen. Jeder Knoten liefert dasselbe
Ausgabeschema, nämlich Metadaten sowie zwei positionsgleiche Listen aus
Verzögerungen und Drücken. Die Aggregation besteht deshalb im Kern aus einer
schlichten Verkettung dieser Listen:

\begin{lstlisting}[style=clusterlst, language=Python, caption={Zusammenführen der Knoten-Ergebnisse, gekürzt (\texttt{merge\_results.py}).}]
all_delays, all_pressures = [], []
for f in sorted(glob.glob("/home/manager/ir_node*.json")):
    with open(f) as file:
        data = json.load(file)
        all_delays    += data["hits"]["delays_seconds"]
        all_pressures += data["hits"]["pressures"]

merged = {
    "metadata": {"sample_rate": 44100, "rays_cast": 1000000 * 8,
                 "rays_received": len(all_delays), "room_name": "MVP_Test_Room_1"},
    "hits": {"delays_seconds": all_delays, "pressures": all_pressures},
}
json.dump(merged, open("/home/manager/ir_merged.json", "w"), indent=2)
\end{lstlisting}

Dass die Verkettung hier die richtige Operation ist, folgt direkt aus dem
Replikat-Modell. Weil die einzelnen Läufe unabhängige Stichproben desselben Raums
sind, ist die Vereinigung aller Treffer statistisch identisch mit dem Ergebnis
eines einzigen Laufs mit acht Millionen Strahlen. Es findet keine Mittelung und
keine Gewichtung statt, sondern nur ein Aneinanderhängen. Das Feld
\texttt{rays\_cast} ist im Skript fest auf $8 \cdot 10^{6}$ gesetzt; solange der in
Abschnitt~\ref{subsec:cluster-grenzen} beschriebene Kopierfehler bei Knoten~6
besteht, fließen tatsächlich nur sieben Knoten in die Aggregation ein, und dieser
Metadatenwert überschätzt die reale Strahlzahl. \texttt{rays\_received} zählt
dagegen die tatsächlich empfangenen Treffer. Die
zusammengeführte Datei trägt exakt dasselbe Schema wie die Ausgabe eines einzelnen
Knotens und ist damit ohne weitere Anpassung für die DSP-Pipeline verwendbar.

\subsection{Datenfluss im Überblick}
\label{subsec:cluster-datenfluss}

Der gesamte verteilte Ablauf lässt sich in fünf Stufen zusammenfassen:

\begin{enumerate}
    \item Der Manager startet per SSH auf jedem der acht Knoten
    \texttt{rust\_service} mit identischer Konfiguration, alle parallel.
    \item Jeder Knoten simuliert unabhängig eine Million Strahlen und schreibt
    lokal seine \texttt{ir\_output.json}.
    \item \texttt{wait} synchronisiert: Die Einsammel-Phase beginnt erst, wenn der
    letzte Knoten fertig ist.
    \item \texttt{scp} überträgt jede \texttt{ir\_output.json} als
    \texttt{ir\_node$k$.json} zum Manager.
    \item \texttt{merge\_results.py} verkettet alle Trefferlisten zu einer
    \texttt{ir\_merged.json}, die anschließend die DSP-Pipeline verarbeitet.
\end{enumerate}

\subsection{Leistungsbetrachtung}
\label{subsec:cluster-leistung}

\subsubsection{Erwarteter Speedup}
\label{subsubsec:speedup}

Bei der Leistung sind zwei Sichtweisen zu trennen. In der umgesetzten
Replikat-Betriebsart rechnet jeder Knoten dieselbe Strahlzahl, sodass acht Knoten
in nahezu derselben Wanduhrzeit die achtfache Stichprobe liefern. Der Gewinn ist
hier keine kürzere Laufzeit, sondern die höhere Genauigkeit aus
Abbildung~\ref{fig:mc-convergence}; man spricht vom schwachen Skalieren im Sinne
von Gustafson, bei dem mit der Rechnerzahl die bearbeitete Problemgröße wächst,
nicht die Geschwindigkeit einer festen Aufgabe. Will man dagegen eine feste
Gesamtstrahlzahl in kürzerer Zeit erreichen, so verteilt man diese auf die Knoten,
sodass jeder nur $8\cdot10^{6}/k$ Strahlen rechnet. Erst in dieser
Partitionierungs-Betriebsart greift das Gesetz von Amdahl. Bezeichnet $p$ den parallelen
Anteil der Arbeit, hier die reine Strahlverfolgung, und $1-p$ den seriellen Rest
aus Prozessstart, Ergebnistransfer und Aggregation, so beträgt der theoretische
Speedup
\[
    S(k) \;=\; \frac{1}{(1-p) + \dfrac{p}{k}} .
\]
Da die Strahlverfolgung die Laufzeit klar dominiert, liegt $p$ nahe bei eins, und
$S(k)$ steigt bis $k=8$ nahezu linear. In der Praxis bleibt der reale Speedup
darunter, und zwar aus drei Gründen: dem bereits erwähnten Straggler-Effekt (der
Lauf ist so schnell wie der langsamste Knoten), dem Overhead für den
SSH-Verbindungsaufbau sowie der Übertragung und Zusammenführung der
Ergebnisdateien.

\subsubsection{Kommunikations- und I/O-Overhead}
\label{subsubsec:overhead}

Der Overhead dieses Ansatzes ist klein, aber nicht null. Pro Knoten fallen ein
SSH-Verbindungsaufbau und eine \texttt{scp}-Übertragung an. Übertragen wird jeweils
eine JSON-Datei, deren Größe mit der Zahl der empfangenen Treffer wächst; bei
einer Million geworfener Strahlen liegt sie im Bereich einiger Megabyte. Da diese
acht Übertragungen erst \emph{nach} der Rechenphase stattfinden und nicht mit ihr
überlappen, addieren sie sich voll zur Gesamtlaufzeit. Bei deutlich mehr Knoten
oder größeren Ergebnisdateien würde dieser sequenzielle Einsammel-Schritt zum
Flaschenhals; eine Überlappung von Rechnung und Transfer oder ein kompakteres
Ausgabeformat als das gut lesbare, aber platzintensive JSON wären dann die
naheliegenden Gegenmaßnahmen.

\subsubsection{Messung am Cluster}
\label{subsubsec:messung}

Die konkreten Laufzeiten hängen von der Laborhardware ab und sind am Cluster zu
messen. Um einen aussagekräftigen Speedup zu erhalten, wird in der
Partitionierungs-Betriebsart gemessen: In \texttt{input\_config.json} wird
\texttt{rays\_to\_cast} auf $8\cdot10^{6}/k$ gesetzt, sodass die Gesamtarbeit über
alle $k$ Knoten konstant bleibt. Dann wird \texttt{run\_cluster.sh} für
verschiedene Knotenzahlen mit der Shell-eigenen Zeitmessung ausgeführt:

\begin{lstlisting}[style=clusterlst, language=bash, caption={Messvorschrift für die Laufzeit am Cluster.}]
time ./run_cluster.sh        # Wanduhrzeit fuer den kompletten verteilten Lauf
\end{lstlisting}

Tabelle~\ref{tab:speedup} nimmt die Messwerte auf. Als Referenz dient der Lauf mit
einem einzelnen Knoten; der Speedup ergibt sich als Verhältnis der Wanduhrzeiten.

\begin{table}[htbp]
    \centering
    \caption{Wanduhrzeit und Speedup des verteilten Laufs (am Cluster zu messen).}
    \label{tab:speedup}
    \begin{tabular}{r r r}
        \toprule
        Knoten $k$ & Wanduhrzeit [s] & Speedup $S$ \\
        \midrule
        1 & \dots & $1{,}00$ \\
        2 & \dots & \dots \\
        4 & \dots & \dots \\
        8 & \dots & \dots \\
        \bottomrule
    \end{tabular}
\end{table}

\subsection{Fehlertoleranz und Sicherheit}
\label{subsec:cluster-fehler}

Zwei Querschnittsthemen verdienen eigene Betrachtung. Beim \textbf{Sicherheits}aspekt
ist die Schlüssel-basierte SSH-Authentifizierung zentral: Der Manager benötigt
passwortlosen Zugriff auf jeden Knoten, was über hinterlegte öffentliche Schlüssel
gelöst ist. Der private Schlüssel des Managers ist damit das entscheidende
Geheimnis; wer ihn besitzt, kann auf allen Knoten Code ausführen. In einem
abgeschlossenen Labornetz ist das vertretbar, für einen offeneren Betrieb wäre eine
engere Rechtevergabe angebracht.

Beim Thema \textbf{Fehlertoleranz} zeigt der aktuelle Stand seine deutlichste
Schwäche, die im folgenden Abschnitt zusammen mit den übrigen Grenzen diskutiert
wird. Der Kern des Problems: Ein ausgefallener Knoten wird nicht erkannt, sondern
schlicht übergangen, sodass die Stichprobengröße unbemerkt sinkt.

\subsection{Grenzen des aktuellen Standes und Ausblick}
\label{subsec:cluster-grenzen}

Die Verteilungsschicht ist als funktionsfähiger Prototyp umgesetzt. Für einen
robusten Betrieb bestehen mehrere klar benannte Schwachstellen, die zugleich die
nächsten Entwicklungsschritte vorgeben.

\begin{itemize}
    \item \textbf{Fest verdrahtete Konfiguration.} IP-Adressen, Knotennamen, die
    Knotenzahl und die Ablagepfade stehen direkt im Skript. Ein Umzug auf andere
    Hardware oder eine andere Knotenzahl erfordert Änderungen am Quelltext. Eine
    Auslagerung dieser Werte in eine Konfigurationsdatei wäre der naheliegende
    erste Schritt.

    \item \textbf{Kopierfehler bei Knoten~6.} Im Einsammel-Schritt wird das
    Ergebnis von Knoten~6 versehentlich nach \texttt{ir\_node5.json} kopiert und
    überschreibt damit das Ergebnis von Knoten~5. Effektiv fließen so nur sieben
    statt acht unabhängige Stichproben in die Aggregation ein. Der Fehler ist ein
    Tippfehler im Zielnamen und leicht zu beheben, führt aber vor Augen, wie
    fehleranfällig die von Hand ausgeschriebene Knotenliste ist -- ein weiteres
    Argument für die genannte Auslagerung in eine Schleife über eine
    Konfigurationsliste.

    \item \textbf{Keine Fehlererkennung.} Fällt ein Knoten aus oder liefert er kein
    Ergebnis, schlägt der zugehörige \texttt{scp}-Aufruf fehl, der Gesamtlauf setzt
    sich aber fort und aggregiert stillschweigend über die vorhandenen Dateien. Ein
    Knotenausfall bleibt damit unbemerkt und verfälscht die Stichprobengröße. Eine
    Prüfung der Rückgabecodes und eine Warnung bei fehlenden Ergebnissen wären hier
    angebracht.

    \item \textbf{Manuelle Kopplung an die DSP-Pipeline.} Die Aggregation erzeugt
    \texttt{ir\_merged.json}, die Python-DSP-Pipeline liest jedoch
    \texttt{ir\_output.json}. Der Übergang zwischen verteiltem Lauf und Faltung ist
    derzeit ein manueller Schritt. Eine saubere Integration würde den Cluster-Lauf
    direkt aus der Orchestrierung heraus anstoßen und die zusammengeführte Datei
    automatisch weiterreichen.

    \item \textbf{Statische Verteilung.} Der Replikat-Ansatz mit gleich großen
    Aufträgen ist einfach und robust, setzt aber homogene Hardware voraus und lastet
    unterschiedlich schnelle Knoten nicht unterschiedlich aus. Ein Modell mit
    dynamischer Arbeitsverteilung -- ein Manager vergibt Ray-Kontingente, sobald ein
    Knoten frei wird -- würde heterogene Hardware besser nutzen und die Knotenzahl
    flexibel machen.
\end{itemize}

Alle genannten Punkte betreffen Robustheit und Bedienbarkeit, nicht die
grundsätzliche Korrektheit des Verfahrens. Die verteilte Ausführung liefert eine
gültige, statistisch dichtere Impulsantwort, und das Manager-Worker-Muster über
SSH und \texttt{scp} erfüllt seinen Zweck im Laborbetrieb.

\subsection{Zusammenfassung}
\label{subsec:cluster-zusammenfassung}

Dieses Kapitel hat die verteilte Ausführung der Strahlverfolgung auf dem
Labor-Cluster beschrieben. Weil die Simulation embarrassingly parallel ist, genügt
ein Replikat-Ansatz: Alle acht Knoten rechnen denselben Auftrag mit unabhängigen
Zufallsströmen, was statistisch einer achtfach größeren Simulation entspricht und
den Schätzfehler der Impulsantwort um rund $\sqrt{8}$ senkt -- ein Verhalten, das
das numerische Experiment in Abbildung~\ref{fig:mc-convergence} bestätigt. Ein
Shell-Skript startet die Knoten per SSH parallel, synchronisiert über \texttt{wait}
und sammelt die Ergebnisse per \texttt{scp} ein; \texttt{merge\_results.py}
verkettet die Trefferlisten zu einer gemeinsamen Datei im gewohnten Schema. Der
Prototyp erfüllt seine Aufgabe, weist aber in Konfigurierbarkeit, Fehlererkennung
und der Anbindung an die DSP-Pipeline die diskutierten Schwächen auf, die den
Rahmen für die Weiterentwicklung abstecken.

%   \end{document}
% =====================================================================

\lstdefinestyle{clusterlst}{
    basicstyle=\ttfamily\small,
    showstringspaces=false,
    breaklines=true,
    numbers=left,
    numberstyle=\tiny,
    frame=single,
    tabsize=2
}

%======================================================================
\section{Verteilte Ausführung auf dem Rechencluster}
\label{sec:cluster}

\subsection{Einordnung und Aufgabenstellung}
\label{subsec:cluster-einordnung}

Die Strahlverfolgung ist der teuerste Schritt der gesamten Pipeline. Für eine
brauchbare Impulsantwort müssen Hunderttausende bis Millionen einzelner
Schallstrahlen durch die Raumgeometrie verfolgt werden, und jeder dieser Strahlen
verlangt mehrere Schnitt- und Reflexionsberechnungen. Der entscheidende Vorteil
dabei: Kein Strahl hängt vom Ergebnis eines anderen ab. Das Problem ist damit
\emph{embarrassingly parallel} und lässt sich ohne Kommunikation zwischen den
Teilrechnungen auf beliebig viele Rechner verteilen.

Dieses Kapitel dokumentiert die Verteilungsschicht des Systems. Sie startet die
Rust-Simulation gleichzeitig auf allen acht Knoten des Labor-Clusters, wartet auf
deren Abschluss, sammelt die Teilergebnisse ein und fügt sie zu einer einzigen,
statistisch dichteren Impulsantwort zusammen. Aus Sicht der übrigen Pipeline
ändert sich dabei nichts: Am Ende steht wieder eine Datei im gewohnten Schema,
die die Signalverarbeitung unverändert weiterverarbeiten kann.

Zur Abgrenzung von den Nachbarkapiteln gehören zwei Punkte. Erstens ist die
\emph{prozessinterne} Nebenläufigkeit der Rust-Engine, also die Verteilung auf
die Kerne eines einzelnen Rechners mittels \texttt{rayon}, an anderer Stelle
beschrieben und bleibt hier außen vor. Sie parallelisiert innerhalb eines
Prozesses; dieses Kapitel parallelisiert über Prozess- und Rechnergrenzen hinweg.
Zweitens geht es nicht um den einzelnen Subprozess-Aufruf aus der Python-Pipeline,
sondern um die Ebene darüber: die gleichzeitige Ausführung auf mehreren physischen
Knoten und die anschließende Aggregation ihrer Ergebnisse.

\subsection{Theoretischer Hintergrund}
\label{subsec:cluster-theorie}

\subsubsection{Monte-Carlo-Strahlverfolgung}
\label{subsubsec:mc}

Die Engine schätzt die Impulsantwort eines Raumes stochastisch. Sie verschickt
zufällig gerichtete Strahlen von der Schallquelle, verfolgt ihre Reflexionen und
registriert, wann und mit welchem Restdruck sie den Empfänger erreichen. Jeder
einzelne Strahl ist eine Zufallsstichprobe; erst die Summe vieler Strahlen ergibt
ein aussagekräftiges Bild der Raumakustik. Solche Monte-Carlo-Verfahren haben eine
charakteristische Eigenschaft: Ihr Schätzfehler sinkt mit der Wurzel der
Stichprobengröße. Verdoppelt man die Zahl der Strahlen, halbiert sich der Fehler
also nicht, sondern sinkt nur um den Faktor $\sqrt{2}$. Genau diese langsame
Konvergenz macht große Strahlzahlen nötig und motiviert die Verteilung auf mehrere
Rechner.

\subsubsection{Zwei Wege der Verteilung}
\label{subsubsec:strategien}

Ein Monte-Carlo-Verfahren lässt sich auf zwei grundsätzlich verschiedene Arten
verteilen. Tabelle~\ref{tab:strategien} stellt sie gegenüber.

\begin{table}[htbp]
    \centering
    \caption{Zwei Strategien zur Verteilung der Strahlverfolgung.}
    \label{tab:strategien}
    \begin{tabular}{p{3.2cm} p{5.2cm} p{4.2cm}}
        \toprule
        & \textbf{Work-Partitioning} & \textbf{Replikat-Parallelismus} \\
        \midrule
        Prinzip & Gesamtzahl der Strahlen wird in disjunkte Anteile zerlegt & Jeder Knoten rechnet denselben Auftrag mit eigenem Zufallsstrom \\
        Koordination & Zuteilung der Anteile nötig & keine, jeder Knoten ist autark \\
        Zufallszahlen & ein globaler Strom, aufgeteilt & unabhängige Ströme pro Knoten \\
        Robustheit & Ausfall trifft einen festen Anteil & Ausfall kostet nur Stichprobengröße \\
        \bottomrule
    \end{tabular}
\end{table}

Das umgesetzte System verwendet den zweiten Weg. Alle acht Knoten starten die
Simulation mit derselben Konfigurationsdatei \texttt{input\_config.json} und
werfen dennoch jeweils andere Strahlen, weil jeder Prozess seinen eigenen
Zufallsstrom besitzt. Der Ansatz ist bewusst schlicht gewählt: Er braucht keine
zentrale Zuteilung und keinen Austausch zwischen den Knoten, sondern nur, dass am
Ende alle Ergebnisse eingesammelt werden.

\subsubsection{Unabhängigkeit der Zufallsströme}
\label{subsubsec:rng}

Die Korrektheit des Replikat-Ansatzes steht und fällt mit einer Bedingung: Die
Zufallszahlen auf den Knoten müssen unabhängig sein. Würde jeder Knoten dieselbe
Folge zufälliger Richtungen erzeugen, so lieferten alle acht exakt dasselbe
Ergebnis, und das Zusammenführen brächte keinen Erkenntnisgewinn, sondern nur
achtfach dieselben Strahlen. Die Rust-Engine schließt das aus, indem sie ihren
Zufallsstrom pro Thread bzw. Prozess frisch aus der Systementropie zieht:

\begin{lstlisting}[style=clusterlst, language=Rust, caption={Initialisierung des Zufallsstroms pro Prozess (\texttt{simulation.rs}).}]
let mut rng = rand::rng();                       // Omnidirectional, random rays
let random_angle: f32 = rng.random_range(0.0..(2.0 * PI));
let direction = Vec2::new(random_angle.cos(), random_angle.sin());
\end{lstlisting}

Weil jeder der acht Prozesse auf seinem Knoten einen anderen Startzustand erhält,
erzeugt jeder eine andere Stichprobe. Acht Knoten mit je einer Million Strahlen
liefern damit acht Millionen \emph{unabhängige} Stichproben derselben Raumakustik.

\subsubsection{Statistische Konvergenz}
\label{subsubsec:konvergenz}

Statistisch verhält sich der verteilte Lauf wie eine einzige, achtfach größere
Simulation. Für einen Monte-Carlo-Schätzer aus $N$ unabhängigen Teilstichproben
gleicher Größe sinkt die Standardabweichung des Mittelwerts mit
\[
    \sigma_N \;=\; \frac{\sigma_1}{\sqrt{N}} .
\]
Acht Knoten reduzieren die Streuung der geschätzten Impulsantwort demnach um den
Faktor $\sqrt{8} \approx 2{,}83$ gegenüber einem einzelnen Knoten.
Abbildung~\ref{fig:mc-convergence} bestätigt dieses Verhalten in einem numerischen
Experiment: Aufgetragen ist der mittlere quadratische Fehler der geschätzten
Impulsantwort über der Knotenzahl. Die Messpunkte folgen sauber der
$1/\sqrt{N}$-Referenzkurve; von einem auf acht Knoten sinkt der Fehler um den
Faktor $2{,}78$, praktisch exakt der theoretisch erwartete Wert.

\begin{figure}[htbp]
    \centering
    \includegraphics[width=0.8\textwidth]{figures/mc_convergence}
    \caption{Mittlerer quadratischer Fehler der geschätzten Impulsantwort über der
    Knotenzahl $N$ (je $5\cdot10^4$ Strahlen, doppelt logarithmisch). Die
    Messpunkte folgen der $1/\sqrt{N}$-Referenz; acht Knoten senken den Fehler um
    den Faktor $2{,}78 \approx \sqrt{8}$.}
    \label{fig:mc-convergence}
\end{figure}

Bemerkenswert ist die Kehrseite dieser Kurve: Um den Fehler zu halbieren, braucht
es die vierfache Knotenzahl. Der Nutzen zusätzlicher Knoten nimmt also ab, was bei
der Wahl der Clustergröße zu berücksichtigen ist. Für die hier angestrebte
Genauigkeit sind acht Knoten ein sinnvoller Kompromiss zwischen Aufwand und
Qualität.

\subsection{Architektur des Clusters}
\label{subsec:cluster-arch}

Das System folgt einem \emph{Manager-Worker}-Muster
(Abbildung~\ref{fig:cluster-arch}). Ein Manager-Knoten koordiniert den Ablauf,
acht Worker-Knoten führen die eigentliche Simulation aus. Der Manager selbst
rechnet nicht mit, sondern startet die Worker, wartet auf sie und führt am Ende
die Ergebnisse zusammen.

\begin{figure}[htbp]
    \centering
    \includegraphics[width=0.82\textwidth]{figures/cluster_arch}
    \caption{Manager-Worker-Architektur. Der Manager startet die acht Knoten
    parallel per \texttt{ssh}, holt nach Abschluss jede Ergebnisdatei per
    \texttt{scp} zurück und führt sie mit \texttt{merge\_results.py} zusammen.}
    \label{fig:cluster-arch}
\end{figure}

Die Kommunikation läuft vollständig über Standard-Werkzeuge des Betriebssystems.
\texttt{ssh} startet die Prozesse aus der Ferne, \texttt{scp} überträgt die
Ergebnisdateien zurück. Ein eigenes Netzwerkprotokoll oder ein
Message-Passing-Framework kommt nicht zum Einsatz. Für einen Stapelbetrieb, bei
dem jeder Knoten genau einmal startet, eine Weile rechnet und ein einziges
Ergebnis zurückliefert, wäre solcher Aufwand nicht gerechtfertigt: Es gibt keine
laufende Zustandssynchronisation und keinen Nachrichtenaustausch während der
Rechnung, den ein Framework verwalten müsste.

Damit das funktioniert, müssen drei Voraussetzungen erfüllt sein. Auf jedem Knoten
liegt dieselbe vorkompilierte Binärdatei \texttt{rust\_service} samt der
Konfiguration \texttt{input\_config.json}. Der Manager kann sich per SSH ohne
interaktive Passworteingabe anmelden, was eine Schlüssel-basierte
Authentifizierung voraussetzt. Und die Knoten sind über feste Adressen im
Labornetz erreichbar.

\subsection{Ablaufsteuerung: \texttt{run\_cluster.sh}}
\label{subsec:cluster-run}

Das Skript \texttt{run\_cluster.sh} steuert den verteilten Lauf in drei Phasen:
paralleles Starten, Warten und Einsammeln.

\begin{lstlisting}[style=clusterlst, language=bash, caption={Verteilter Start und Einsammeln der Ergebnisse (gekürzt).}]
# Phase 1: alle Knoten gleichzeitig starten (Hintergrundprozesse)
ssh node1@139.6.65.29 "./rust_service --config-file input_config.json" &
ssh node2@139.6.65.18 "./rust_service --config-file input_config.json" &
# ... node3 bis node8 analog ...

wait                       # Phase 2: auf Abschluss ALLER Knoten warten

# Phase 3: Ergebnis jedes Knotens zum Manager kopieren
scp node1@139.6.65.29:~/ir_output.json ~/ir_node1.json
scp node2@139.6.65.18:~/ir_output.json ~/ir_node2.json
# ... node3 bis node8 analog ...
\end{lstlisting}

Zwei Eigenschaften der Shell nutzt der Ablauf gezielt. Das kaufmännische Und
(\texttt{\&}) am Zeilenende schickt jeden \texttt{ssh}-Aufruf in den Hintergrund.
Die acht Simulationen laufen dadurch nahezu gleichzeitig an, statt nacheinander
abgearbeitet zu werden. Das darauffolgende \texttt{wait} blockiert, bis
\emph{alle} Hintergrundprozesse beendet sind, und bildet so eine
Synchronisationsbarriere zwischen der Rechen- und der Einsammel-Phase. Erst wenn
der letzte Knoten fertig ist, beginnt \texttt{scp}, von jedem Knoten dessen
\texttt{ir\_output.json} zu holen und unter einem knotenspezifischen Namen auf dem
Manager abzulegen.

Diese Barriere hat eine wichtige Konsequenz für die Laufzeit: Der Gesamtlauf ist
so langsam wie der \emph{langsamste} Knoten. Solange die Hardware homogen ist,
fällt das kaum ins Gewicht; auf ungleichen Knoten würde ein einzelner langsamer
Rechner den ganzen Lauf ausbremsen, während die übrigen sieben bereits fertig
warten. Dieses Straggler-Problem ist der zentrale Grund, warum ein statischer
Replikat-Ansatz mit gleich großen Aufträgen homogene Hardware voraussetzt.

\subsection{Aggregation der Teilergebnisse: \texttt{merge\_results.py}}
\label{subsec:cluster-merge}

Sobald alle Teilergebnisse auf dem Manager liegen, führt \texttt{merge\_results.py}
sie zu einer einzigen Impulsantwort zusammen. Jeder Knoten liefert dasselbe
Ausgabeschema, nämlich Metadaten sowie zwei positionsgleiche Listen aus
Verzögerungen und Drücken. Die Aggregation besteht deshalb im Kern aus einer
schlichten Verkettung dieser Listen:

\begin{lstlisting}[style=clusterlst, language=Python, caption={Zusammenführen der Knoten-Ergebnisse, gekürzt (\texttt{merge\_results.py}).}]
all_delays, all_pressures = [], []
for f in sorted(glob.glob("/home/manager/ir_node*.json")):
    with open(f) as file:
        data = json.load(file)
        all_delays    += data["hits"]["delays_seconds"]
        all_pressures += data["hits"]["pressures"]

merged = {
    "metadata": {"sample_rate": 44100, "rays_cast": 1000000 * 8,
                 "rays_received": len(all_delays), "room_name": "MVP_Test_Room_1"},
    "hits": {"delays_seconds": all_delays, "pressures": all_pressures},
}
json.dump(merged, open("/home/manager/ir_merged.json", "w"), indent=2)
\end{lstlisting}

Dass die Verkettung hier die richtige Operation ist, folgt direkt aus dem
Replikat-Modell. Weil die einzelnen Läufe unabhängige Stichproben desselben Raums
sind, ist die Vereinigung aller Treffer statistisch identisch mit dem Ergebnis
eines einzigen Laufs mit acht Millionen Strahlen. Es findet keine Mittelung und
keine Gewichtung statt, sondern nur ein Aneinanderhängen. Das Feld
\texttt{rays\_cast} ist im Skript fest auf $8 \cdot 10^{6}$ gesetzt; solange der in
Abschnitt~\ref{subsec:cluster-grenzen} beschriebene Kopierfehler bei Knoten~6
besteht, fließen tatsächlich nur sieben Knoten in die Aggregation ein, und dieser
Metadatenwert überschätzt die reale Strahlzahl. \texttt{rays\_received} zählt
dagegen die tatsächlich empfangenen Treffer. Die
zusammengeführte Datei trägt exakt dasselbe Schema wie die Ausgabe eines einzelnen
Knotens und ist damit ohne weitere Anpassung für die DSP-Pipeline verwendbar.

\subsection{Datenfluss im Überblick}
\label{subsec:cluster-datenfluss}

Der gesamte verteilte Ablauf lässt sich in fünf Stufen zusammenfassen:

\begin{enumerate}
    \item Der Manager startet per SSH auf jedem der acht Knoten
    \texttt{rust\_service} mit identischer Konfiguration, alle parallel.
    \item Jeder Knoten simuliert unabhängig eine Million Strahlen und schreibt
    lokal seine \texttt{ir\_output.json}.
    \item \texttt{wait} synchronisiert: Die Einsammel-Phase beginnt erst, wenn der
    letzte Knoten fertig ist.
    \item \texttt{scp} überträgt jede \texttt{ir\_output.json} als
    \texttt{ir\_node$k$.json} zum Manager.
    \item \texttt{merge\_results.py} verkettet alle Trefferlisten zu einer
    \texttt{ir\_merged.json}, die anschließend die DSP-Pipeline verarbeitet.
\end{enumerate}

\subsection{Leistungsbetrachtung}
\label{subsec:cluster-leistung}

\subsubsection{Erwarteter Speedup}
\label{subsubsec:speedup}

Bei der Leistung sind zwei Sichtweisen zu trennen. In der umgesetzten
Replikat-Betriebsart rechnet jeder Knoten dieselbe Strahlzahl, sodass acht Knoten
in nahezu derselben Wanduhrzeit die achtfache Stichprobe liefern. Der Gewinn ist
hier keine kürzere Laufzeit, sondern die höhere Genauigkeit aus
Abbildung~\ref{fig:mc-convergence}; man spricht vom schwachen Skalieren im Sinne
von Gustafson, bei dem mit der Rechnerzahl die bearbeitete Problemgröße wächst,
nicht die Geschwindigkeit einer festen Aufgabe. Will man dagegen eine feste
Gesamtstrahlzahl in kürzerer Zeit erreichen, so verteilt man diese auf die Knoten,
sodass jeder nur $8\cdot10^{6}/k$ Strahlen rechnet. Erst in dieser
Partitionierungs-Betriebsart greift das Gesetz von Amdahl. Bezeichnet $p$ den parallelen
Anteil der Arbeit, hier die reine Strahlverfolgung, und $1-p$ den seriellen Rest
aus Prozessstart, Ergebnistransfer und Aggregation, so beträgt der theoretische
Speedup
\[
    S(k) \;=\; \frac{1}{(1-p) + \dfrac{p}{k}} .
\]
Da die Strahlverfolgung die Laufzeit klar dominiert, liegt $p$ nahe bei eins, und
$S(k)$ steigt bis $k=8$ nahezu linear. In der Praxis bleibt der reale Speedup
darunter, und zwar aus drei Gründen: dem bereits erwähnten Straggler-Effekt (der
Lauf ist so schnell wie der langsamste Knoten), dem Overhead für den
SSH-Verbindungsaufbau sowie der Übertragung und Zusammenführung der
Ergebnisdateien.

\subsubsection{Kommunikations- und I/O-Overhead}
\label{subsubsec:overhead}

Der Overhead dieses Ansatzes ist klein, aber nicht null. Pro Knoten fallen ein
SSH-Verbindungsaufbau und eine \texttt{scp}-Übertragung an. Übertragen wird jeweils
eine JSON-Datei, deren Größe mit der Zahl der empfangenen Treffer wächst; bei
einer Million geworfener Strahlen liegt sie im Bereich einiger Megabyte. Da diese
acht Übertragungen erst \emph{nach} der Rechenphase stattfinden und nicht mit ihr
überlappen, addieren sie sich voll zur Gesamtlaufzeit. Bei deutlich mehr Knoten
oder größeren Ergebnisdateien würde dieser sequenzielle Einsammel-Schritt zum
Flaschenhals; eine Überlappung von Rechnung und Transfer oder ein kompakteres
Ausgabeformat als das gut lesbare, aber platzintensive JSON wären dann die
naheliegenden Gegenmaßnahmen.

\subsubsection{Messung am Cluster}
\label{subsubsec:messung}

Die konkreten Laufzeiten hängen von der Laborhardware ab und sind am Cluster zu
messen. Um einen aussagekräftigen Speedup zu erhalten, wird in der
Partitionierungs-Betriebsart gemessen: In \texttt{input\_config.json} wird
\texttt{rays\_to\_cast} auf $8\cdot10^{6}/k$ gesetzt, sodass die Gesamtarbeit über
alle $k$ Knoten konstant bleibt. Dann wird \texttt{run\_cluster.sh} für
verschiedene Knotenzahlen mit der Shell-eigenen Zeitmessung ausgeführt:

\begin{lstlisting}[style=clusterlst, language=bash, caption={Messvorschrift für die Laufzeit am Cluster.}]
time ./run_cluster.sh        # Wanduhrzeit fuer den kompletten verteilten Lauf
\end{lstlisting}

Tabelle~\ref{tab:speedup} nimmt die Messwerte auf. Als Referenz dient der Lauf mit
einem einzelnen Knoten; der Speedup ergibt sich als Verhältnis der Wanduhrzeiten.

\begin{table}[htbp]
    \centering
    \caption{Wanduhrzeit und Speedup des verteilten Laufs (am Cluster zu messen).}
    \label{tab:speedup}
    \begin{tabular}{r r r}
        \toprule
        Knoten $k$ & Wanduhrzeit [s] & Speedup $S$ \\
        \midrule
        1 & \dots & $1{,}00$ \\
        2 & \dots & \dots \\
        4 & \dots & \dots \\
        8 & \dots & \dots \\
        \bottomrule
    \end{tabular}
\end{table}

\subsection{Fehlertoleranz und Sicherheit}
\label{subsec:cluster-fehler}

Zwei Querschnittsthemen verdienen eigene Betrachtung. Beim \textbf{Sicherheits}aspekt
ist die Schlüssel-basierte SSH-Authentifizierung zentral: Der Manager benötigt
passwortlosen Zugriff auf jeden Knoten, was über hinterlegte öffentliche Schlüssel
gelöst ist. Der private Schlüssel des Managers ist damit das entscheidende
Geheimnis; wer ihn besitzt, kann auf allen Knoten Code ausführen. In einem
abgeschlossenen Labornetz ist das vertretbar, für einen offeneren Betrieb wäre eine
engere Rechtevergabe angebracht.

Beim Thema \textbf{Fehlertoleranz} zeigt der aktuelle Stand seine deutlichste
Schwäche, die im folgenden Abschnitt zusammen mit den übrigen Grenzen diskutiert
wird. Der Kern des Problems: Ein ausgefallener Knoten wird nicht erkannt, sondern
schlicht übergangen, sodass die Stichprobengröße unbemerkt sinkt.

\subsection{Grenzen des aktuellen Standes und Ausblick}
\label{subsec:cluster-grenzen}

Die Verteilungsschicht ist als funktionsfähiger Prototyp umgesetzt. Für einen
robusten Betrieb bestehen mehrere klar benannte Schwachstellen, die zugleich die
nächsten Entwicklungsschritte vorgeben.

\begin{itemize}
    \item \textbf{Fest verdrahtete Konfiguration.} IP-Adressen, Knotennamen, die
    Knotenzahl und die Ablagepfade stehen direkt im Skript. Ein Umzug auf andere
    Hardware oder eine andere Knotenzahl erfordert Änderungen am Quelltext. Eine
    Auslagerung dieser Werte in eine Konfigurationsdatei wäre der naheliegende
    erste Schritt.

    \item \textbf{Kopierfehler bei Knoten~6.} Im Einsammel-Schritt wird das
    Ergebnis von Knoten~6 versehentlich nach \texttt{ir\_node5.json} kopiert und
    überschreibt damit das Ergebnis von Knoten~5. Effektiv fließen so nur sieben
    statt acht unabhängige Stichproben in die Aggregation ein. Der Fehler ist ein
    Tippfehler im Zielnamen und leicht zu beheben, führt aber vor Augen, wie
    fehleranfällig die von Hand ausgeschriebene Knotenliste ist -- ein weiteres
    Argument für die genannte Auslagerung in eine Schleife über eine
    Konfigurationsliste.

    \item \textbf{Keine Fehlererkennung.} Fällt ein Knoten aus oder liefert er kein
    Ergebnis, schlägt der zugehörige \texttt{scp}-Aufruf fehl, der Gesamtlauf setzt
    sich aber fort und aggregiert stillschweigend über die vorhandenen Dateien. Ein
    Knotenausfall bleibt damit unbemerkt und verfälscht die Stichprobengröße. Eine
    Prüfung der Rückgabecodes und eine Warnung bei fehlenden Ergebnissen wären hier
    angebracht.

    \item \textbf{Manuelle Kopplung an die DSP-Pipeline.} Die Aggregation erzeugt
    \texttt{ir\_merged.json}, die Python-DSP-Pipeline liest jedoch
    \texttt{ir\_output.json}. Der Übergang zwischen verteiltem Lauf und Faltung ist
    derzeit ein manueller Schritt. Eine saubere Integration würde den Cluster-Lauf
    direkt aus der Orchestrierung heraus anstoßen und die zusammengeführte Datei
    automatisch weiterreichen.

    \item \textbf{Statische Verteilung.} Der Replikat-Ansatz mit gleich großen
    Aufträgen ist einfach und robust, setzt aber homogene Hardware voraus und lastet
    unterschiedlich schnelle Knoten nicht unterschiedlich aus. Ein Modell mit
    dynamischer Arbeitsverteilung -- ein Manager vergibt Ray-Kontingente, sobald ein
    Knoten frei wird -- würde heterogene Hardware besser nutzen und die Knotenzahl
    flexibel machen.
\end{itemize}

Alle genannten Punkte betreffen Robustheit und Bedienbarkeit, nicht die
grundsätzliche Korrektheit des Verfahrens. Die verteilte Ausführung liefert eine
gültige, statistisch dichtere Impulsantwort, und das Manager-Worker-Muster über
SSH und \texttt{scp} erfüllt seinen Zweck im Laborbetrieb.

\subsection{Zusammenfassung}
\label{subsec:cluster-zusammenfassung}

Dieses Kapitel hat die verteilte Ausführung der Strahlverfolgung auf dem
Labor-Cluster beschrieben. Weil die Simulation embarrassingly parallel ist, genügt
ein Replikat-Ansatz: Alle acht Knoten rechnen denselben Auftrag mit unabhängigen
Zufallsströmen, was statistisch einer achtfach größeren Simulation entspricht und
den Schätzfehler der Impulsantwort um rund $\sqrt{8}$ senkt -- ein Verhalten, das
das numerische Experiment in Abbildung~\ref{fig:mc-convergence} bestätigt. Ein
Shell-Skript startet die Knoten per SSH parallel, synchronisiert über \texttt{wait}
und sammelt die Ergebnisse per \texttt{scp} ein; \texttt{merge\_results.py}
verkettet die Trefferlisten zu einer gemeinsamen Datei im gewohnten Schema. Der
Prototyp erfüllt seine Aufgabe, weist aber in Konfigurierbarkeit, Fehlererkennung
und der Anbindung an die DSP-Pipeline die diskutierten Schwächen auf, die den
Rahmen für die Weiterentwicklung abstecken.

%   \end{document}
% =====================================================================

\lstdefinestyle{clusterlst}{
    basicstyle=\ttfamily\small,
    showstringspaces=false,
    breaklines=true,
    numbers=left,
    numberstyle=\tiny,
    frame=single,
    tabsize=2
}

%======================================================================
\section{Verteilte Ausführung auf dem Rechencluster}
\label{sec:cluster}

\subsection{Einordnung und Aufgabenstellung}
\label{subsec:cluster-einordnung}

Die Strahlverfolgung ist der teuerste Schritt der gesamten Pipeline. Für eine
brauchbare Impulsantwort müssen Hunderttausende bis Millionen einzelner
Schallstrahlen durch die Raumgeometrie verfolgt werden, und jeder dieser Strahlen
verlangt mehrere Schnitt- und Reflexionsberechnungen. Der entscheidende Vorteil
dabei: Kein Strahl hängt vom Ergebnis eines anderen ab. Das Problem ist damit
\emph{embarrassingly parallel} und lässt sich ohne Kommunikation zwischen den
Teilrechnungen auf beliebig viele Rechner verteilen.

Dieses Kapitel dokumentiert die Verteilungsschicht des Systems. Sie startet die
Rust-Simulation gleichzeitig auf allen acht Knoten des Labor-Clusters, wartet auf
deren Abschluss, sammelt die Teilergebnisse ein und fügt sie zu einer einzigen,
statistisch dichteren Impulsantwort zusammen. Aus Sicht der übrigen Pipeline
ändert sich dabei nichts: Am Ende steht wieder eine Datei im gewohnten Schema,
die die Signalverarbeitung unverändert weiterverarbeiten kann.

Zur Abgrenzung von den Nachbarkapiteln gehören zwei Punkte. Erstens ist die
\emph{prozessinterne} Nebenläufigkeit der Rust-Engine, also die Verteilung auf
die Kerne eines einzelnen Rechners mittels \texttt{rayon}, an anderer Stelle
beschrieben und bleibt hier außen vor. Sie parallelisiert innerhalb eines
Prozesses; dieses Kapitel parallelisiert über Prozess- und Rechnergrenzen hinweg.
Zweitens geht es nicht um den einzelnen Subprozess-Aufruf aus der Python-Pipeline,
sondern um die Ebene darüber: die gleichzeitige Ausführung auf mehreren physischen
Knoten und die anschließende Aggregation ihrer Ergebnisse.

\subsection{Theoretischer Hintergrund}
\label{subsec:cluster-theorie}

\subsubsection{Monte-Carlo-Strahlverfolgung}
\label{subsubsec:mc}

Die Engine schätzt die Impulsantwort eines Raumes stochastisch. Sie verschickt
zufällig gerichtete Strahlen von der Schallquelle, verfolgt ihre Reflexionen und
registriert, wann und mit welchem Restdruck sie den Empfänger erreichen. Jeder
einzelne Strahl ist eine Zufallsstichprobe; erst die Summe vieler Strahlen ergibt
ein aussagekräftiges Bild der Raumakustik. Solche Monte-Carlo-Verfahren haben eine
charakteristische Eigenschaft: Ihr Schätzfehler sinkt mit der Wurzel der
Stichprobengröße. Verdoppelt man die Zahl der Strahlen, halbiert sich der Fehler
also nicht, sondern sinkt nur um den Faktor $\sqrt{2}$. Genau diese langsame
Konvergenz macht große Strahlzahlen nötig und motiviert die Verteilung auf mehrere
Rechner.

\subsubsection{Zwei Wege der Verteilung}
\label{subsubsec:strategien}

Ein Monte-Carlo-Verfahren lässt sich auf zwei grundsätzlich verschiedene Arten
verteilen. Tabelle~\ref{tab:strategien} stellt sie gegenüber.

\begin{table}[htbp]
    \centering
    \caption{Zwei Strategien zur Verteilung der Strahlverfolgung.}
    \label{tab:strategien}
    \begin{tabular}{p{3.2cm} p{5.2cm} p{4.2cm}}
        \toprule
        & \textbf{Work-Partitioning} & \textbf{Replikat-Parallelismus} \\
        \midrule
        Prinzip & Gesamtzahl der Strahlen wird in disjunkte Anteile zerlegt & Jeder Knoten rechnet denselben Auftrag mit eigenem Zufallsstrom \\
        Koordination & Zuteilung der Anteile nötig & keine, jeder Knoten ist autark \\
        Zufallszahlen & ein globaler Strom, aufgeteilt & unabhängige Ströme pro Knoten \\
        Robustheit & Ausfall trifft einen festen Anteil & Ausfall kostet nur Stichprobengröße \\
        \bottomrule
    \end{tabular}
\end{table}

Das umgesetzte System verwendet den zweiten Weg. Alle acht Knoten starten die
Simulation mit derselben Konfigurationsdatei \texttt{input\_config.json} und
werfen dennoch jeweils andere Strahlen, weil jeder Prozess seinen eigenen
Zufallsstrom besitzt. Der Ansatz ist bewusst schlicht gewählt: Er braucht keine
zentrale Zuteilung und keinen Austausch zwischen den Knoten, sondern nur, dass am
Ende alle Ergebnisse eingesammelt werden.

\subsubsection{Unabhängigkeit der Zufallsströme}
\label{subsubsec:rng}

Die Korrektheit des Replikat-Ansatzes steht und fällt mit einer Bedingung: Die
Zufallszahlen auf den Knoten müssen unabhängig sein. Würde jeder Knoten dieselbe
Folge zufälliger Richtungen erzeugen, so lieferten alle acht exakt dasselbe
Ergebnis, und das Zusammenführen brächte keinen Erkenntnisgewinn, sondern nur
achtfach dieselben Strahlen. Die Rust-Engine schließt das aus, indem sie ihren
Zufallsstrom pro Thread bzw. Prozess frisch aus der Systementropie zieht:

\begin{lstlisting}[style=clusterlst, language=Rust, caption={Initialisierung des Zufallsstroms pro Prozess (\texttt{simulation.rs}).}]
let mut rng = rand::rng();                       // Omnidirectional, random rays
let random_angle: f32 = rng.random_range(0.0..(2.0 * PI));
let direction = Vec2::new(random_angle.cos(), random_angle.sin());
\end{lstlisting}

Weil jeder der acht Prozesse auf seinem Knoten einen anderen Startzustand erhält,
erzeugt jeder eine andere Stichprobe. Acht Knoten mit je einer Million Strahlen
liefern damit acht Millionen \emph{unabhängige} Stichproben derselben Raumakustik.

\subsubsection{Statistische Konvergenz}
\label{subsubsec:konvergenz}

Statistisch verhält sich der verteilte Lauf wie eine einzige, achtfach größere
Simulation. Für einen Monte-Carlo-Schätzer aus $N$ unabhängigen Teilstichproben
gleicher Größe sinkt die Standardabweichung des Mittelwerts mit
\[
    \sigma_N \;=\; \frac{\sigma_1}{\sqrt{N}} .
\]
Acht Knoten reduzieren die Streuung der geschätzten Impulsantwort demnach um den
Faktor $\sqrt{8} \approx 2{,}83$ gegenüber einem einzelnen Knoten.
Abbildung~\ref{fig:mc-convergence} bestätigt dieses Verhalten in einem numerischen
Experiment: Aufgetragen ist der mittlere quadratische Fehler der geschätzten
Impulsantwort über der Knotenzahl. Die Messpunkte folgen sauber der
$1/\sqrt{N}$-Referenzkurve; von einem auf acht Knoten sinkt der Fehler um den
Faktor $2{,}78$, praktisch exakt der theoretisch erwartete Wert.

\begin{figure}[htbp]
    \centering
    \includegraphics[width=0.8\textwidth]{figures/mc_convergence}
    \caption{Mittlerer quadratischer Fehler der geschätzten Impulsantwort über der
    Knotenzahl $N$ (je $5\cdot10^4$ Strahlen, doppelt logarithmisch). Die
    Messpunkte folgen der $1/\sqrt{N}$-Referenz; acht Knoten senken den Fehler um
    den Faktor $2{,}78 \approx \sqrt{8}$.}
    \label{fig:mc-convergence}
\end{figure}

Bemerkenswert ist die Kehrseite dieser Kurve: Um den Fehler zu halbieren, braucht
es die vierfache Knotenzahl. Der Nutzen zusätzlicher Knoten nimmt also ab, was bei
der Wahl der Clustergröße zu berücksichtigen ist. Für die hier angestrebte
Genauigkeit sind acht Knoten ein sinnvoller Kompromiss zwischen Aufwand und
Qualität.

\subsection{Architektur des Clusters}
\label{subsec:cluster-arch}

Das System folgt einem \emph{Manager-Worker}-Muster
(Abbildung~\ref{fig:cluster-arch}). Ein Manager-Knoten koordiniert den Ablauf,
acht Worker-Knoten führen die eigentliche Simulation aus. Der Manager selbst
rechnet nicht mit, sondern startet die Worker, wartet auf sie und führt am Ende
die Ergebnisse zusammen.

\begin{figure}[htbp]
    \centering
    \includegraphics[width=0.82\textwidth]{figures/cluster_arch}
    \caption{Manager-Worker-Architektur. Der Manager startet die acht Knoten
    parallel per \texttt{ssh}, holt nach Abschluss jede Ergebnisdatei per
    \texttt{scp} zurück und führt sie mit \texttt{merge\_results.py} zusammen.}
    \label{fig:cluster-arch}
\end{figure}

Die Kommunikation läuft vollständig über Standard-Werkzeuge des Betriebssystems.
\texttt{ssh} startet die Prozesse aus der Ferne, \texttt{scp} überträgt die
Ergebnisdateien zurück. Ein eigenes Netzwerkprotokoll oder ein
Message-Passing-Framework kommt nicht zum Einsatz. Für einen Stapelbetrieb, bei
dem jeder Knoten genau einmal startet, eine Weile rechnet und ein einziges
Ergebnis zurückliefert, wäre solcher Aufwand nicht gerechtfertigt: Es gibt keine
laufende Zustandssynchronisation und keinen Nachrichtenaustausch während der
Rechnung, den ein Framework verwalten müsste.

Damit das funktioniert, müssen drei Voraussetzungen erfüllt sein. Auf jedem Knoten
liegt dieselbe vorkompilierte Binärdatei \texttt{rust\_service} samt der
Konfiguration \texttt{input\_config.json}. Der Manager kann sich per SSH ohne
interaktive Passworteingabe anmelden, was eine Schlüssel-basierte
Authentifizierung voraussetzt. Und die Knoten sind über feste Adressen im
Labornetz erreichbar.

\subsection{Ablaufsteuerung: \texttt{run\_cluster.sh}}
\label{subsec:cluster-run}

Das Skript \texttt{run\_cluster.sh} steuert den verteilten Lauf in drei Phasen:
paralleles Starten, Warten und Einsammeln.

\begin{lstlisting}[style=clusterlst, language=bash, caption={Verteilter Start und Einsammeln der Ergebnisse (gekürzt).}]
# Phase 1: alle Knoten gleichzeitig starten (Hintergrundprozesse)
ssh node1@139.6.65.29 "./rust_service --config-file input_config.json" &
ssh node2@139.6.65.18 "./rust_service --config-file input_config.json" &
# ... node3 bis node8 analog ...

wait                       # Phase 2: auf Abschluss ALLER Knoten warten

# Phase 3: Ergebnis jedes Knotens zum Manager kopieren
scp node1@139.6.65.29:~/ir_output.json ~/ir_node1.json
scp node2@139.6.65.18:~/ir_output.json ~/ir_node2.json
# ... node3 bis node8 analog ...
\end{lstlisting}

Zwei Eigenschaften der Shell nutzt der Ablauf gezielt. Das kaufmännische Und
(\texttt{\&}) am Zeilenende schickt jeden \texttt{ssh}-Aufruf in den Hintergrund.
Die acht Simulationen laufen dadurch nahezu gleichzeitig an, statt nacheinander
abgearbeitet zu werden. Das darauffolgende \texttt{wait} blockiert, bis
\emph{alle} Hintergrundprozesse beendet sind, und bildet so eine
Synchronisationsbarriere zwischen der Rechen- und der Einsammel-Phase. Erst wenn
der letzte Knoten fertig ist, beginnt \texttt{scp}, von jedem Knoten dessen
\texttt{ir\_output.json} zu holen und unter einem knotenspezifischen Namen auf dem
Manager abzulegen.

Diese Barriere hat eine wichtige Konsequenz für die Laufzeit: Der Gesamtlauf ist
so langsam wie der \emph{langsamste} Knoten. Solange die Hardware homogen ist,
fällt das kaum ins Gewicht; auf ungleichen Knoten würde ein einzelner langsamer
Rechner den ganzen Lauf ausbremsen, während die übrigen sieben bereits fertig
warten. Dieses Straggler-Problem ist der zentrale Grund, warum ein statischer
Replikat-Ansatz mit gleich großen Aufträgen homogene Hardware voraussetzt.

\subsection{Aggregation der Teilergebnisse: \texttt{merge\_results.py}}
\label{subsec:cluster-merge}

Sobald alle Teilergebnisse auf dem Manager liegen, führt \texttt{merge\_results.py}
sie zu einer einzigen Impulsantwort zusammen. Jeder Knoten liefert dasselbe
Ausgabeschema, nämlich Metadaten sowie zwei positionsgleiche Listen aus
Verzögerungen und Drücken. Die Aggregation besteht deshalb im Kern aus einer
schlichten Verkettung dieser Listen:

\begin{lstlisting}[style=clusterlst, language=Python, caption={Zusammenführen der Knoten-Ergebnisse, gekürzt (\texttt{merge\_results.py}).}]
all_delays, all_pressures = [], []
for f in sorted(glob.glob("/home/manager/ir_node*.json")):
    with open(f) as file:
        data = json.load(file)
        all_delays    += data["hits"]["delays_seconds"]
        all_pressures += data["hits"]["pressures"]

merged = {
    "metadata": {"sample_rate": 44100, "rays_cast": 1000000 * 8,
                 "rays_received": len(all_delays), "room_name": "MVP_Test_Room_1"},
    "hits": {"delays_seconds": all_delays, "pressures": all_pressures},
}
json.dump(merged, open("/home/manager/ir_merged.json", "w"), indent=2)
\end{lstlisting}

Dass die Verkettung hier die richtige Operation ist, folgt direkt aus dem
Replikat-Modell. Weil die einzelnen Läufe unabhängige Stichproben desselben Raums
sind, ist die Vereinigung aller Treffer statistisch identisch mit dem Ergebnis
eines einzigen Laufs mit acht Millionen Strahlen. Es findet keine Mittelung und
keine Gewichtung statt, sondern nur ein Aneinanderhängen. Das Feld
\texttt{rays\_cast} ist im Skript fest auf $8 \cdot 10^{6}$ gesetzt; solange der in
Abschnitt~\ref{subsec:cluster-grenzen} beschriebene Kopierfehler bei Knoten~6
besteht, fließen tatsächlich nur sieben Knoten in die Aggregation ein, und dieser
Metadatenwert überschätzt die reale Strahlzahl. \texttt{rays\_received} zählt
dagegen die tatsächlich empfangenen Treffer. Die
zusammengeführte Datei trägt exakt dasselbe Schema wie die Ausgabe eines einzelnen
Knotens und ist damit ohne weitere Anpassung für die DSP-Pipeline verwendbar.

\subsection{Datenfluss im Überblick}
\label{subsec:cluster-datenfluss}

Der gesamte verteilte Ablauf lässt sich in fünf Stufen zusammenfassen:

\begin{enumerate}
    \item Der Manager startet per SSH auf jedem der acht Knoten
    \texttt{rust\_service} mit identischer Konfiguration, alle parallel.
    \item Jeder Knoten simuliert unabhängig eine Million Strahlen und schreibt
    lokal seine \texttt{ir\_output.json}.
    \item \texttt{wait} synchronisiert: Die Einsammel-Phase beginnt erst, wenn der
    letzte Knoten fertig ist.
    \item \texttt{scp} überträgt jede \texttt{ir\_output.json} als
    \texttt{ir\_node$k$.json} zum Manager.
    \item \texttt{merge\_results.py} verkettet alle Trefferlisten zu einer
    \texttt{ir\_merged.json}, die anschließend die DSP-Pipeline verarbeitet.
\end{enumerate}

\subsection{Leistungsbetrachtung}
\label{subsec:cluster-leistung}

\subsubsection{Erwarteter Speedup}
\label{subsubsec:speedup}

Bei der Leistung sind zwei Sichtweisen zu trennen. In der umgesetzten
Replikat-Betriebsart rechnet jeder Knoten dieselbe Strahlzahl, sodass acht Knoten
in nahezu derselben Wanduhrzeit die achtfache Stichprobe liefern. Der Gewinn ist
hier keine kürzere Laufzeit, sondern die höhere Genauigkeit aus
Abbildung~\ref{fig:mc-convergence}; man spricht vom schwachen Skalieren im Sinne
von Gustafson, bei dem mit der Rechnerzahl die bearbeitete Problemgröße wächst,
nicht die Geschwindigkeit einer festen Aufgabe. Will man dagegen eine feste
Gesamtstrahlzahl in kürzerer Zeit erreichen, so verteilt man diese auf die Knoten,
sodass jeder nur $8\cdot10^{6}/k$ Strahlen rechnet. Erst in dieser
Partitionierungs-Betriebsart greift das Gesetz von Amdahl. Bezeichnet $p$ den parallelen
Anteil der Arbeit, hier die reine Strahlverfolgung, und $1-p$ den seriellen Rest
aus Prozessstart, Ergebnistransfer und Aggregation, so beträgt der theoretische
Speedup
\[
    S(k) \;=\; \frac{1}{(1-p) + \dfrac{p}{k}} .
\]
Da die Strahlverfolgung die Laufzeit klar dominiert, liegt $p$ nahe bei eins, und
$S(k)$ steigt bis $k=8$ nahezu linear. In der Praxis bleibt der reale Speedup
darunter, und zwar aus drei Gründen: dem bereits erwähnten Straggler-Effekt (der
Lauf ist so schnell wie der langsamste Knoten), dem Overhead für den
SSH-Verbindungsaufbau sowie der Übertragung und Zusammenführung der
Ergebnisdateien.

\subsubsection{Kommunikations- und I/O-Overhead}
\label{subsubsec:overhead}

Der Overhead dieses Ansatzes ist klein, aber nicht null. Pro Knoten fallen ein
SSH-Verbindungsaufbau und eine \texttt{scp}-Übertragung an. Übertragen wird jeweils
eine JSON-Datei, deren Größe mit der Zahl der empfangenen Treffer wächst; bei
einer Million geworfener Strahlen liegt sie im Bereich einiger Megabyte. Da diese
acht Übertragungen erst \emph{nach} der Rechenphase stattfinden und nicht mit ihr
überlappen, addieren sie sich voll zur Gesamtlaufzeit. Bei deutlich mehr Knoten
oder größeren Ergebnisdateien würde dieser sequenzielle Einsammel-Schritt zum
Flaschenhals; eine Überlappung von Rechnung und Transfer oder ein kompakteres
Ausgabeformat als das gut lesbare, aber platzintensive JSON wären dann die
naheliegenden Gegenmaßnahmen.

\subsubsection{Messung am Cluster}
\label{subsubsec:messung}

Die konkreten Laufzeiten hängen von der Laborhardware ab und sind am Cluster zu
messen. Um einen aussagekräftigen Speedup zu erhalten, wird in der
Partitionierungs-Betriebsart gemessen: In \texttt{input\_config.json} wird
\texttt{rays\_to\_cast} auf $8\cdot10^{6}/k$ gesetzt, sodass die Gesamtarbeit über
alle $k$ Knoten konstant bleibt. Dann wird \texttt{run\_cluster.sh} für
verschiedene Knotenzahlen mit der Shell-eigenen Zeitmessung ausgeführt:

\begin{lstlisting}[style=clusterlst, language=bash, caption={Messvorschrift für die Laufzeit am Cluster.}]
time ./run_cluster.sh        # Wanduhrzeit fuer den kompletten verteilten Lauf
\end{lstlisting}

Tabelle~\ref{tab:speedup} nimmt die Messwerte auf. Als Referenz dient der Lauf mit
einem einzelnen Knoten; der Speedup ergibt sich als Verhältnis der Wanduhrzeiten.

\begin{table}[htbp]
    \centering
    \caption{Wanduhrzeit und Speedup des verteilten Laufs (am Cluster zu messen).}
    \label{tab:speedup}
    \begin{tabular}{r r r}
        \toprule
        Knoten $k$ & Wanduhrzeit [s] & Speedup $S$ \\
        \midrule
        1 & \dots & $1{,}00$ \\
        2 & \dots & \dots \\
        4 & \dots & \dots \\
        8 & \dots & \dots \\
        \bottomrule
    \end{tabular}
\end{table}

\subsection{Fehlertoleranz und Sicherheit}
\label{subsec:cluster-fehler}

Zwei Querschnittsthemen verdienen eigene Betrachtung. Beim \textbf{Sicherheits}aspekt
ist die Schlüssel-basierte SSH-Authentifizierung zentral: Der Manager benötigt
passwortlosen Zugriff auf jeden Knoten, was über hinterlegte öffentliche Schlüssel
gelöst ist. Der private Schlüssel des Managers ist damit das entscheidende
Geheimnis; wer ihn besitzt, kann auf allen Knoten Code ausführen. In einem
abgeschlossenen Labornetz ist das vertretbar, für einen offeneren Betrieb wäre eine
engere Rechtevergabe angebracht.

Beim Thema \textbf{Fehlertoleranz} zeigt der aktuelle Stand seine deutlichste
Schwäche, die im folgenden Abschnitt zusammen mit den übrigen Grenzen diskutiert
wird. Der Kern des Problems: Ein ausgefallener Knoten wird nicht erkannt, sondern
schlicht übergangen, sodass die Stichprobengröße unbemerkt sinkt.

\subsection{Grenzen des aktuellen Standes und Ausblick}
\label{subsec:cluster-grenzen}

Die Verteilungsschicht ist als funktionsfähiger Prototyp umgesetzt. Für einen
robusten Betrieb bestehen mehrere klar benannte Schwachstellen, die zugleich die
nächsten Entwicklungsschritte vorgeben.

\begin{itemize}
    \item \textbf{Fest verdrahtete Konfiguration.} IP-Adressen, Knotennamen, die
    Knotenzahl und die Ablagepfade stehen direkt im Skript. Ein Umzug auf andere
    Hardware oder eine andere Knotenzahl erfordert Änderungen am Quelltext. Eine
    Auslagerung dieser Werte in eine Konfigurationsdatei wäre der naheliegende
    erste Schritt.

    \item \textbf{Kopierfehler bei Knoten~6.} Im Einsammel-Schritt wird das
    Ergebnis von Knoten~6 versehentlich nach \texttt{ir\_node5.json} kopiert und
    überschreibt damit das Ergebnis von Knoten~5. Effektiv fließen so nur sieben
    statt acht unabhängige Stichproben in die Aggregation ein. Der Fehler ist ein
    Tippfehler im Zielnamen und leicht zu beheben, führt aber vor Augen, wie
    fehleranfällig die von Hand ausgeschriebene Knotenliste ist -- ein weiteres
    Argument für die genannte Auslagerung in eine Schleife über eine
    Konfigurationsliste.

    \item \textbf{Keine Fehlererkennung.} Fällt ein Knoten aus oder liefert er kein
    Ergebnis, schlägt der zugehörige \texttt{scp}-Aufruf fehl, der Gesamtlauf setzt
    sich aber fort und aggregiert stillschweigend über die vorhandenen Dateien. Ein
    Knotenausfall bleibt damit unbemerkt und verfälscht die Stichprobengröße. Eine
    Prüfung der Rückgabecodes und eine Warnung bei fehlenden Ergebnissen wären hier
    angebracht.

    \item \textbf{Manuelle Kopplung an die DSP-Pipeline.} Die Aggregation erzeugt
    \texttt{ir\_merged.json}, die Python-DSP-Pipeline liest jedoch
    \texttt{ir\_output.json}. Der Übergang zwischen verteiltem Lauf und Faltung ist
    derzeit ein manueller Schritt. Eine saubere Integration würde den Cluster-Lauf
    direkt aus der Orchestrierung heraus anstoßen und die zusammengeführte Datei
    automatisch weiterreichen.

    \item \textbf{Statische Verteilung.} Der Replikat-Ansatz mit gleich großen
    Aufträgen ist einfach und robust, setzt aber homogene Hardware voraus und lastet
    unterschiedlich schnelle Knoten nicht unterschiedlich aus. Ein Modell mit
    dynamischer Arbeitsverteilung -- ein Manager vergibt Ray-Kontingente, sobald ein
    Knoten frei wird -- würde heterogene Hardware besser nutzen und die Knotenzahl
    flexibel machen.
\end{itemize}

Alle genannten Punkte betreffen Robustheit und Bedienbarkeit, nicht die
grundsätzliche Korrektheit des Verfahrens. Die verteilte Ausführung liefert eine
gültige, statistisch dichtere Impulsantwort, und das Manager-Worker-Muster über
SSH und \texttt{scp} erfüllt seinen Zweck im Laborbetrieb.

\subsection{Zusammenfassung}
\label{subsec:cluster-zusammenfassung}

Dieses Kapitel hat die verteilte Ausführung der Strahlverfolgung auf dem
Labor-Cluster beschrieben. Weil die Simulation embarrassingly parallel ist, genügt
ein Replikat-Ansatz: Alle acht Knoten rechnen denselben Auftrag mit unabhängigen
Zufallsströmen, was statistisch einer achtfach größeren Simulation entspricht und
den Schätzfehler der Impulsantwort um rund $\sqrt{8}$ senkt -- ein Verhalten, das
das numerische Experiment in Abbildung~\ref{fig:mc-convergence} bestätigt. Ein
Shell-Skript startet die Knoten per SSH parallel, synchronisiert über \texttt{wait}
und sammelt die Ergebnisse per \texttt{scp} ein; \texttt{merge\_results.py}
verkettet die Trefferlisten zu einer gemeinsamen Datei im gewohnten Schema. Der
Prototyp erfüllt seine Aufgabe, weist aber in Konfigurierbarkeit, Fehlererkennung
und der Anbindung an die DSP-Pipeline die diskutierten Schwächen auf, die den
Rahmen für die Weiterentwicklung abstecken.
 im Master.
% Benoetigte Pakete: amsmath, listings, graphicx, booktabs.
% Fuer saubere Optik im Master: fontenc, babel[ngerman], microtype.
% Abbildungen: latexDocs/figures/cluster_arch.png, mc_convergence.png
%
% STANDALONE-VORSCHAU:
%   \documentclass[12pt]{article}
%   \usepackage[utf8]{inputenc}\usepackage[T1]{fontenc}
%   \usepackage[ngerman]{babel}\usepackage{microtype}
%   \usepackage{amsmath}\usepackage{listings}
%   \usepackage{graphicx}\usepackage{booktabs}
%   \begin{document}
%   % =====================================================================
% Kapitel: Verteilte Ausfuehrung auf dem Rechencluster
% Autor: Mier
%
% EINBINDEN: % =====================================================================
% Kapitel: Verteilte Ausfuehrung auf dem Rechencluster
% Autor: Mier
%
% EINBINDEN: % =====================================================================
% Kapitel: Verteilte Ausfuehrung auf dem Rechencluster
% Autor: Mier
%
% EINBINDEN: \input{Kapitel_Cluster_Mier} im Master.
% Benoetigte Pakete: amsmath, listings, graphicx, booktabs.
% Fuer saubere Optik im Master: fontenc, babel[ngerman], microtype.
% Abbildungen: latexDocs/figures/cluster_arch.png, mc_convergence.png
%
% STANDALONE-VORSCHAU:
%   \documentclass[12pt]{article}
%   \usepackage[utf8]{inputenc}\usepackage[T1]{fontenc}
%   \usepackage[ngerman]{babel}\usepackage{microtype}
%   \usepackage{amsmath}\usepackage{listings}
%   \usepackage{graphicx}\usepackage{booktabs}
%   \begin{document}
%   \input{Kapitel_Cluster_Mier}
%   \end{document}
% =====================================================================

\lstdefinestyle{clusterlst}{
    basicstyle=\ttfamily\small,
    showstringspaces=false,
    breaklines=true,
    numbers=left,
    numberstyle=\tiny,
    frame=single,
    tabsize=2
}

%======================================================================
\section{Verteilte Ausführung auf dem Rechencluster}
\label{sec:cluster}

\subsection{Einordnung und Aufgabenstellung}
\label{subsec:cluster-einordnung}

Die Strahlverfolgung ist der teuerste Schritt der gesamten Pipeline. Für eine
brauchbare Impulsantwort müssen Hunderttausende bis Millionen einzelner
Schallstrahlen durch die Raumgeometrie verfolgt werden, und jeder dieser Strahlen
verlangt mehrere Schnitt- und Reflexionsberechnungen. Der entscheidende Vorteil
dabei: Kein Strahl hängt vom Ergebnis eines anderen ab. Das Problem ist damit
\emph{embarrassingly parallel} und lässt sich ohne Kommunikation zwischen den
Teilrechnungen auf beliebig viele Rechner verteilen.

Dieses Kapitel dokumentiert die Verteilungsschicht des Systems. Sie startet die
Rust-Simulation gleichzeitig auf allen acht Knoten des Labor-Clusters, wartet auf
deren Abschluss, sammelt die Teilergebnisse ein und fügt sie zu einer einzigen,
statistisch dichteren Impulsantwort zusammen. Aus Sicht der übrigen Pipeline
ändert sich dabei nichts: Am Ende steht wieder eine Datei im gewohnten Schema,
die die Signalverarbeitung unverändert weiterverarbeiten kann.

Zur Abgrenzung von den Nachbarkapiteln gehören zwei Punkte. Erstens ist die
\emph{prozessinterne} Nebenläufigkeit der Rust-Engine, also die Verteilung auf
die Kerne eines einzelnen Rechners mittels \texttt{rayon}, an anderer Stelle
beschrieben und bleibt hier außen vor. Sie parallelisiert innerhalb eines
Prozesses; dieses Kapitel parallelisiert über Prozess- und Rechnergrenzen hinweg.
Zweitens geht es nicht um den einzelnen Subprozess-Aufruf aus der Python-Pipeline,
sondern um die Ebene darüber: die gleichzeitige Ausführung auf mehreren physischen
Knoten und die anschließende Aggregation ihrer Ergebnisse.

\subsection{Theoretischer Hintergrund}
\label{subsec:cluster-theorie}

\subsubsection{Monte-Carlo-Strahlverfolgung}
\label{subsubsec:mc}

Die Engine schätzt die Impulsantwort eines Raumes stochastisch. Sie verschickt
zufällig gerichtete Strahlen von der Schallquelle, verfolgt ihre Reflexionen und
registriert, wann und mit welchem Restdruck sie den Empfänger erreichen. Jeder
einzelne Strahl ist eine Zufallsstichprobe; erst die Summe vieler Strahlen ergibt
ein aussagekräftiges Bild der Raumakustik. Solche Monte-Carlo-Verfahren haben eine
charakteristische Eigenschaft: Ihr Schätzfehler sinkt mit der Wurzel der
Stichprobengröße. Verdoppelt man die Zahl der Strahlen, halbiert sich der Fehler
also nicht, sondern sinkt nur um den Faktor $\sqrt{2}$. Genau diese langsame
Konvergenz macht große Strahlzahlen nötig und motiviert die Verteilung auf mehrere
Rechner.

\subsubsection{Zwei Wege der Verteilung}
\label{subsubsec:strategien}

Ein Monte-Carlo-Verfahren lässt sich auf zwei grundsätzlich verschiedene Arten
verteilen. Tabelle~\ref{tab:strategien} stellt sie gegenüber.

\begin{table}[htbp]
    \centering
    \caption{Zwei Strategien zur Verteilung der Strahlverfolgung.}
    \label{tab:strategien}
    \begin{tabular}{p{3.2cm} p{5.2cm} p{4.2cm}}
        \toprule
        & \textbf{Work-Partitioning} & \textbf{Replikat-Parallelismus} \\
        \midrule
        Prinzip & Gesamtzahl der Strahlen wird in disjunkte Anteile zerlegt & Jeder Knoten rechnet denselben Auftrag mit eigenem Zufallsstrom \\
        Koordination & Zuteilung der Anteile nötig & keine, jeder Knoten ist autark \\
        Zufallszahlen & ein globaler Strom, aufgeteilt & unabhängige Ströme pro Knoten \\
        Robustheit & Ausfall trifft einen festen Anteil & Ausfall kostet nur Stichprobengröße \\
        \bottomrule
    \end{tabular}
\end{table}

Das umgesetzte System verwendet den zweiten Weg. Alle acht Knoten starten die
Simulation mit derselben Konfigurationsdatei \texttt{input\_config.json} und
werfen dennoch jeweils andere Strahlen, weil jeder Prozess seinen eigenen
Zufallsstrom besitzt. Der Ansatz ist bewusst schlicht gewählt: Er braucht keine
zentrale Zuteilung und keinen Austausch zwischen den Knoten, sondern nur, dass am
Ende alle Ergebnisse eingesammelt werden.

\subsubsection{Unabhängigkeit der Zufallsströme}
\label{subsubsec:rng}

Die Korrektheit des Replikat-Ansatzes steht und fällt mit einer Bedingung: Die
Zufallszahlen auf den Knoten müssen unabhängig sein. Würde jeder Knoten dieselbe
Folge zufälliger Richtungen erzeugen, so lieferten alle acht exakt dasselbe
Ergebnis, und das Zusammenführen brächte keinen Erkenntnisgewinn, sondern nur
achtfach dieselben Strahlen. Die Rust-Engine schließt das aus, indem sie ihren
Zufallsstrom pro Thread bzw. Prozess frisch aus der Systementropie zieht:

\begin{lstlisting}[style=clusterlst, language=Rust, caption={Initialisierung des Zufallsstroms pro Prozess (\texttt{simulation.rs}).}]
let mut rng = rand::rng();                       // Omnidirectional, random rays
let random_angle: f32 = rng.random_range(0.0..(2.0 * PI));
let direction = Vec2::new(random_angle.cos(), random_angle.sin());
\end{lstlisting}

Weil jeder der acht Prozesse auf seinem Knoten einen anderen Startzustand erhält,
erzeugt jeder eine andere Stichprobe. Acht Knoten mit je einer Million Strahlen
liefern damit acht Millionen \emph{unabhängige} Stichproben derselben Raumakustik.

\subsubsection{Statistische Konvergenz}
\label{subsubsec:konvergenz}

Statistisch verhält sich der verteilte Lauf wie eine einzige, achtfach größere
Simulation. Für einen Monte-Carlo-Schätzer aus $N$ unabhängigen Teilstichproben
gleicher Größe sinkt die Standardabweichung des Mittelwerts mit
\[
    \sigma_N \;=\; \frac{\sigma_1}{\sqrt{N}} .
\]
Acht Knoten reduzieren die Streuung der geschätzten Impulsantwort demnach um den
Faktor $\sqrt{8} \approx 2{,}83$ gegenüber einem einzelnen Knoten.
Abbildung~\ref{fig:mc-convergence} bestätigt dieses Verhalten in einem numerischen
Experiment: Aufgetragen ist der mittlere quadratische Fehler der geschätzten
Impulsantwort über der Knotenzahl. Die Messpunkte folgen sauber der
$1/\sqrt{N}$-Referenzkurve; von einem auf acht Knoten sinkt der Fehler um den
Faktor $2{,}78$, praktisch exakt der theoretisch erwartete Wert.

\begin{figure}[htbp]
    \centering
    \includegraphics[width=0.8\textwidth]{figures/mc_convergence}
    \caption{Mittlerer quadratischer Fehler der geschätzten Impulsantwort über der
    Knotenzahl $N$ (je $5\cdot10^4$ Strahlen, doppelt logarithmisch). Die
    Messpunkte folgen der $1/\sqrt{N}$-Referenz; acht Knoten senken den Fehler um
    den Faktor $2{,}78 \approx \sqrt{8}$.}
    \label{fig:mc-convergence}
\end{figure}

Bemerkenswert ist die Kehrseite dieser Kurve: Um den Fehler zu halbieren, braucht
es die vierfache Knotenzahl. Der Nutzen zusätzlicher Knoten nimmt also ab, was bei
der Wahl der Clustergröße zu berücksichtigen ist. Für die hier angestrebte
Genauigkeit sind acht Knoten ein sinnvoller Kompromiss zwischen Aufwand und
Qualität.

\subsection{Architektur des Clusters}
\label{subsec:cluster-arch}

Das System folgt einem \emph{Manager-Worker}-Muster
(Abbildung~\ref{fig:cluster-arch}). Ein Manager-Knoten koordiniert den Ablauf,
acht Worker-Knoten führen die eigentliche Simulation aus. Der Manager selbst
rechnet nicht mit, sondern startet die Worker, wartet auf sie und führt am Ende
die Ergebnisse zusammen.

\begin{figure}[htbp]
    \centering
    \includegraphics[width=0.82\textwidth]{figures/cluster_arch}
    \caption{Manager-Worker-Architektur. Der Manager startet die acht Knoten
    parallel per \texttt{ssh}, holt nach Abschluss jede Ergebnisdatei per
    \texttt{scp} zurück und führt sie mit \texttt{merge\_results.py} zusammen.}
    \label{fig:cluster-arch}
\end{figure}

Die Kommunikation läuft vollständig über Standard-Werkzeuge des Betriebssystems.
\texttt{ssh} startet die Prozesse aus der Ferne, \texttt{scp} überträgt die
Ergebnisdateien zurück. Ein eigenes Netzwerkprotokoll oder ein
Message-Passing-Framework kommt nicht zum Einsatz. Für einen Stapelbetrieb, bei
dem jeder Knoten genau einmal startet, eine Weile rechnet und ein einziges
Ergebnis zurückliefert, wäre solcher Aufwand nicht gerechtfertigt: Es gibt keine
laufende Zustandssynchronisation und keinen Nachrichtenaustausch während der
Rechnung, den ein Framework verwalten müsste.

Damit das funktioniert, müssen drei Voraussetzungen erfüllt sein. Auf jedem Knoten
liegt dieselbe vorkompilierte Binärdatei \texttt{rust\_service} samt der
Konfiguration \texttt{input\_config.json}. Der Manager kann sich per SSH ohne
interaktive Passworteingabe anmelden, was eine Schlüssel-basierte
Authentifizierung voraussetzt. Und die Knoten sind über feste Adressen im
Labornetz erreichbar.

\subsection{Ablaufsteuerung: \texttt{run\_cluster.sh}}
\label{subsec:cluster-run}

Das Skript \texttt{run\_cluster.sh} steuert den verteilten Lauf in drei Phasen:
paralleles Starten, Warten und Einsammeln.

\begin{lstlisting}[style=clusterlst, language=bash, caption={Verteilter Start und Einsammeln der Ergebnisse (gekürzt).}]
# Phase 1: alle Knoten gleichzeitig starten (Hintergrundprozesse)
ssh node1@139.6.65.29 "./rust_service --config-file input_config.json" &
ssh node2@139.6.65.18 "./rust_service --config-file input_config.json" &
# ... node3 bis node8 analog ...

wait                       # Phase 2: auf Abschluss ALLER Knoten warten

# Phase 3: Ergebnis jedes Knotens zum Manager kopieren
scp node1@139.6.65.29:~/ir_output.json ~/ir_node1.json
scp node2@139.6.65.18:~/ir_output.json ~/ir_node2.json
# ... node3 bis node8 analog ...
\end{lstlisting}

Zwei Eigenschaften der Shell nutzt der Ablauf gezielt. Das kaufmännische Und
(\texttt{\&}) am Zeilenende schickt jeden \texttt{ssh}-Aufruf in den Hintergrund.
Die acht Simulationen laufen dadurch nahezu gleichzeitig an, statt nacheinander
abgearbeitet zu werden. Das darauffolgende \texttt{wait} blockiert, bis
\emph{alle} Hintergrundprozesse beendet sind, und bildet so eine
Synchronisationsbarriere zwischen der Rechen- und der Einsammel-Phase. Erst wenn
der letzte Knoten fertig ist, beginnt \texttt{scp}, von jedem Knoten dessen
\texttt{ir\_output.json} zu holen und unter einem knotenspezifischen Namen auf dem
Manager abzulegen.

Diese Barriere hat eine wichtige Konsequenz für die Laufzeit: Der Gesamtlauf ist
so langsam wie der \emph{langsamste} Knoten. Solange die Hardware homogen ist,
fällt das kaum ins Gewicht; auf ungleichen Knoten würde ein einzelner langsamer
Rechner den ganzen Lauf ausbremsen, während die übrigen sieben bereits fertig
warten. Dieses Straggler-Problem ist der zentrale Grund, warum ein statischer
Replikat-Ansatz mit gleich großen Aufträgen homogene Hardware voraussetzt.

\subsection{Aggregation der Teilergebnisse: \texttt{merge\_results.py}}
\label{subsec:cluster-merge}

Sobald alle Teilergebnisse auf dem Manager liegen, führt \texttt{merge\_results.py}
sie zu einer einzigen Impulsantwort zusammen. Jeder Knoten liefert dasselbe
Ausgabeschema, nämlich Metadaten sowie zwei positionsgleiche Listen aus
Verzögerungen und Drücken. Die Aggregation besteht deshalb im Kern aus einer
schlichten Verkettung dieser Listen:

\begin{lstlisting}[style=clusterlst, language=Python, caption={Zusammenführen der Knoten-Ergebnisse, gekürzt (\texttt{merge\_results.py}).}]
all_delays, all_pressures = [], []
for f in sorted(glob.glob("/home/manager/ir_node*.json")):
    with open(f) as file:
        data = json.load(file)
        all_delays    += data["hits"]["delays_seconds"]
        all_pressures += data["hits"]["pressures"]

merged = {
    "metadata": {"sample_rate": 44100, "rays_cast": 1000000 * 8,
                 "rays_received": len(all_delays), "room_name": "MVP_Test_Room_1"},
    "hits": {"delays_seconds": all_delays, "pressures": all_pressures},
}
json.dump(merged, open("/home/manager/ir_merged.json", "w"), indent=2)
\end{lstlisting}

Dass die Verkettung hier die richtige Operation ist, folgt direkt aus dem
Replikat-Modell. Weil die einzelnen Läufe unabhängige Stichproben desselben Raums
sind, ist die Vereinigung aller Treffer statistisch identisch mit dem Ergebnis
eines einzigen Laufs mit acht Millionen Strahlen. Es findet keine Mittelung und
keine Gewichtung statt, sondern nur ein Aneinanderhängen. Das Feld
\texttt{rays\_cast} ist im Skript fest auf $8 \cdot 10^{6}$ gesetzt; solange der in
Abschnitt~\ref{subsec:cluster-grenzen} beschriebene Kopierfehler bei Knoten~6
besteht, fließen tatsächlich nur sieben Knoten in die Aggregation ein, und dieser
Metadatenwert überschätzt die reale Strahlzahl. \texttt{rays\_received} zählt
dagegen die tatsächlich empfangenen Treffer. Die
zusammengeführte Datei trägt exakt dasselbe Schema wie die Ausgabe eines einzelnen
Knotens und ist damit ohne weitere Anpassung für die DSP-Pipeline verwendbar.

\subsection{Datenfluss im Überblick}
\label{subsec:cluster-datenfluss}

Der gesamte verteilte Ablauf lässt sich in fünf Stufen zusammenfassen:

\begin{enumerate}
    \item Der Manager startet per SSH auf jedem der acht Knoten
    \texttt{rust\_service} mit identischer Konfiguration, alle parallel.
    \item Jeder Knoten simuliert unabhängig eine Million Strahlen und schreibt
    lokal seine \texttt{ir\_output.json}.
    \item \texttt{wait} synchronisiert: Die Einsammel-Phase beginnt erst, wenn der
    letzte Knoten fertig ist.
    \item \texttt{scp} überträgt jede \texttt{ir\_output.json} als
    \texttt{ir\_node$k$.json} zum Manager.
    \item \texttt{merge\_results.py} verkettet alle Trefferlisten zu einer
    \texttt{ir\_merged.json}, die anschließend die DSP-Pipeline verarbeitet.
\end{enumerate}

\subsection{Leistungsbetrachtung}
\label{subsec:cluster-leistung}

\subsubsection{Erwarteter Speedup}
\label{subsubsec:speedup}

Bei der Leistung sind zwei Sichtweisen zu trennen. In der umgesetzten
Replikat-Betriebsart rechnet jeder Knoten dieselbe Strahlzahl, sodass acht Knoten
in nahezu derselben Wanduhrzeit die achtfache Stichprobe liefern. Der Gewinn ist
hier keine kürzere Laufzeit, sondern die höhere Genauigkeit aus
Abbildung~\ref{fig:mc-convergence}; man spricht vom schwachen Skalieren im Sinne
von Gustafson, bei dem mit der Rechnerzahl die bearbeitete Problemgröße wächst,
nicht die Geschwindigkeit einer festen Aufgabe. Will man dagegen eine feste
Gesamtstrahlzahl in kürzerer Zeit erreichen, so verteilt man diese auf die Knoten,
sodass jeder nur $8\cdot10^{6}/k$ Strahlen rechnet. Erst in dieser
Partitionierungs-Betriebsart greift das Gesetz von Amdahl. Bezeichnet $p$ den parallelen
Anteil der Arbeit, hier die reine Strahlverfolgung, und $1-p$ den seriellen Rest
aus Prozessstart, Ergebnistransfer und Aggregation, so beträgt der theoretische
Speedup
\[
    S(k) \;=\; \frac{1}{(1-p) + \dfrac{p}{k}} .
\]
Da die Strahlverfolgung die Laufzeit klar dominiert, liegt $p$ nahe bei eins, und
$S(k)$ steigt bis $k=8$ nahezu linear. In der Praxis bleibt der reale Speedup
darunter, und zwar aus drei Gründen: dem bereits erwähnten Straggler-Effekt (der
Lauf ist so schnell wie der langsamste Knoten), dem Overhead für den
SSH-Verbindungsaufbau sowie der Übertragung und Zusammenführung der
Ergebnisdateien.

\subsubsection{Kommunikations- und I/O-Overhead}
\label{subsubsec:overhead}

Der Overhead dieses Ansatzes ist klein, aber nicht null. Pro Knoten fallen ein
SSH-Verbindungsaufbau und eine \texttt{scp}-Übertragung an. Übertragen wird jeweils
eine JSON-Datei, deren Größe mit der Zahl der empfangenen Treffer wächst; bei
einer Million geworfener Strahlen liegt sie im Bereich einiger Megabyte. Da diese
acht Übertragungen erst \emph{nach} der Rechenphase stattfinden und nicht mit ihr
überlappen, addieren sie sich voll zur Gesamtlaufzeit. Bei deutlich mehr Knoten
oder größeren Ergebnisdateien würde dieser sequenzielle Einsammel-Schritt zum
Flaschenhals; eine Überlappung von Rechnung und Transfer oder ein kompakteres
Ausgabeformat als das gut lesbare, aber platzintensive JSON wären dann die
naheliegenden Gegenmaßnahmen.

\subsubsection{Messung am Cluster}
\label{subsubsec:messung}

Die konkreten Laufzeiten hängen von der Laborhardware ab und sind am Cluster zu
messen. Um einen aussagekräftigen Speedup zu erhalten, wird in der
Partitionierungs-Betriebsart gemessen: In \texttt{input\_config.json} wird
\texttt{rays\_to\_cast} auf $8\cdot10^{6}/k$ gesetzt, sodass die Gesamtarbeit über
alle $k$ Knoten konstant bleibt. Dann wird \texttt{run\_cluster.sh} für
verschiedene Knotenzahlen mit der Shell-eigenen Zeitmessung ausgeführt:

\begin{lstlisting}[style=clusterlst, language=bash, caption={Messvorschrift für die Laufzeit am Cluster.}]
time ./run_cluster.sh        # Wanduhrzeit fuer den kompletten verteilten Lauf
\end{lstlisting}

Tabelle~\ref{tab:speedup} nimmt die Messwerte auf. Als Referenz dient der Lauf mit
einem einzelnen Knoten; der Speedup ergibt sich als Verhältnis der Wanduhrzeiten.

\begin{table}[htbp]
    \centering
    \caption{Wanduhrzeit und Speedup des verteilten Laufs (am Cluster zu messen).}
    \label{tab:speedup}
    \begin{tabular}{r r r}
        \toprule
        Knoten $k$ & Wanduhrzeit [s] & Speedup $S$ \\
        \midrule
        1 & \dots & $1{,}00$ \\
        2 & \dots & \dots \\
        4 & \dots & \dots \\
        8 & \dots & \dots \\
        \bottomrule
    \end{tabular}
\end{table}

\subsection{Fehlertoleranz und Sicherheit}
\label{subsec:cluster-fehler}

Zwei Querschnittsthemen verdienen eigene Betrachtung. Beim \textbf{Sicherheits}aspekt
ist die Schlüssel-basierte SSH-Authentifizierung zentral: Der Manager benötigt
passwortlosen Zugriff auf jeden Knoten, was über hinterlegte öffentliche Schlüssel
gelöst ist. Der private Schlüssel des Managers ist damit das entscheidende
Geheimnis; wer ihn besitzt, kann auf allen Knoten Code ausführen. In einem
abgeschlossenen Labornetz ist das vertretbar, für einen offeneren Betrieb wäre eine
engere Rechtevergabe angebracht.

Beim Thema \textbf{Fehlertoleranz} zeigt der aktuelle Stand seine deutlichste
Schwäche, die im folgenden Abschnitt zusammen mit den übrigen Grenzen diskutiert
wird. Der Kern des Problems: Ein ausgefallener Knoten wird nicht erkannt, sondern
schlicht übergangen, sodass die Stichprobengröße unbemerkt sinkt.

\subsection{Grenzen des aktuellen Standes und Ausblick}
\label{subsec:cluster-grenzen}

Die Verteilungsschicht ist als funktionsfähiger Prototyp umgesetzt. Für einen
robusten Betrieb bestehen mehrere klar benannte Schwachstellen, die zugleich die
nächsten Entwicklungsschritte vorgeben.

\begin{itemize}
    \item \textbf{Fest verdrahtete Konfiguration.} IP-Adressen, Knotennamen, die
    Knotenzahl und die Ablagepfade stehen direkt im Skript. Ein Umzug auf andere
    Hardware oder eine andere Knotenzahl erfordert Änderungen am Quelltext. Eine
    Auslagerung dieser Werte in eine Konfigurationsdatei wäre der naheliegende
    erste Schritt.

    \item \textbf{Kopierfehler bei Knoten~6.} Im Einsammel-Schritt wird das
    Ergebnis von Knoten~6 versehentlich nach \texttt{ir\_node5.json} kopiert und
    überschreibt damit das Ergebnis von Knoten~5. Effektiv fließen so nur sieben
    statt acht unabhängige Stichproben in die Aggregation ein. Der Fehler ist ein
    Tippfehler im Zielnamen und leicht zu beheben, führt aber vor Augen, wie
    fehleranfällig die von Hand ausgeschriebene Knotenliste ist -- ein weiteres
    Argument für die genannte Auslagerung in eine Schleife über eine
    Konfigurationsliste.

    \item \textbf{Keine Fehlererkennung.} Fällt ein Knoten aus oder liefert er kein
    Ergebnis, schlägt der zugehörige \texttt{scp}-Aufruf fehl, der Gesamtlauf setzt
    sich aber fort und aggregiert stillschweigend über die vorhandenen Dateien. Ein
    Knotenausfall bleibt damit unbemerkt und verfälscht die Stichprobengröße. Eine
    Prüfung der Rückgabecodes und eine Warnung bei fehlenden Ergebnissen wären hier
    angebracht.

    \item \textbf{Manuelle Kopplung an die DSP-Pipeline.} Die Aggregation erzeugt
    \texttt{ir\_merged.json}, die Python-DSP-Pipeline liest jedoch
    \texttt{ir\_output.json}. Der Übergang zwischen verteiltem Lauf und Faltung ist
    derzeit ein manueller Schritt. Eine saubere Integration würde den Cluster-Lauf
    direkt aus der Orchestrierung heraus anstoßen und die zusammengeführte Datei
    automatisch weiterreichen.

    \item \textbf{Statische Verteilung.} Der Replikat-Ansatz mit gleich großen
    Aufträgen ist einfach und robust, setzt aber homogene Hardware voraus und lastet
    unterschiedlich schnelle Knoten nicht unterschiedlich aus. Ein Modell mit
    dynamischer Arbeitsverteilung -- ein Manager vergibt Ray-Kontingente, sobald ein
    Knoten frei wird -- würde heterogene Hardware besser nutzen und die Knotenzahl
    flexibel machen.
\end{itemize}

Alle genannten Punkte betreffen Robustheit und Bedienbarkeit, nicht die
grundsätzliche Korrektheit des Verfahrens. Die verteilte Ausführung liefert eine
gültige, statistisch dichtere Impulsantwort, und das Manager-Worker-Muster über
SSH und \texttt{scp} erfüllt seinen Zweck im Laborbetrieb.

\subsection{Zusammenfassung}
\label{subsec:cluster-zusammenfassung}

Dieses Kapitel hat die verteilte Ausführung der Strahlverfolgung auf dem
Labor-Cluster beschrieben. Weil die Simulation embarrassingly parallel ist, genügt
ein Replikat-Ansatz: Alle acht Knoten rechnen denselben Auftrag mit unabhängigen
Zufallsströmen, was statistisch einer achtfach größeren Simulation entspricht und
den Schätzfehler der Impulsantwort um rund $\sqrt{8}$ senkt -- ein Verhalten, das
das numerische Experiment in Abbildung~\ref{fig:mc-convergence} bestätigt. Ein
Shell-Skript startet die Knoten per SSH parallel, synchronisiert über \texttt{wait}
und sammelt die Ergebnisse per \texttt{scp} ein; \texttt{merge\_results.py}
verkettet die Trefferlisten zu einer gemeinsamen Datei im gewohnten Schema. Der
Prototyp erfüllt seine Aufgabe, weist aber in Konfigurierbarkeit, Fehlererkennung
und der Anbindung an die DSP-Pipeline die diskutierten Schwächen auf, die den
Rahmen für die Weiterentwicklung abstecken.
 im Master.
% Benoetigte Pakete: amsmath, listings, graphicx, booktabs.
% Fuer saubere Optik im Master: fontenc, babel[ngerman], microtype.
% Abbildungen: latexDocs/figures/cluster_arch.png, mc_convergence.png
%
% STANDALONE-VORSCHAU:
%   \documentclass[12pt]{article}
%   \usepackage[utf8]{inputenc}\usepackage[T1]{fontenc}
%   \usepackage[ngerman]{babel}\usepackage{microtype}
%   \usepackage{amsmath}\usepackage{listings}
%   \usepackage{graphicx}\usepackage{booktabs}
%   \begin{document}
%   % =====================================================================
% Kapitel: Verteilte Ausfuehrung auf dem Rechencluster
% Autor: Mier
%
% EINBINDEN: \input{Kapitel_Cluster_Mier} im Master.
% Benoetigte Pakete: amsmath, listings, graphicx, booktabs.
% Fuer saubere Optik im Master: fontenc, babel[ngerman], microtype.
% Abbildungen: latexDocs/figures/cluster_arch.png, mc_convergence.png
%
% STANDALONE-VORSCHAU:
%   \documentclass[12pt]{article}
%   \usepackage[utf8]{inputenc}\usepackage[T1]{fontenc}
%   \usepackage[ngerman]{babel}\usepackage{microtype}
%   \usepackage{amsmath}\usepackage{listings}
%   \usepackage{graphicx}\usepackage{booktabs}
%   \begin{document}
%   \input{Kapitel_Cluster_Mier}
%   \end{document}
% =====================================================================

\lstdefinestyle{clusterlst}{
    basicstyle=\ttfamily\small,
    showstringspaces=false,
    breaklines=true,
    numbers=left,
    numberstyle=\tiny,
    frame=single,
    tabsize=2
}

%======================================================================
\section{Verteilte Ausführung auf dem Rechencluster}
\label{sec:cluster}

\subsection{Einordnung und Aufgabenstellung}
\label{subsec:cluster-einordnung}

Die Strahlverfolgung ist der teuerste Schritt der gesamten Pipeline. Für eine
brauchbare Impulsantwort müssen Hunderttausende bis Millionen einzelner
Schallstrahlen durch die Raumgeometrie verfolgt werden, und jeder dieser Strahlen
verlangt mehrere Schnitt- und Reflexionsberechnungen. Der entscheidende Vorteil
dabei: Kein Strahl hängt vom Ergebnis eines anderen ab. Das Problem ist damit
\emph{embarrassingly parallel} und lässt sich ohne Kommunikation zwischen den
Teilrechnungen auf beliebig viele Rechner verteilen.

Dieses Kapitel dokumentiert die Verteilungsschicht des Systems. Sie startet die
Rust-Simulation gleichzeitig auf allen acht Knoten des Labor-Clusters, wartet auf
deren Abschluss, sammelt die Teilergebnisse ein und fügt sie zu einer einzigen,
statistisch dichteren Impulsantwort zusammen. Aus Sicht der übrigen Pipeline
ändert sich dabei nichts: Am Ende steht wieder eine Datei im gewohnten Schema,
die die Signalverarbeitung unverändert weiterverarbeiten kann.

Zur Abgrenzung von den Nachbarkapiteln gehören zwei Punkte. Erstens ist die
\emph{prozessinterne} Nebenläufigkeit der Rust-Engine, also die Verteilung auf
die Kerne eines einzelnen Rechners mittels \texttt{rayon}, an anderer Stelle
beschrieben und bleibt hier außen vor. Sie parallelisiert innerhalb eines
Prozesses; dieses Kapitel parallelisiert über Prozess- und Rechnergrenzen hinweg.
Zweitens geht es nicht um den einzelnen Subprozess-Aufruf aus der Python-Pipeline,
sondern um die Ebene darüber: die gleichzeitige Ausführung auf mehreren physischen
Knoten und die anschließende Aggregation ihrer Ergebnisse.

\subsection{Theoretischer Hintergrund}
\label{subsec:cluster-theorie}

\subsubsection{Monte-Carlo-Strahlverfolgung}
\label{subsubsec:mc}

Die Engine schätzt die Impulsantwort eines Raumes stochastisch. Sie verschickt
zufällig gerichtete Strahlen von der Schallquelle, verfolgt ihre Reflexionen und
registriert, wann und mit welchem Restdruck sie den Empfänger erreichen. Jeder
einzelne Strahl ist eine Zufallsstichprobe; erst die Summe vieler Strahlen ergibt
ein aussagekräftiges Bild der Raumakustik. Solche Monte-Carlo-Verfahren haben eine
charakteristische Eigenschaft: Ihr Schätzfehler sinkt mit der Wurzel der
Stichprobengröße. Verdoppelt man die Zahl der Strahlen, halbiert sich der Fehler
also nicht, sondern sinkt nur um den Faktor $\sqrt{2}$. Genau diese langsame
Konvergenz macht große Strahlzahlen nötig und motiviert die Verteilung auf mehrere
Rechner.

\subsubsection{Zwei Wege der Verteilung}
\label{subsubsec:strategien}

Ein Monte-Carlo-Verfahren lässt sich auf zwei grundsätzlich verschiedene Arten
verteilen. Tabelle~\ref{tab:strategien} stellt sie gegenüber.

\begin{table}[htbp]
    \centering
    \caption{Zwei Strategien zur Verteilung der Strahlverfolgung.}
    \label{tab:strategien}
    \begin{tabular}{p{3.2cm} p{5.2cm} p{4.2cm}}
        \toprule
        & \textbf{Work-Partitioning} & \textbf{Replikat-Parallelismus} \\
        \midrule
        Prinzip & Gesamtzahl der Strahlen wird in disjunkte Anteile zerlegt & Jeder Knoten rechnet denselben Auftrag mit eigenem Zufallsstrom \\
        Koordination & Zuteilung der Anteile nötig & keine, jeder Knoten ist autark \\
        Zufallszahlen & ein globaler Strom, aufgeteilt & unabhängige Ströme pro Knoten \\
        Robustheit & Ausfall trifft einen festen Anteil & Ausfall kostet nur Stichprobengröße \\
        \bottomrule
    \end{tabular}
\end{table}

Das umgesetzte System verwendet den zweiten Weg. Alle acht Knoten starten die
Simulation mit derselben Konfigurationsdatei \texttt{input\_config.json} und
werfen dennoch jeweils andere Strahlen, weil jeder Prozess seinen eigenen
Zufallsstrom besitzt. Der Ansatz ist bewusst schlicht gewählt: Er braucht keine
zentrale Zuteilung und keinen Austausch zwischen den Knoten, sondern nur, dass am
Ende alle Ergebnisse eingesammelt werden.

\subsubsection{Unabhängigkeit der Zufallsströme}
\label{subsubsec:rng}

Die Korrektheit des Replikat-Ansatzes steht und fällt mit einer Bedingung: Die
Zufallszahlen auf den Knoten müssen unabhängig sein. Würde jeder Knoten dieselbe
Folge zufälliger Richtungen erzeugen, so lieferten alle acht exakt dasselbe
Ergebnis, und das Zusammenführen brächte keinen Erkenntnisgewinn, sondern nur
achtfach dieselben Strahlen. Die Rust-Engine schließt das aus, indem sie ihren
Zufallsstrom pro Thread bzw. Prozess frisch aus der Systementropie zieht:

\begin{lstlisting}[style=clusterlst, language=Rust, caption={Initialisierung des Zufallsstroms pro Prozess (\texttt{simulation.rs}).}]
let mut rng = rand::rng();                       // Omnidirectional, random rays
let random_angle: f32 = rng.random_range(0.0..(2.0 * PI));
let direction = Vec2::new(random_angle.cos(), random_angle.sin());
\end{lstlisting}

Weil jeder der acht Prozesse auf seinem Knoten einen anderen Startzustand erhält,
erzeugt jeder eine andere Stichprobe. Acht Knoten mit je einer Million Strahlen
liefern damit acht Millionen \emph{unabhängige} Stichproben derselben Raumakustik.

\subsubsection{Statistische Konvergenz}
\label{subsubsec:konvergenz}

Statistisch verhält sich der verteilte Lauf wie eine einzige, achtfach größere
Simulation. Für einen Monte-Carlo-Schätzer aus $N$ unabhängigen Teilstichproben
gleicher Größe sinkt die Standardabweichung des Mittelwerts mit
\[
    \sigma_N \;=\; \frac{\sigma_1}{\sqrt{N}} .
\]
Acht Knoten reduzieren die Streuung der geschätzten Impulsantwort demnach um den
Faktor $\sqrt{8} \approx 2{,}83$ gegenüber einem einzelnen Knoten.
Abbildung~\ref{fig:mc-convergence} bestätigt dieses Verhalten in einem numerischen
Experiment: Aufgetragen ist der mittlere quadratische Fehler der geschätzten
Impulsantwort über der Knotenzahl. Die Messpunkte folgen sauber der
$1/\sqrt{N}$-Referenzkurve; von einem auf acht Knoten sinkt der Fehler um den
Faktor $2{,}78$, praktisch exakt der theoretisch erwartete Wert.

\begin{figure}[htbp]
    \centering
    \includegraphics[width=0.8\textwidth]{figures/mc_convergence}
    \caption{Mittlerer quadratischer Fehler der geschätzten Impulsantwort über der
    Knotenzahl $N$ (je $5\cdot10^4$ Strahlen, doppelt logarithmisch). Die
    Messpunkte folgen der $1/\sqrt{N}$-Referenz; acht Knoten senken den Fehler um
    den Faktor $2{,}78 \approx \sqrt{8}$.}
    \label{fig:mc-convergence}
\end{figure}

Bemerkenswert ist die Kehrseite dieser Kurve: Um den Fehler zu halbieren, braucht
es die vierfache Knotenzahl. Der Nutzen zusätzlicher Knoten nimmt also ab, was bei
der Wahl der Clustergröße zu berücksichtigen ist. Für die hier angestrebte
Genauigkeit sind acht Knoten ein sinnvoller Kompromiss zwischen Aufwand und
Qualität.

\subsection{Architektur des Clusters}
\label{subsec:cluster-arch}

Das System folgt einem \emph{Manager-Worker}-Muster
(Abbildung~\ref{fig:cluster-arch}). Ein Manager-Knoten koordiniert den Ablauf,
acht Worker-Knoten führen die eigentliche Simulation aus. Der Manager selbst
rechnet nicht mit, sondern startet die Worker, wartet auf sie und führt am Ende
die Ergebnisse zusammen.

\begin{figure}[htbp]
    \centering
    \includegraphics[width=0.82\textwidth]{figures/cluster_arch}
    \caption{Manager-Worker-Architektur. Der Manager startet die acht Knoten
    parallel per \texttt{ssh}, holt nach Abschluss jede Ergebnisdatei per
    \texttt{scp} zurück und führt sie mit \texttt{merge\_results.py} zusammen.}
    \label{fig:cluster-arch}
\end{figure}

Die Kommunikation läuft vollständig über Standard-Werkzeuge des Betriebssystems.
\texttt{ssh} startet die Prozesse aus der Ferne, \texttt{scp} überträgt die
Ergebnisdateien zurück. Ein eigenes Netzwerkprotokoll oder ein
Message-Passing-Framework kommt nicht zum Einsatz. Für einen Stapelbetrieb, bei
dem jeder Knoten genau einmal startet, eine Weile rechnet und ein einziges
Ergebnis zurückliefert, wäre solcher Aufwand nicht gerechtfertigt: Es gibt keine
laufende Zustandssynchronisation und keinen Nachrichtenaustausch während der
Rechnung, den ein Framework verwalten müsste.

Damit das funktioniert, müssen drei Voraussetzungen erfüllt sein. Auf jedem Knoten
liegt dieselbe vorkompilierte Binärdatei \texttt{rust\_service} samt der
Konfiguration \texttt{input\_config.json}. Der Manager kann sich per SSH ohne
interaktive Passworteingabe anmelden, was eine Schlüssel-basierte
Authentifizierung voraussetzt. Und die Knoten sind über feste Adressen im
Labornetz erreichbar.

\subsection{Ablaufsteuerung: \texttt{run\_cluster.sh}}
\label{subsec:cluster-run}

Das Skript \texttt{run\_cluster.sh} steuert den verteilten Lauf in drei Phasen:
paralleles Starten, Warten und Einsammeln.

\begin{lstlisting}[style=clusterlst, language=bash, caption={Verteilter Start und Einsammeln der Ergebnisse (gekürzt).}]
# Phase 1: alle Knoten gleichzeitig starten (Hintergrundprozesse)
ssh node1@139.6.65.29 "./rust_service --config-file input_config.json" &
ssh node2@139.6.65.18 "./rust_service --config-file input_config.json" &
# ... node3 bis node8 analog ...

wait                       # Phase 2: auf Abschluss ALLER Knoten warten

# Phase 3: Ergebnis jedes Knotens zum Manager kopieren
scp node1@139.6.65.29:~/ir_output.json ~/ir_node1.json
scp node2@139.6.65.18:~/ir_output.json ~/ir_node2.json
# ... node3 bis node8 analog ...
\end{lstlisting}

Zwei Eigenschaften der Shell nutzt der Ablauf gezielt. Das kaufmännische Und
(\texttt{\&}) am Zeilenende schickt jeden \texttt{ssh}-Aufruf in den Hintergrund.
Die acht Simulationen laufen dadurch nahezu gleichzeitig an, statt nacheinander
abgearbeitet zu werden. Das darauffolgende \texttt{wait} blockiert, bis
\emph{alle} Hintergrundprozesse beendet sind, und bildet so eine
Synchronisationsbarriere zwischen der Rechen- und der Einsammel-Phase. Erst wenn
der letzte Knoten fertig ist, beginnt \texttt{scp}, von jedem Knoten dessen
\texttt{ir\_output.json} zu holen und unter einem knotenspezifischen Namen auf dem
Manager abzulegen.

Diese Barriere hat eine wichtige Konsequenz für die Laufzeit: Der Gesamtlauf ist
so langsam wie der \emph{langsamste} Knoten. Solange die Hardware homogen ist,
fällt das kaum ins Gewicht; auf ungleichen Knoten würde ein einzelner langsamer
Rechner den ganzen Lauf ausbremsen, während die übrigen sieben bereits fertig
warten. Dieses Straggler-Problem ist der zentrale Grund, warum ein statischer
Replikat-Ansatz mit gleich großen Aufträgen homogene Hardware voraussetzt.

\subsection{Aggregation der Teilergebnisse: \texttt{merge\_results.py}}
\label{subsec:cluster-merge}

Sobald alle Teilergebnisse auf dem Manager liegen, führt \texttt{merge\_results.py}
sie zu einer einzigen Impulsantwort zusammen. Jeder Knoten liefert dasselbe
Ausgabeschema, nämlich Metadaten sowie zwei positionsgleiche Listen aus
Verzögerungen und Drücken. Die Aggregation besteht deshalb im Kern aus einer
schlichten Verkettung dieser Listen:

\begin{lstlisting}[style=clusterlst, language=Python, caption={Zusammenführen der Knoten-Ergebnisse, gekürzt (\texttt{merge\_results.py}).}]
all_delays, all_pressures = [], []
for f in sorted(glob.glob("/home/manager/ir_node*.json")):
    with open(f) as file:
        data = json.load(file)
        all_delays    += data["hits"]["delays_seconds"]
        all_pressures += data["hits"]["pressures"]

merged = {
    "metadata": {"sample_rate": 44100, "rays_cast": 1000000 * 8,
                 "rays_received": len(all_delays), "room_name": "MVP_Test_Room_1"},
    "hits": {"delays_seconds": all_delays, "pressures": all_pressures},
}
json.dump(merged, open("/home/manager/ir_merged.json", "w"), indent=2)
\end{lstlisting}

Dass die Verkettung hier die richtige Operation ist, folgt direkt aus dem
Replikat-Modell. Weil die einzelnen Läufe unabhängige Stichproben desselben Raums
sind, ist die Vereinigung aller Treffer statistisch identisch mit dem Ergebnis
eines einzigen Laufs mit acht Millionen Strahlen. Es findet keine Mittelung und
keine Gewichtung statt, sondern nur ein Aneinanderhängen. Das Feld
\texttt{rays\_cast} ist im Skript fest auf $8 \cdot 10^{6}$ gesetzt; solange der in
Abschnitt~\ref{subsec:cluster-grenzen} beschriebene Kopierfehler bei Knoten~6
besteht, fließen tatsächlich nur sieben Knoten in die Aggregation ein, und dieser
Metadatenwert überschätzt die reale Strahlzahl. \texttt{rays\_received} zählt
dagegen die tatsächlich empfangenen Treffer. Die
zusammengeführte Datei trägt exakt dasselbe Schema wie die Ausgabe eines einzelnen
Knotens und ist damit ohne weitere Anpassung für die DSP-Pipeline verwendbar.

\subsection{Datenfluss im Überblick}
\label{subsec:cluster-datenfluss}

Der gesamte verteilte Ablauf lässt sich in fünf Stufen zusammenfassen:

\begin{enumerate}
    \item Der Manager startet per SSH auf jedem der acht Knoten
    \texttt{rust\_service} mit identischer Konfiguration, alle parallel.
    \item Jeder Knoten simuliert unabhängig eine Million Strahlen und schreibt
    lokal seine \texttt{ir\_output.json}.
    \item \texttt{wait} synchronisiert: Die Einsammel-Phase beginnt erst, wenn der
    letzte Knoten fertig ist.
    \item \texttt{scp} überträgt jede \texttt{ir\_output.json} als
    \texttt{ir\_node$k$.json} zum Manager.
    \item \texttt{merge\_results.py} verkettet alle Trefferlisten zu einer
    \texttt{ir\_merged.json}, die anschließend die DSP-Pipeline verarbeitet.
\end{enumerate}

\subsection{Leistungsbetrachtung}
\label{subsec:cluster-leistung}

\subsubsection{Erwarteter Speedup}
\label{subsubsec:speedup}

Bei der Leistung sind zwei Sichtweisen zu trennen. In der umgesetzten
Replikat-Betriebsart rechnet jeder Knoten dieselbe Strahlzahl, sodass acht Knoten
in nahezu derselben Wanduhrzeit die achtfache Stichprobe liefern. Der Gewinn ist
hier keine kürzere Laufzeit, sondern die höhere Genauigkeit aus
Abbildung~\ref{fig:mc-convergence}; man spricht vom schwachen Skalieren im Sinne
von Gustafson, bei dem mit der Rechnerzahl die bearbeitete Problemgröße wächst,
nicht die Geschwindigkeit einer festen Aufgabe. Will man dagegen eine feste
Gesamtstrahlzahl in kürzerer Zeit erreichen, so verteilt man diese auf die Knoten,
sodass jeder nur $8\cdot10^{6}/k$ Strahlen rechnet. Erst in dieser
Partitionierungs-Betriebsart greift das Gesetz von Amdahl. Bezeichnet $p$ den parallelen
Anteil der Arbeit, hier die reine Strahlverfolgung, und $1-p$ den seriellen Rest
aus Prozessstart, Ergebnistransfer und Aggregation, so beträgt der theoretische
Speedup
\[
    S(k) \;=\; \frac{1}{(1-p) + \dfrac{p}{k}} .
\]
Da die Strahlverfolgung die Laufzeit klar dominiert, liegt $p$ nahe bei eins, und
$S(k)$ steigt bis $k=8$ nahezu linear. In der Praxis bleibt der reale Speedup
darunter, und zwar aus drei Gründen: dem bereits erwähnten Straggler-Effekt (der
Lauf ist so schnell wie der langsamste Knoten), dem Overhead für den
SSH-Verbindungsaufbau sowie der Übertragung und Zusammenführung der
Ergebnisdateien.

\subsubsection{Kommunikations- und I/O-Overhead}
\label{subsubsec:overhead}

Der Overhead dieses Ansatzes ist klein, aber nicht null. Pro Knoten fallen ein
SSH-Verbindungsaufbau und eine \texttt{scp}-Übertragung an. Übertragen wird jeweils
eine JSON-Datei, deren Größe mit der Zahl der empfangenen Treffer wächst; bei
einer Million geworfener Strahlen liegt sie im Bereich einiger Megabyte. Da diese
acht Übertragungen erst \emph{nach} der Rechenphase stattfinden und nicht mit ihr
überlappen, addieren sie sich voll zur Gesamtlaufzeit. Bei deutlich mehr Knoten
oder größeren Ergebnisdateien würde dieser sequenzielle Einsammel-Schritt zum
Flaschenhals; eine Überlappung von Rechnung und Transfer oder ein kompakteres
Ausgabeformat als das gut lesbare, aber platzintensive JSON wären dann die
naheliegenden Gegenmaßnahmen.

\subsubsection{Messung am Cluster}
\label{subsubsec:messung}

Die konkreten Laufzeiten hängen von der Laborhardware ab und sind am Cluster zu
messen. Um einen aussagekräftigen Speedup zu erhalten, wird in der
Partitionierungs-Betriebsart gemessen: In \texttt{input\_config.json} wird
\texttt{rays\_to\_cast} auf $8\cdot10^{6}/k$ gesetzt, sodass die Gesamtarbeit über
alle $k$ Knoten konstant bleibt. Dann wird \texttt{run\_cluster.sh} für
verschiedene Knotenzahlen mit der Shell-eigenen Zeitmessung ausgeführt:

\begin{lstlisting}[style=clusterlst, language=bash, caption={Messvorschrift für die Laufzeit am Cluster.}]
time ./run_cluster.sh        # Wanduhrzeit fuer den kompletten verteilten Lauf
\end{lstlisting}

Tabelle~\ref{tab:speedup} nimmt die Messwerte auf. Als Referenz dient der Lauf mit
einem einzelnen Knoten; der Speedup ergibt sich als Verhältnis der Wanduhrzeiten.

\begin{table}[htbp]
    \centering
    \caption{Wanduhrzeit und Speedup des verteilten Laufs (am Cluster zu messen).}
    \label{tab:speedup}
    \begin{tabular}{r r r}
        \toprule
        Knoten $k$ & Wanduhrzeit [s] & Speedup $S$ \\
        \midrule
        1 & \dots & $1{,}00$ \\
        2 & \dots & \dots \\
        4 & \dots & \dots \\
        8 & \dots & \dots \\
        \bottomrule
    \end{tabular}
\end{table}

\subsection{Fehlertoleranz und Sicherheit}
\label{subsec:cluster-fehler}

Zwei Querschnittsthemen verdienen eigene Betrachtung. Beim \textbf{Sicherheits}aspekt
ist die Schlüssel-basierte SSH-Authentifizierung zentral: Der Manager benötigt
passwortlosen Zugriff auf jeden Knoten, was über hinterlegte öffentliche Schlüssel
gelöst ist. Der private Schlüssel des Managers ist damit das entscheidende
Geheimnis; wer ihn besitzt, kann auf allen Knoten Code ausführen. In einem
abgeschlossenen Labornetz ist das vertretbar, für einen offeneren Betrieb wäre eine
engere Rechtevergabe angebracht.

Beim Thema \textbf{Fehlertoleranz} zeigt der aktuelle Stand seine deutlichste
Schwäche, die im folgenden Abschnitt zusammen mit den übrigen Grenzen diskutiert
wird. Der Kern des Problems: Ein ausgefallener Knoten wird nicht erkannt, sondern
schlicht übergangen, sodass die Stichprobengröße unbemerkt sinkt.

\subsection{Grenzen des aktuellen Standes und Ausblick}
\label{subsec:cluster-grenzen}

Die Verteilungsschicht ist als funktionsfähiger Prototyp umgesetzt. Für einen
robusten Betrieb bestehen mehrere klar benannte Schwachstellen, die zugleich die
nächsten Entwicklungsschritte vorgeben.

\begin{itemize}
    \item \textbf{Fest verdrahtete Konfiguration.} IP-Adressen, Knotennamen, die
    Knotenzahl und die Ablagepfade stehen direkt im Skript. Ein Umzug auf andere
    Hardware oder eine andere Knotenzahl erfordert Änderungen am Quelltext. Eine
    Auslagerung dieser Werte in eine Konfigurationsdatei wäre der naheliegende
    erste Schritt.

    \item \textbf{Kopierfehler bei Knoten~6.} Im Einsammel-Schritt wird das
    Ergebnis von Knoten~6 versehentlich nach \texttt{ir\_node5.json} kopiert und
    überschreibt damit das Ergebnis von Knoten~5. Effektiv fließen so nur sieben
    statt acht unabhängige Stichproben in die Aggregation ein. Der Fehler ist ein
    Tippfehler im Zielnamen und leicht zu beheben, führt aber vor Augen, wie
    fehleranfällig die von Hand ausgeschriebene Knotenliste ist -- ein weiteres
    Argument für die genannte Auslagerung in eine Schleife über eine
    Konfigurationsliste.

    \item \textbf{Keine Fehlererkennung.} Fällt ein Knoten aus oder liefert er kein
    Ergebnis, schlägt der zugehörige \texttt{scp}-Aufruf fehl, der Gesamtlauf setzt
    sich aber fort und aggregiert stillschweigend über die vorhandenen Dateien. Ein
    Knotenausfall bleibt damit unbemerkt und verfälscht die Stichprobengröße. Eine
    Prüfung der Rückgabecodes und eine Warnung bei fehlenden Ergebnissen wären hier
    angebracht.

    \item \textbf{Manuelle Kopplung an die DSP-Pipeline.} Die Aggregation erzeugt
    \texttt{ir\_merged.json}, die Python-DSP-Pipeline liest jedoch
    \texttt{ir\_output.json}. Der Übergang zwischen verteiltem Lauf und Faltung ist
    derzeit ein manueller Schritt. Eine saubere Integration würde den Cluster-Lauf
    direkt aus der Orchestrierung heraus anstoßen und die zusammengeführte Datei
    automatisch weiterreichen.

    \item \textbf{Statische Verteilung.} Der Replikat-Ansatz mit gleich großen
    Aufträgen ist einfach und robust, setzt aber homogene Hardware voraus und lastet
    unterschiedlich schnelle Knoten nicht unterschiedlich aus. Ein Modell mit
    dynamischer Arbeitsverteilung -- ein Manager vergibt Ray-Kontingente, sobald ein
    Knoten frei wird -- würde heterogene Hardware besser nutzen und die Knotenzahl
    flexibel machen.
\end{itemize}

Alle genannten Punkte betreffen Robustheit und Bedienbarkeit, nicht die
grundsätzliche Korrektheit des Verfahrens. Die verteilte Ausführung liefert eine
gültige, statistisch dichtere Impulsantwort, und das Manager-Worker-Muster über
SSH und \texttt{scp} erfüllt seinen Zweck im Laborbetrieb.

\subsection{Zusammenfassung}
\label{subsec:cluster-zusammenfassung}

Dieses Kapitel hat die verteilte Ausführung der Strahlverfolgung auf dem
Labor-Cluster beschrieben. Weil die Simulation embarrassingly parallel ist, genügt
ein Replikat-Ansatz: Alle acht Knoten rechnen denselben Auftrag mit unabhängigen
Zufallsströmen, was statistisch einer achtfach größeren Simulation entspricht und
den Schätzfehler der Impulsantwort um rund $\sqrt{8}$ senkt -- ein Verhalten, das
das numerische Experiment in Abbildung~\ref{fig:mc-convergence} bestätigt. Ein
Shell-Skript startet die Knoten per SSH parallel, synchronisiert über \texttt{wait}
und sammelt die Ergebnisse per \texttt{scp} ein; \texttt{merge\_results.py}
verkettet die Trefferlisten zu einer gemeinsamen Datei im gewohnten Schema. Der
Prototyp erfüllt seine Aufgabe, weist aber in Konfigurierbarkeit, Fehlererkennung
und der Anbindung an die DSP-Pipeline die diskutierten Schwächen auf, die den
Rahmen für die Weiterentwicklung abstecken.

%   \end{document}
% =====================================================================

\lstdefinestyle{clusterlst}{
    basicstyle=\ttfamily\small,
    showstringspaces=false,
    breaklines=true,
    numbers=left,
    numberstyle=\tiny,
    frame=single,
    tabsize=2
}

%======================================================================
\section{Verteilte Ausführung auf dem Rechencluster}
\label{sec:cluster}

\subsection{Einordnung und Aufgabenstellung}
\label{subsec:cluster-einordnung}

Die Strahlverfolgung ist der teuerste Schritt der gesamten Pipeline. Für eine
brauchbare Impulsantwort müssen Hunderttausende bis Millionen einzelner
Schallstrahlen durch die Raumgeometrie verfolgt werden, und jeder dieser Strahlen
verlangt mehrere Schnitt- und Reflexionsberechnungen. Der entscheidende Vorteil
dabei: Kein Strahl hängt vom Ergebnis eines anderen ab. Das Problem ist damit
\emph{embarrassingly parallel} und lässt sich ohne Kommunikation zwischen den
Teilrechnungen auf beliebig viele Rechner verteilen.

Dieses Kapitel dokumentiert die Verteilungsschicht des Systems. Sie startet die
Rust-Simulation gleichzeitig auf allen acht Knoten des Labor-Clusters, wartet auf
deren Abschluss, sammelt die Teilergebnisse ein und fügt sie zu einer einzigen,
statistisch dichteren Impulsantwort zusammen. Aus Sicht der übrigen Pipeline
ändert sich dabei nichts: Am Ende steht wieder eine Datei im gewohnten Schema,
die die Signalverarbeitung unverändert weiterverarbeiten kann.

Zur Abgrenzung von den Nachbarkapiteln gehören zwei Punkte. Erstens ist die
\emph{prozessinterne} Nebenläufigkeit der Rust-Engine, also die Verteilung auf
die Kerne eines einzelnen Rechners mittels \texttt{rayon}, an anderer Stelle
beschrieben und bleibt hier außen vor. Sie parallelisiert innerhalb eines
Prozesses; dieses Kapitel parallelisiert über Prozess- und Rechnergrenzen hinweg.
Zweitens geht es nicht um den einzelnen Subprozess-Aufruf aus der Python-Pipeline,
sondern um die Ebene darüber: die gleichzeitige Ausführung auf mehreren physischen
Knoten und die anschließende Aggregation ihrer Ergebnisse.

\subsection{Theoretischer Hintergrund}
\label{subsec:cluster-theorie}

\subsubsection{Monte-Carlo-Strahlverfolgung}
\label{subsubsec:mc}

Die Engine schätzt die Impulsantwort eines Raumes stochastisch. Sie verschickt
zufällig gerichtete Strahlen von der Schallquelle, verfolgt ihre Reflexionen und
registriert, wann und mit welchem Restdruck sie den Empfänger erreichen. Jeder
einzelne Strahl ist eine Zufallsstichprobe; erst die Summe vieler Strahlen ergibt
ein aussagekräftiges Bild der Raumakustik. Solche Monte-Carlo-Verfahren haben eine
charakteristische Eigenschaft: Ihr Schätzfehler sinkt mit der Wurzel der
Stichprobengröße. Verdoppelt man die Zahl der Strahlen, halbiert sich der Fehler
also nicht, sondern sinkt nur um den Faktor $\sqrt{2}$. Genau diese langsame
Konvergenz macht große Strahlzahlen nötig und motiviert die Verteilung auf mehrere
Rechner.

\subsubsection{Zwei Wege der Verteilung}
\label{subsubsec:strategien}

Ein Monte-Carlo-Verfahren lässt sich auf zwei grundsätzlich verschiedene Arten
verteilen. Tabelle~\ref{tab:strategien} stellt sie gegenüber.

\begin{table}[htbp]
    \centering
    \caption{Zwei Strategien zur Verteilung der Strahlverfolgung.}
    \label{tab:strategien}
    \begin{tabular}{p{3.2cm} p{5.2cm} p{4.2cm}}
        \toprule
        & \textbf{Work-Partitioning} & \textbf{Replikat-Parallelismus} \\
        \midrule
        Prinzip & Gesamtzahl der Strahlen wird in disjunkte Anteile zerlegt & Jeder Knoten rechnet denselben Auftrag mit eigenem Zufallsstrom \\
        Koordination & Zuteilung der Anteile nötig & keine, jeder Knoten ist autark \\
        Zufallszahlen & ein globaler Strom, aufgeteilt & unabhängige Ströme pro Knoten \\
        Robustheit & Ausfall trifft einen festen Anteil & Ausfall kostet nur Stichprobengröße \\
        \bottomrule
    \end{tabular}
\end{table}

Das umgesetzte System verwendet den zweiten Weg. Alle acht Knoten starten die
Simulation mit derselben Konfigurationsdatei \texttt{input\_config.json} und
werfen dennoch jeweils andere Strahlen, weil jeder Prozess seinen eigenen
Zufallsstrom besitzt. Der Ansatz ist bewusst schlicht gewählt: Er braucht keine
zentrale Zuteilung und keinen Austausch zwischen den Knoten, sondern nur, dass am
Ende alle Ergebnisse eingesammelt werden.

\subsubsection{Unabhängigkeit der Zufallsströme}
\label{subsubsec:rng}

Die Korrektheit des Replikat-Ansatzes steht und fällt mit einer Bedingung: Die
Zufallszahlen auf den Knoten müssen unabhängig sein. Würde jeder Knoten dieselbe
Folge zufälliger Richtungen erzeugen, so lieferten alle acht exakt dasselbe
Ergebnis, und das Zusammenführen brächte keinen Erkenntnisgewinn, sondern nur
achtfach dieselben Strahlen. Die Rust-Engine schließt das aus, indem sie ihren
Zufallsstrom pro Thread bzw. Prozess frisch aus der Systementropie zieht:

\begin{lstlisting}[style=clusterlst, language=Rust, caption={Initialisierung des Zufallsstroms pro Prozess (\texttt{simulation.rs}).}]
let mut rng = rand::rng();                       // Omnidirectional, random rays
let random_angle: f32 = rng.random_range(0.0..(2.0 * PI));
let direction = Vec2::new(random_angle.cos(), random_angle.sin());
\end{lstlisting}

Weil jeder der acht Prozesse auf seinem Knoten einen anderen Startzustand erhält,
erzeugt jeder eine andere Stichprobe. Acht Knoten mit je einer Million Strahlen
liefern damit acht Millionen \emph{unabhängige} Stichproben derselben Raumakustik.

\subsubsection{Statistische Konvergenz}
\label{subsubsec:konvergenz}

Statistisch verhält sich der verteilte Lauf wie eine einzige, achtfach größere
Simulation. Für einen Monte-Carlo-Schätzer aus $N$ unabhängigen Teilstichproben
gleicher Größe sinkt die Standardabweichung des Mittelwerts mit
\[
    \sigma_N \;=\; \frac{\sigma_1}{\sqrt{N}} .
\]
Acht Knoten reduzieren die Streuung der geschätzten Impulsantwort demnach um den
Faktor $\sqrt{8} \approx 2{,}83$ gegenüber einem einzelnen Knoten.
Abbildung~\ref{fig:mc-convergence} bestätigt dieses Verhalten in einem numerischen
Experiment: Aufgetragen ist der mittlere quadratische Fehler der geschätzten
Impulsantwort über der Knotenzahl. Die Messpunkte folgen sauber der
$1/\sqrt{N}$-Referenzkurve; von einem auf acht Knoten sinkt der Fehler um den
Faktor $2{,}78$, praktisch exakt der theoretisch erwartete Wert.

\begin{figure}[htbp]
    \centering
    \includegraphics[width=0.8\textwidth]{figures/mc_convergence}
    \caption{Mittlerer quadratischer Fehler der geschätzten Impulsantwort über der
    Knotenzahl $N$ (je $5\cdot10^4$ Strahlen, doppelt logarithmisch). Die
    Messpunkte folgen der $1/\sqrt{N}$-Referenz; acht Knoten senken den Fehler um
    den Faktor $2{,}78 \approx \sqrt{8}$.}
    \label{fig:mc-convergence}
\end{figure}

Bemerkenswert ist die Kehrseite dieser Kurve: Um den Fehler zu halbieren, braucht
es die vierfache Knotenzahl. Der Nutzen zusätzlicher Knoten nimmt also ab, was bei
der Wahl der Clustergröße zu berücksichtigen ist. Für die hier angestrebte
Genauigkeit sind acht Knoten ein sinnvoller Kompromiss zwischen Aufwand und
Qualität.

\subsection{Architektur des Clusters}
\label{subsec:cluster-arch}

Das System folgt einem \emph{Manager-Worker}-Muster
(Abbildung~\ref{fig:cluster-arch}). Ein Manager-Knoten koordiniert den Ablauf,
acht Worker-Knoten führen die eigentliche Simulation aus. Der Manager selbst
rechnet nicht mit, sondern startet die Worker, wartet auf sie und führt am Ende
die Ergebnisse zusammen.

\begin{figure}[htbp]
    \centering
    \includegraphics[width=0.82\textwidth]{figures/cluster_arch}
    \caption{Manager-Worker-Architektur. Der Manager startet die acht Knoten
    parallel per \texttt{ssh}, holt nach Abschluss jede Ergebnisdatei per
    \texttt{scp} zurück und führt sie mit \texttt{merge\_results.py} zusammen.}
    \label{fig:cluster-arch}
\end{figure}

Die Kommunikation läuft vollständig über Standard-Werkzeuge des Betriebssystems.
\texttt{ssh} startet die Prozesse aus der Ferne, \texttt{scp} überträgt die
Ergebnisdateien zurück. Ein eigenes Netzwerkprotokoll oder ein
Message-Passing-Framework kommt nicht zum Einsatz. Für einen Stapelbetrieb, bei
dem jeder Knoten genau einmal startet, eine Weile rechnet und ein einziges
Ergebnis zurückliefert, wäre solcher Aufwand nicht gerechtfertigt: Es gibt keine
laufende Zustandssynchronisation und keinen Nachrichtenaustausch während der
Rechnung, den ein Framework verwalten müsste.

Damit das funktioniert, müssen drei Voraussetzungen erfüllt sein. Auf jedem Knoten
liegt dieselbe vorkompilierte Binärdatei \texttt{rust\_service} samt der
Konfiguration \texttt{input\_config.json}. Der Manager kann sich per SSH ohne
interaktive Passworteingabe anmelden, was eine Schlüssel-basierte
Authentifizierung voraussetzt. Und die Knoten sind über feste Adressen im
Labornetz erreichbar.

\subsection{Ablaufsteuerung: \texttt{run\_cluster.sh}}
\label{subsec:cluster-run}

Das Skript \texttt{run\_cluster.sh} steuert den verteilten Lauf in drei Phasen:
paralleles Starten, Warten und Einsammeln.

\begin{lstlisting}[style=clusterlst, language=bash, caption={Verteilter Start und Einsammeln der Ergebnisse (gekürzt).}]
# Phase 1: alle Knoten gleichzeitig starten (Hintergrundprozesse)
ssh node1@139.6.65.29 "./rust_service --config-file input_config.json" &
ssh node2@139.6.65.18 "./rust_service --config-file input_config.json" &
# ... node3 bis node8 analog ...

wait                       # Phase 2: auf Abschluss ALLER Knoten warten

# Phase 3: Ergebnis jedes Knotens zum Manager kopieren
scp node1@139.6.65.29:~/ir_output.json ~/ir_node1.json
scp node2@139.6.65.18:~/ir_output.json ~/ir_node2.json
# ... node3 bis node8 analog ...
\end{lstlisting}

Zwei Eigenschaften der Shell nutzt der Ablauf gezielt. Das kaufmännische Und
(\texttt{\&}) am Zeilenende schickt jeden \texttt{ssh}-Aufruf in den Hintergrund.
Die acht Simulationen laufen dadurch nahezu gleichzeitig an, statt nacheinander
abgearbeitet zu werden. Das darauffolgende \texttt{wait} blockiert, bis
\emph{alle} Hintergrundprozesse beendet sind, und bildet so eine
Synchronisationsbarriere zwischen der Rechen- und der Einsammel-Phase. Erst wenn
der letzte Knoten fertig ist, beginnt \texttt{scp}, von jedem Knoten dessen
\texttt{ir\_output.json} zu holen und unter einem knotenspezifischen Namen auf dem
Manager abzulegen.

Diese Barriere hat eine wichtige Konsequenz für die Laufzeit: Der Gesamtlauf ist
so langsam wie der \emph{langsamste} Knoten. Solange die Hardware homogen ist,
fällt das kaum ins Gewicht; auf ungleichen Knoten würde ein einzelner langsamer
Rechner den ganzen Lauf ausbremsen, während die übrigen sieben bereits fertig
warten. Dieses Straggler-Problem ist der zentrale Grund, warum ein statischer
Replikat-Ansatz mit gleich großen Aufträgen homogene Hardware voraussetzt.

\subsection{Aggregation der Teilergebnisse: \texttt{merge\_results.py}}
\label{subsec:cluster-merge}

Sobald alle Teilergebnisse auf dem Manager liegen, führt \texttt{merge\_results.py}
sie zu einer einzigen Impulsantwort zusammen. Jeder Knoten liefert dasselbe
Ausgabeschema, nämlich Metadaten sowie zwei positionsgleiche Listen aus
Verzögerungen und Drücken. Die Aggregation besteht deshalb im Kern aus einer
schlichten Verkettung dieser Listen:

\begin{lstlisting}[style=clusterlst, language=Python, caption={Zusammenführen der Knoten-Ergebnisse, gekürzt (\texttt{merge\_results.py}).}]
all_delays, all_pressures = [], []
for f in sorted(glob.glob("/home/manager/ir_node*.json")):
    with open(f) as file:
        data = json.load(file)
        all_delays    += data["hits"]["delays_seconds"]
        all_pressures += data["hits"]["pressures"]

merged = {
    "metadata": {"sample_rate": 44100, "rays_cast": 1000000 * 8,
                 "rays_received": len(all_delays), "room_name": "MVP_Test_Room_1"},
    "hits": {"delays_seconds": all_delays, "pressures": all_pressures},
}
json.dump(merged, open("/home/manager/ir_merged.json", "w"), indent=2)
\end{lstlisting}

Dass die Verkettung hier die richtige Operation ist, folgt direkt aus dem
Replikat-Modell. Weil die einzelnen Läufe unabhängige Stichproben desselben Raums
sind, ist die Vereinigung aller Treffer statistisch identisch mit dem Ergebnis
eines einzigen Laufs mit acht Millionen Strahlen. Es findet keine Mittelung und
keine Gewichtung statt, sondern nur ein Aneinanderhängen. Das Feld
\texttt{rays\_cast} ist im Skript fest auf $8 \cdot 10^{6}$ gesetzt; solange der in
Abschnitt~\ref{subsec:cluster-grenzen} beschriebene Kopierfehler bei Knoten~6
besteht, fließen tatsächlich nur sieben Knoten in die Aggregation ein, und dieser
Metadatenwert überschätzt die reale Strahlzahl. \texttt{rays\_received} zählt
dagegen die tatsächlich empfangenen Treffer. Die
zusammengeführte Datei trägt exakt dasselbe Schema wie die Ausgabe eines einzelnen
Knotens und ist damit ohne weitere Anpassung für die DSP-Pipeline verwendbar.

\subsection{Datenfluss im Überblick}
\label{subsec:cluster-datenfluss}

Der gesamte verteilte Ablauf lässt sich in fünf Stufen zusammenfassen:

\begin{enumerate}
    \item Der Manager startet per SSH auf jedem der acht Knoten
    \texttt{rust\_service} mit identischer Konfiguration, alle parallel.
    \item Jeder Knoten simuliert unabhängig eine Million Strahlen und schreibt
    lokal seine \texttt{ir\_output.json}.
    \item \texttt{wait} synchronisiert: Die Einsammel-Phase beginnt erst, wenn der
    letzte Knoten fertig ist.
    \item \texttt{scp} überträgt jede \texttt{ir\_output.json} als
    \texttt{ir\_node$k$.json} zum Manager.
    \item \texttt{merge\_results.py} verkettet alle Trefferlisten zu einer
    \texttt{ir\_merged.json}, die anschließend die DSP-Pipeline verarbeitet.
\end{enumerate}

\subsection{Leistungsbetrachtung}
\label{subsec:cluster-leistung}

\subsubsection{Erwarteter Speedup}
\label{subsubsec:speedup}

Bei der Leistung sind zwei Sichtweisen zu trennen. In der umgesetzten
Replikat-Betriebsart rechnet jeder Knoten dieselbe Strahlzahl, sodass acht Knoten
in nahezu derselben Wanduhrzeit die achtfache Stichprobe liefern. Der Gewinn ist
hier keine kürzere Laufzeit, sondern die höhere Genauigkeit aus
Abbildung~\ref{fig:mc-convergence}; man spricht vom schwachen Skalieren im Sinne
von Gustafson, bei dem mit der Rechnerzahl die bearbeitete Problemgröße wächst,
nicht die Geschwindigkeit einer festen Aufgabe. Will man dagegen eine feste
Gesamtstrahlzahl in kürzerer Zeit erreichen, so verteilt man diese auf die Knoten,
sodass jeder nur $8\cdot10^{6}/k$ Strahlen rechnet. Erst in dieser
Partitionierungs-Betriebsart greift das Gesetz von Amdahl. Bezeichnet $p$ den parallelen
Anteil der Arbeit, hier die reine Strahlverfolgung, und $1-p$ den seriellen Rest
aus Prozessstart, Ergebnistransfer und Aggregation, so beträgt der theoretische
Speedup
\[
    S(k) \;=\; \frac{1}{(1-p) + \dfrac{p}{k}} .
\]
Da die Strahlverfolgung die Laufzeit klar dominiert, liegt $p$ nahe bei eins, und
$S(k)$ steigt bis $k=8$ nahezu linear. In der Praxis bleibt der reale Speedup
darunter, und zwar aus drei Gründen: dem bereits erwähnten Straggler-Effekt (der
Lauf ist so schnell wie der langsamste Knoten), dem Overhead für den
SSH-Verbindungsaufbau sowie der Übertragung und Zusammenführung der
Ergebnisdateien.

\subsubsection{Kommunikations- und I/O-Overhead}
\label{subsubsec:overhead}

Der Overhead dieses Ansatzes ist klein, aber nicht null. Pro Knoten fallen ein
SSH-Verbindungsaufbau und eine \texttt{scp}-Übertragung an. Übertragen wird jeweils
eine JSON-Datei, deren Größe mit der Zahl der empfangenen Treffer wächst; bei
einer Million geworfener Strahlen liegt sie im Bereich einiger Megabyte. Da diese
acht Übertragungen erst \emph{nach} der Rechenphase stattfinden und nicht mit ihr
überlappen, addieren sie sich voll zur Gesamtlaufzeit. Bei deutlich mehr Knoten
oder größeren Ergebnisdateien würde dieser sequenzielle Einsammel-Schritt zum
Flaschenhals; eine Überlappung von Rechnung und Transfer oder ein kompakteres
Ausgabeformat als das gut lesbare, aber platzintensive JSON wären dann die
naheliegenden Gegenmaßnahmen.

\subsubsection{Messung am Cluster}
\label{subsubsec:messung}

Die konkreten Laufzeiten hängen von der Laborhardware ab und sind am Cluster zu
messen. Um einen aussagekräftigen Speedup zu erhalten, wird in der
Partitionierungs-Betriebsart gemessen: In \texttt{input\_config.json} wird
\texttt{rays\_to\_cast} auf $8\cdot10^{6}/k$ gesetzt, sodass die Gesamtarbeit über
alle $k$ Knoten konstant bleibt. Dann wird \texttt{run\_cluster.sh} für
verschiedene Knotenzahlen mit der Shell-eigenen Zeitmessung ausgeführt:

\begin{lstlisting}[style=clusterlst, language=bash, caption={Messvorschrift für die Laufzeit am Cluster.}]
time ./run_cluster.sh        # Wanduhrzeit fuer den kompletten verteilten Lauf
\end{lstlisting}

Tabelle~\ref{tab:speedup} nimmt die Messwerte auf. Als Referenz dient der Lauf mit
einem einzelnen Knoten; der Speedup ergibt sich als Verhältnis der Wanduhrzeiten.

\begin{table}[htbp]
    \centering
    \caption{Wanduhrzeit und Speedup des verteilten Laufs (am Cluster zu messen).}
    \label{tab:speedup}
    \begin{tabular}{r r r}
        \toprule
        Knoten $k$ & Wanduhrzeit [s] & Speedup $S$ \\
        \midrule
        1 & \dots & $1{,}00$ \\
        2 & \dots & \dots \\
        4 & \dots & \dots \\
        8 & \dots & \dots \\
        \bottomrule
    \end{tabular}
\end{table}

\subsection{Fehlertoleranz und Sicherheit}
\label{subsec:cluster-fehler}

Zwei Querschnittsthemen verdienen eigene Betrachtung. Beim \textbf{Sicherheits}aspekt
ist die Schlüssel-basierte SSH-Authentifizierung zentral: Der Manager benötigt
passwortlosen Zugriff auf jeden Knoten, was über hinterlegte öffentliche Schlüssel
gelöst ist. Der private Schlüssel des Managers ist damit das entscheidende
Geheimnis; wer ihn besitzt, kann auf allen Knoten Code ausführen. In einem
abgeschlossenen Labornetz ist das vertretbar, für einen offeneren Betrieb wäre eine
engere Rechtevergabe angebracht.

Beim Thema \textbf{Fehlertoleranz} zeigt der aktuelle Stand seine deutlichste
Schwäche, die im folgenden Abschnitt zusammen mit den übrigen Grenzen diskutiert
wird. Der Kern des Problems: Ein ausgefallener Knoten wird nicht erkannt, sondern
schlicht übergangen, sodass die Stichprobengröße unbemerkt sinkt.

\subsection{Grenzen des aktuellen Standes und Ausblick}
\label{subsec:cluster-grenzen}

Die Verteilungsschicht ist als funktionsfähiger Prototyp umgesetzt. Für einen
robusten Betrieb bestehen mehrere klar benannte Schwachstellen, die zugleich die
nächsten Entwicklungsschritte vorgeben.

\begin{itemize}
    \item \textbf{Fest verdrahtete Konfiguration.} IP-Adressen, Knotennamen, die
    Knotenzahl und die Ablagepfade stehen direkt im Skript. Ein Umzug auf andere
    Hardware oder eine andere Knotenzahl erfordert Änderungen am Quelltext. Eine
    Auslagerung dieser Werte in eine Konfigurationsdatei wäre der naheliegende
    erste Schritt.

    \item \textbf{Kopierfehler bei Knoten~6.} Im Einsammel-Schritt wird das
    Ergebnis von Knoten~6 versehentlich nach \texttt{ir\_node5.json} kopiert und
    überschreibt damit das Ergebnis von Knoten~5. Effektiv fließen so nur sieben
    statt acht unabhängige Stichproben in die Aggregation ein. Der Fehler ist ein
    Tippfehler im Zielnamen und leicht zu beheben, führt aber vor Augen, wie
    fehleranfällig die von Hand ausgeschriebene Knotenliste ist -- ein weiteres
    Argument für die genannte Auslagerung in eine Schleife über eine
    Konfigurationsliste.

    \item \textbf{Keine Fehlererkennung.} Fällt ein Knoten aus oder liefert er kein
    Ergebnis, schlägt der zugehörige \texttt{scp}-Aufruf fehl, der Gesamtlauf setzt
    sich aber fort und aggregiert stillschweigend über die vorhandenen Dateien. Ein
    Knotenausfall bleibt damit unbemerkt und verfälscht die Stichprobengröße. Eine
    Prüfung der Rückgabecodes und eine Warnung bei fehlenden Ergebnissen wären hier
    angebracht.

    \item \textbf{Manuelle Kopplung an die DSP-Pipeline.} Die Aggregation erzeugt
    \texttt{ir\_merged.json}, die Python-DSP-Pipeline liest jedoch
    \texttt{ir\_output.json}. Der Übergang zwischen verteiltem Lauf und Faltung ist
    derzeit ein manueller Schritt. Eine saubere Integration würde den Cluster-Lauf
    direkt aus der Orchestrierung heraus anstoßen und die zusammengeführte Datei
    automatisch weiterreichen.

    \item \textbf{Statische Verteilung.} Der Replikat-Ansatz mit gleich großen
    Aufträgen ist einfach und robust, setzt aber homogene Hardware voraus und lastet
    unterschiedlich schnelle Knoten nicht unterschiedlich aus. Ein Modell mit
    dynamischer Arbeitsverteilung -- ein Manager vergibt Ray-Kontingente, sobald ein
    Knoten frei wird -- würde heterogene Hardware besser nutzen und die Knotenzahl
    flexibel machen.
\end{itemize}

Alle genannten Punkte betreffen Robustheit und Bedienbarkeit, nicht die
grundsätzliche Korrektheit des Verfahrens. Die verteilte Ausführung liefert eine
gültige, statistisch dichtere Impulsantwort, und das Manager-Worker-Muster über
SSH und \texttt{scp} erfüllt seinen Zweck im Laborbetrieb.

\subsection{Zusammenfassung}
\label{subsec:cluster-zusammenfassung}

Dieses Kapitel hat die verteilte Ausführung der Strahlverfolgung auf dem
Labor-Cluster beschrieben. Weil die Simulation embarrassingly parallel ist, genügt
ein Replikat-Ansatz: Alle acht Knoten rechnen denselben Auftrag mit unabhängigen
Zufallsströmen, was statistisch einer achtfach größeren Simulation entspricht und
den Schätzfehler der Impulsantwort um rund $\sqrt{8}$ senkt -- ein Verhalten, das
das numerische Experiment in Abbildung~\ref{fig:mc-convergence} bestätigt. Ein
Shell-Skript startet die Knoten per SSH parallel, synchronisiert über \texttt{wait}
und sammelt die Ergebnisse per \texttt{scp} ein; \texttt{merge\_results.py}
verkettet die Trefferlisten zu einer gemeinsamen Datei im gewohnten Schema. Der
Prototyp erfüllt seine Aufgabe, weist aber in Konfigurierbarkeit, Fehlererkennung
und der Anbindung an die DSP-Pipeline die diskutierten Schwächen auf, die den
Rahmen für die Weiterentwicklung abstecken.
 im Master.
% Benoetigte Pakete: amsmath, listings, graphicx, booktabs.
% Fuer saubere Optik im Master: fontenc, babel[ngerman], microtype.
% Abbildungen: latexDocs/figures/cluster_arch.png, mc_convergence.png
%
% STANDALONE-VORSCHAU:
%   \documentclass[12pt]{article}
%   \usepackage[utf8]{inputenc}\usepackage[T1]{fontenc}
%   \usepackage[ngerman]{babel}\usepackage{microtype}
%   \usepackage{amsmath}\usepackage{listings}
%   \usepackage{graphicx}\usepackage{booktabs}
%   \begin{document}
%   % =====================================================================
% Kapitel: Verteilte Ausfuehrung auf dem Rechencluster
% Autor: Mier
%
% EINBINDEN: % =====================================================================
% Kapitel: Verteilte Ausfuehrung auf dem Rechencluster
% Autor: Mier
%
% EINBINDEN: \input{Kapitel_Cluster_Mier} im Master.
% Benoetigte Pakete: amsmath, listings, graphicx, booktabs.
% Fuer saubere Optik im Master: fontenc, babel[ngerman], microtype.
% Abbildungen: latexDocs/figures/cluster_arch.png, mc_convergence.png
%
% STANDALONE-VORSCHAU:
%   \documentclass[12pt]{article}
%   \usepackage[utf8]{inputenc}\usepackage[T1]{fontenc}
%   \usepackage[ngerman]{babel}\usepackage{microtype}
%   \usepackage{amsmath}\usepackage{listings}
%   \usepackage{graphicx}\usepackage{booktabs}
%   \begin{document}
%   \input{Kapitel_Cluster_Mier}
%   \end{document}
% =====================================================================

\lstdefinestyle{clusterlst}{
    basicstyle=\ttfamily\small,
    showstringspaces=false,
    breaklines=true,
    numbers=left,
    numberstyle=\tiny,
    frame=single,
    tabsize=2
}

%======================================================================
\section{Verteilte Ausführung auf dem Rechencluster}
\label{sec:cluster}

\subsection{Einordnung und Aufgabenstellung}
\label{subsec:cluster-einordnung}

Die Strahlverfolgung ist der teuerste Schritt der gesamten Pipeline. Für eine
brauchbare Impulsantwort müssen Hunderttausende bis Millionen einzelner
Schallstrahlen durch die Raumgeometrie verfolgt werden, und jeder dieser Strahlen
verlangt mehrere Schnitt- und Reflexionsberechnungen. Der entscheidende Vorteil
dabei: Kein Strahl hängt vom Ergebnis eines anderen ab. Das Problem ist damit
\emph{embarrassingly parallel} und lässt sich ohne Kommunikation zwischen den
Teilrechnungen auf beliebig viele Rechner verteilen.

Dieses Kapitel dokumentiert die Verteilungsschicht des Systems. Sie startet die
Rust-Simulation gleichzeitig auf allen acht Knoten des Labor-Clusters, wartet auf
deren Abschluss, sammelt die Teilergebnisse ein und fügt sie zu einer einzigen,
statistisch dichteren Impulsantwort zusammen. Aus Sicht der übrigen Pipeline
ändert sich dabei nichts: Am Ende steht wieder eine Datei im gewohnten Schema,
die die Signalverarbeitung unverändert weiterverarbeiten kann.

Zur Abgrenzung von den Nachbarkapiteln gehören zwei Punkte. Erstens ist die
\emph{prozessinterne} Nebenläufigkeit der Rust-Engine, also die Verteilung auf
die Kerne eines einzelnen Rechners mittels \texttt{rayon}, an anderer Stelle
beschrieben und bleibt hier außen vor. Sie parallelisiert innerhalb eines
Prozesses; dieses Kapitel parallelisiert über Prozess- und Rechnergrenzen hinweg.
Zweitens geht es nicht um den einzelnen Subprozess-Aufruf aus der Python-Pipeline,
sondern um die Ebene darüber: die gleichzeitige Ausführung auf mehreren physischen
Knoten und die anschließende Aggregation ihrer Ergebnisse.

\subsection{Theoretischer Hintergrund}
\label{subsec:cluster-theorie}

\subsubsection{Monte-Carlo-Strahlverfolgung}
\label{subsubsec:mc}

Die Engine schätzt die Impulsantwort eines Raumes stochastisch. Sie verschickt
zufällig gerichtete Strahlen von der Schallquelle, verfolgt ihre Reflexionen und
registriert, wann und mit welchem Restdruck sie den Empfänger erreichen. Jeder
einzelne Strahl ist eine Zufallsstichprobe; erst die Summe vieler Strahlen ergibt
ein aussagekräftiges Bild der Raumakustik. Solche Monte-Carlo-Verfahren haben eine
charakteristische Eigenschaft: Ihr Schätzfehler sinkt mit der Wurzel der
Stichprobengröße. Verdoppelt man die Zahl der Strahlen, halbiert sich der Fehler
also nicht, sondern sinkt nur um den Faktor $\sqrt{2}$. Genau diese langsame
Konvergenz macht große Strahlzahlen nötig und motiviert die Verteilung auf mehrere
Rechner.

\subsubsection{Zwei Wege der Verteilung}
\label{subsubsec:strategien}

Ein Monte-Carlo-Verfahren lässt sich auf zwei grundsätzlich verschiedene Arten
verteilen. Tabelle~\ref{tab:strategien} stellt sie gegenüber.

\begin{table}[htbp]
    \centering
    \caption{Zwei Strategien zur Verteilung der Strahlverfolgung.}
    \label{tab:strategien}
    \begin{tabular}{p{3.2cm} p{5.2cm} p{4.2cm}}
        \toprule
        & \textbf{Work-Partitioning} & \textbf{Replikat-Parallelismus} \\
        \midrule
        Prinzip & Gesamtzahl der Strahlen wird in disjunkte Anteile zerlegt & Jeder Knoten rechnet denselben Auftrag mit eigenem Zufallsstrom \\
        Koordination & Zuteilung der Anteile nötig & keine, jeder Knoten ist autark \\
        Zufallszahlen & ein globaler Strom, aufgeteilt & unabhängige Ströme pro Knoten \\
        Robustheit & Ausfall trifft einen festen Anteil & Ausfall kostet nur Stichprobengröße \\
        \bottomrule
    \end{tabular}
\end{table}

Das umgesetzte System verwendet den zweiten Weg. Alle acht Knoten starten die
Simulation mit derselben Konfigurationsdatei \texttt{input\_config.json} und
werfen dennoch jeweils andere Strahlen, weil jeder Prozess seinen eigenen
Zufallsstrom besitzt. Der Ansatz ist bewusst schlicht gewählt: Er braucht keine
zentrale Zuteilung und keinen Austausch zwischen den Knoten, sondern nur, dass am
Ende alle Ergebnisse eingesammelt werden.

\subsubsection{Unabhängigkeit der Zufallsströme}
\label{subsubsec:rng}

Die Korrektheit des Replikat-Ansatzes steht und fällt mit einer Bedingung: Die
Zufallszahlen auf den Knoten müssen unabhängig sein. Würde jeder Knoten dieselbe
Folge zufälliger Richtungen erzeugen, so lieferten alle acht exakt dasselbe
Ergebnis, und das Zusammenführen brächte keinen Erkenntnisgewinn, sondern nur
achtfach dieselben Strahlen. Die Rust-Engine schließt das aus, indem sie ihren
Zufallsstrom pro Thread bzw. Prozess frisch aus der Systementropie zieht:

\begin{lstlisting}[style=clusterlst, language=Rust, caption={Initialisierung des Zufallsstroms pro Prozess (\texttt{simulation.rs}).}]
let mut rng = rand::rng();                       // Omnidirectional, random rays
let random_angle: f32 = rng.random_range(0.0..(2.0 * PI));
let direction = Vec2::new(random_angle.cos(), random_angle.sin());
\end{lstlisting}

Weil jeder der acht Prozesse auf seinem Knoten einen anderen Startzustand erhält,
erzeugt jeder eine andere Stichprobe. Acht Knoten mit je einer Million Strahlen
liefern damit acht Millionen \emph{unabhängige} Stichproben derselben Raumakustik.

\subsubsection{Statistische Konvergenz}
\label{subsubsec:konvergenz}

Statistisch verhält sich der verteilte Lauf wie eine einzige, achtfach größere
Simulation. Für einen Monte-Carlo-Schätzer aus $N$ unabhängigen Teilstichproben
gleicher Größe sinkt die Standardabweichung des Mittelwerts mit
\[
    \sigma_N \;=\; \frac{\sigma_1}{\sqrt{N}} .
\]
Acht Knoten reduzieren die Streuung der geschätzten Impulsantwort demnach um den
Faktor $\sqrt{8} \approx 2{,}83$ gegenüber einem einzelnen Knoten.
Abbildung~\ref{fig:mc-convergence} bestätigt dieses Verhalten in einem numerischen
Experiment: Aufgetragen ist der mittlere quadratische Fehler der geschätzten
Impulsantwort über der Knotenzahl. Die Messpunkte folgen sauber der
$1/\sqrt{N}$-Referenzkurve; von einem auf acht Knoten sinkt der Fehler um den
Faktor $2{,}78$, praktisch exakt der theoretisch erwartete Wert.

\begin{figure}[htbp]
    \centering
    \includegraphics[width=0.8\textwidth]{figures/mc_convergence}
    \caption{Mittlerer quadratischer Fehler der geschätzten Impulsantwort über der
    Knotenzahl $N$ (je $5\cdot10^4$ Strahlen, doppelt logarithmisch). Die
    Messpunkte folgen der $1/\sqrt{N}$-Referenz; acht Knoten senken den Fehler um
    den Faktor $2{,}78 \approx \sqrt{8}$.}
    \label{fig:mc-convergence}
\end{figure}

Bemerkenswert ist die Kehrseite dieser Kurve: Um den Fehler zu halbieren, braucht
es die vierfache Knotenzahl. Der Nutzen zusätzlicher Knoten nimmt also ab, was bei
der Wahl der Clustergröße zu berücksichtigen ist. Für die hier angestrebte
Genauigkeit sind acht Knoten ein sinnvoller Kompromiss zwischen Aufwand und
Qualität.

\subsection{Architektur des Clusters}
\label{subsec:cluster-arch}

Das System folgt einem \emph{Manager-Worker}-Muster
(Abbildung~\ref{fig:cluster-arch}). Ein Manager-Knoten koordiniert den Ablauf,
acht Worker-Knoten führen die eigentliche Simulation aus. Der Manager selbst
rechnet nicht mit, sondern startet die Worker, wartet auf sie und führt am Ende
die Ergebnisse zusammen.

\begin{figure}[htbp]
    \centering
    \includegraphics[width=0.82\textwidth]{figures/cluster_arch}
    \caption{Manager-Worker-Architektur. Der Manager startet die acht Knoten
    parallel per \texttt{ssh}, holt nach Abschluss jede Ergebnisdatei per
    \texttt{scp} zurück und führt sie mit \texttt{merge\_results.py} zusammen.}
    \label{fig:cluster-arch}
\end{figure}

Die Kommunikation läuft vollständig über Standard-Werkzeuge des Betriebssystems.
\texttt{ssh} startet die Prozesse aus der Ferne, \texttt{scp} überträgt die
Ergebnisdateien zurück. Ein eigenes Netzwerkprotokoll oder ein
Message-Passing-Framework kommt nicht zum Einsatz. Für einen Stapelbetrieb, bei
dem jeder Knoten genau einmal startet, eine Weile rechnet und ein einziges
Ergebnis zurückliefert, wäre solcher Aufwand nicht gerechtfertigt: Es gibt keine
laufende Zustandssynchronisation und keinen Nachrichtenaustausch während der
Rechnung, den ein Framework verwalten müsste.

Damit das funktioniert, müssen drei Voraussetzungen erfüllt sein. Auf jedem Knoten
liegt dieselbe vorkompilierte Binärdatei \texttt{rust\_service} samt der
Konfiguration \texttt{input\_config.json}. Der Manager kann sich per SSH ohne
interaktive Passworteingabe anmelden, was eine Schlüssel-basierte
Authentifizierung voraussetzt. Und die Knoten sind über feste Adressen im
Labornetz erreichbar.

\subsection{Ablaufsteuerung: \texttt{run\_cluster.sh}}
\label{subsec:cluster-run}

Das Skript \texttt{run\_cluster.sh} steuert den verteilten Lauf in drei Phasen:
paralleles Starten, Warten und Einsammeln.

\begin{lstlisting}[style=clusterlst, language=bash, caption={Verteilter Start und Einsammeln der Ergebnisse (gekürzt).}]
# Phase 1: alle Knoten gleichzeitig starten (Hintergrundprozesse)
ssh node1@139.6.65.29 "./rust_service --config-file input_config.json" &
ssh node2@139.6.65.18 "./rust_service --config-file input_config.json" &
# ... node3 bis node8 analog ...

wait                       # Phase 2: auf Abschluss ALLER Knoten warten

# Phase 3: Ergebnis jedes Knotens zum Manager kopieren
scp node1@139.6.65.29:~/ir_output.json ~/ir_node1.json
scp node2@139.6.65.18:~/ir_output.json ~/ir_node2.json
# ... node3 bis node8 analog ...
\end{lstlisting}

Zwei Eigenschaften der Shell nutzt der Ablauf gezielt. Das kaufmännische Und
(\texttt{\&}) am Zeilenende schickt jeden \texttt{ssh}-Aufruf in den Hintergrund.
Die acht Simulationen laufen dadurch nahezu gleichzeitig an, statt nacheinander
abgearbeitet zu werden. Das darauffolgende \texttt{wait} blockiert, bis
\emph{alle} Hintergrundprozesse beendet sind, und bildet so eine
Synchronisationsbarriere zwischen der Rechen- und der Einsammel-Phase. Erst wenn
der letzte Knoten fertig ist, beginnt \texttt{scp}, von jedem Knoten dessen
\texttt{ir\_output.json} zu holen und unter einem knotenspezifischen Namen auf dem
Manager abzulegen.

Diese Barriere hat eine wichtige Konsequenz für die Laufzeit: Der Gesamtlauf ist
so langsam wie der \emph{langsamste} Knoten. Solange die Hardware homogen ist,
fällt das kaum ins Gewicht; auf ungleichen Knoten würde ein einzelner langsamer
Rechner den ganzen Lauf ausbremsen, während die übrigen sieben bereits fertig
warten. Dieses Straggler-Problem ist der zentrale Grund, warum ein statischer
Replikat-Ansatz mit gleich großen Aufträgen homogene Hardware voraussetzt.

\subsection{Aggregation der Teilergebnisse: \texttt{merge\_results.py}}
\label{subsec:cluster-merge}

Sobald alle Teilergebnisse auf dem Manager liegen, führt \texttt{merge\_results.py}
sie zu einer einzigen Impulsantwort zusammen. Jeder Knoten liefert dasselbe
Ausgabeschema, nämlich Metadaten sowie zwei positionsgleiche Listen aus
Verzögerungen und Drücken. Die Aggregation besteht deshalb im Kern aus einer
schlichten Verkettung dieser Listen:

\begin{lstlisting}[style=clusterlst, language=Python, caption={Zusammenführen der Knoten-Ergebnisse, gekürzt (\texttt{merge\_results.py}).}]
all_delays, all_pressures = [], []
for f in sorted(glob.glob("/home/manager/ir_node*.json")):
    with open(f) as file:
        data = json.load(file)
        all_delays    += data["hits"]["delays_seconds"]
        all_pressures += data["hits"]["pressures"]

merged = {
    "metadata": {"sample_rate": 44100, "rays_cast": 1000000 * 8,
                 "rays_received": len(all_delays), "room_name": "MVP_Test_Room_1"},
    "hits": {"delays_seconds": all_delays, "pressures": all_pressures},
}
json.dump(merged, open("/home/manager/ir_merged.json", "w"), indent=2)
\end{lstlisting}

Dass die Verkettung hier die richtige Operation ist, folgt direkt aus dem
Replikat-Modell. Weil die einzelnen Läufe unabhängige Stichproben desselben Raums
sind, ist die Vereinigung aller Treffer statistisch identisch mit dem Ergebnis
eines einzigen Laufs mit acht Millionen Strahlen. Es findet keine Mittelung und
keine Gewichtung statt, sondern nur ein Aneinanderhängen. Das Feld
\texttt{rays\_cast} ist im Skript fest auf $8 \cdot 10^{6}$ gesetzt; solange der in
Abschnitt~\ref{subsec:cluster-grenzen} beschriebene Kopierfehler bei Knoten~6
besteht, fließen tatsächlich nur sieben Knoten in die Aggregation ein, und dieser
Metadatenwert überschätzt die reale Strahlzahl. \texttt{rays\_received} zählt
dagegen die tatsächlich empfangenen Treffer. Die
zusammengeführte Datei trägt exakt dasselbe Schema wie die Ausgabe eines einzelnen
Knotens und ist damit ohne weitere Anpassung für die DSP-Pipeline verwendbar.

\subsection{Datenfluss im Überblick}
\label{subsec:cluster-datenfluss}

Der gesamte verteilte Ablauf lässt sich in fünf Stufen zusammenfassen:

\begin{enumerate}
    \item Der Manager startet per SSH auf jedem der acht Knoten
    \texttt{rust\_service} mit identischer Konfiguration, alle parallel.
    \item Jeder Knoten simuliert unabhängig eine Million Strahlen und schreibt
    lokal seine \texttt{ir\_output.json}.
    \item \texttt{wait} synchronisiert: Die Einsammel-Phase beginnt erst, wenn der
    letzte Knoten fertig ist.
    \item \texttt{scp} überträgt jede \texttt{ir\_output.json} als
    \texttt{ir\_node$k$.json} zum Manager.
    \item \texttt{merge\_results.py} verkettet alle Trefferlisten zu einer
    \texttt{ir\_merged.json}, die anschließend die DSP-Pipeline verarbeitet.
\end{enumerate}

\subsection{Leistungsbetrachtung}
\label{subsec:cluster-leistung}

\subsubsection{Erwarteter Speedup}
\label{subsubsec:speedup}

Bei der Leistung sind zwei Sichtweisen zu trennen. In der umgesetzten
Replikat-Betriebsart rechnet jeder Knoten dieselbe Strahlzahl, sodass acht Knoten
in nahezu derselben Wanduhrzeit die achtfache Stichprobe liefern. Der Gewinn ist
hier keine kürzere Laufzeit, sondern die höhere Genauigkeit aus
Abbildung~\ref{fig:mc-convergence}; man spricht vom schwachen Skalieren im Sinne
von Gustafson, bei dem mit der Rechnerzahl die bearbeitete Problemgröße wächst,
nicht die Geschwindigkeit einer festen Aufgabe. Will man dagegen eine feste
Gesamtstrahlzahl in kürzerer Zeit erreichen, so verteilt man diese auf die Knoten,
sodass jeder nur $8\cdot10^{6}/k$ Strahlen rechnet. Erst in dieser
Partitionierungs-Betriebsart greift das Gesetz von Amdahl. Bezeichnet $p$ den parallelen
Anteil der Arbeit, hier die reine Strahlverfolgung, und $1-p$ den seriellen Rest
aus Prozessstart, Ergebnistransfer und Aggregation, so beträgt der theoretische
Speedup
\[
    S(k) \;=\; \frac{1}{(1-p) + \dfrac{p}{k}} .
\]
Da die Strahlverfolgung die Laufzeit klar dominiert, liegt $p$ nahe bei eins, und
$S(k)$ steigt bis $k=8$ nahezu linear. In der Praxis bleibt der reale Speedup
darunter, und zwar aus drei Gründen: dem bereits erwähnten Straggler-Effekt (der
Lauf ist so schnell wie der langsamste Knoten), dem Overhead für den
SSH-Verbindungsaufbau sowie der Übertragung und Zusammenführung der
Ergebnisdateien.

\subsubsection{Kommunikations- und I/O-Overhead}
\label{subsubsec:overhead}

Der Overhead dieses Ansatzes ist klein, aber nicht null. Pro Knoten fallen ein
SSH-Verbindungsaufbau und eine \texttt{scp}-Übertragung an. Übertragen wird jeweils
eine JSON-Datei, deren Größe mit der Zahl der empfangenen Treffer wächst; bei
einer Million geworfener Strahlen liegt sie im Bereich einiger Megabyte. Da diese
acht Übertragungen erst \emph{nach} der Rechenphase stattfinden und nicht mit ihr
überlappen, addieren sie sich voll zur Gesamtlaufzeit. Bei deutlich mehr Knoten
oder größeren Ergebnisdateien würde dieser sequenzielle Einsammel-Schritt zum
Flaschenhals; eine Überlappung von Rechnung und Transfer oder ein kompakteres
Ausgabeformat als das gut lesbare, aber platzintensive JSON wären dann die
naheliegenden Gegenmaßnahmen.

\subsubsection{Messung am Cluster}
\label{subsubsec:messung}

Die konkreten Laufzeiten hängen von der Laborhardware ab und sind am Cluster zu
messen. Um einen aussagekräftigen Speedup zu erhalten, wird in der
Partitionierungs-Betriebsart gemessen: In \texttt{input\_config.json} wird
\texttt{rays\_to\_cast} auf $8\cdot10^{6}/k$ gesetzt, sodass die Gesamtarbeit über
alle $k$ Knoten konstant bleibt. Dann wird \texttt{run\_cluster.sh} für
verschiedene Knotenzahlen mit der Shell-eigenen Zeitmessung ausgeführt:

\begin{lstlisting}[style=clusterlst, language=bash, caption={Messvorschrift für die Laufzeit am Cluster.}]
time ./run_cluster.sh        # Wanduhrzeit fuer den kompletten verteilten Lauf
\end{lstlisting}

Tabelle~\ref{tab:speedup} nimmt die Messwerte auf. Als Referenz dient der Lauf mit
einem einzelnen Knoten; der Speedup ergibt sich als Verhältnis der Wanduhrzeiten.

\begin{table}[htbp]
    \centering
    \caption{Wanduhrzeit und Speedup des verteilten Laufs (am Cluster zu messen).}
    \label{tab:speedup}
    \begin{tabular}{r r r}
        \toprule
        Knoten $k$ & Wanduhrzeit [s] & Speedup $S$ \\
        \midrule
        1 & \dots & $1{,}00$ \\
        2 & \dots & \dots \\
        4 & \dots & \dots \\
        8 & \dots & \dots \\
        \bottomrule
    \end{tabular}
\end{table}

\subsection{Fehlertoleranz und Sicherheit}
\label{subsec:cluster-fehler}

Zwei Querschnittsthemen verdienen eigene Betrachtung. Beim \textbf{Sicherheits}aspekt
ist die Schlüssel-basierte SSH-Authentifizierung zentral: Der Manager benötigt
passwortlosen Zugriff auf jeden Knoten, was über hinterlegte öffentliche Schlüssel
gelöst ist. Der private Schlüssel des Managers ist damit das entscheidende
Geheimnis; wer ihn besitzt, kann auf allen Knoten Code ausführen. In einem
abgeschlossenen Labornetz ist das vertretbar, für einen offeneren Betrieb wäre eine
engere Rechtevergabe angebracht.

Beim Thema \textbf{Fehlertoleranz} zeigt der aktuelle Stand seine deutlichste
Schwäche, die im folgenden Abschnitt zusammen mit den übrigen Grenzen diskutiert
wird. Der Kern des Problems: Ein ausgefallener Knoten wird nicht erkannt, sondern
schlicht übergangen, sodass die Stichprobengröße unbemerkt sinkt.

\subsection{Grenzen des aktuellen Standes und Ausblick}
\label{subsec:cluster-grenzen}

Die Verteilungsschicht ist als funktionsfähiger Prototyp umgesetzt. Für einen
robusten Betrieb bestehen mehrere klar benannte Schwachstellen, die zugleich die
nächsten Entwicklungsschritte vorgeben.

\begin{itemize}
    \item \textbf{Fest verdrahtete Konfiguration.} IP-Adressen, Knotennamen, die
    Knotenzahl und die Ablagepfade stehen direkt im Skript. Ein Umzug auf andere
    Hardware oder eine andere Knotenzahl erfordert Änderungen am Quelltext. Eine
    Auslagerung dieser Werte in eine Konfigurationsdatei wäre der naheliegende
    erste Schritt.

    \item \textbf{Kopierfehler bei Knoten~6.} Im Einsammel-Schritt wird das
    Ergebnis von Knoten~6 versehentlich nach \texttt{ir\_node5.json} kopiert und
    überschreibt damit das Ergebnis von Knoten~5. Effektiv fließen so nur sieben
    statt acht unabhängige Stichproben in die Aggregation ein. Der Fehler ist ein
    Tippfehler im Zielnamen und leicht zu beheben, führt aber vor Augen, wie
    fehleranfällig die von Hand ausgeschriebene Knotenliste ist -- ein weiteres
    Argument für die genannte Auslagerung in eine Schleife über eine
    Konfigurationsliste.

    \item \textbf{Keine Fehlererkennung.} Fällt ein Knoten aus oder liefert er kein
    Ergebnis, schlägt der zugehörige \texttt{scp}-Aufruf fehl, der Gesamtlauf setzt
    sich aber fort und aggregiert stillschweigend über die vorhandenen Dateien. Ein
    Knotenausfall bleibt damit unbemerkt und verfälscht die Stichprobengröße. Eine
    Prüfung der Rückgabecodes und eine Warnung bei fehlenden Ergebnissen wären hier
    angebracht.

    \item \textbf{Manuelle Kopplung an die DSP-Pipeline.} Die Aggregation erzeugt
    \texttt{ir\_merged.json}, die Python-DSP-Pipeline liest jedoch
    \texttt{ir\_output.json}. Der Übergang zwischen verteiltem Lauf und Faltung ist
    derzeit ein manueller Schritt. Eine saubere Integration würde den Cluster-Lauf
    direkt aus der Orchestrierung heraus anstoßen und die zusammengeführte Datei
    automatisch weiterreichen.

    \item \textbf{Statische Verteilung.} Der Replikat-Ansatz mit gleich großen
    Aufträgen ist einfach und robust, setzt aber homogene Hardware voraus und lastet
    unterschiedlich schnelle Knoten nicht unterschiedlich aus. Ein Modell mit
    dynamischer Arbeitsverteilung -- ein Manager vergibt Ray-Kontingente, sobald ein
    Knoten frei wird -- würde heterogene Hardware besser nutzen und die Knotenzahl
    flexibel machen.
\end{itemize}

Alle genannten Punkte betreffen Robustheit und Bedienbarkeit, nicht die
grundsätzliche Korrektheit des Verfahrens. Die verteilte Ausführung liefert eine
gültige, statistisch dichtere Impulsantwort, und das Manager-Worker-Muster über
SSH und \texttt{scp} erfüllt seinen Zweck im Laborbetrieb.

\subsection{Zusammenfassung}
\label{subsec:cluster-zusammenfassung}

Dieses Kapitel hat die verteilte Ausführung der Strahlverfolgung auf dem
Labor-Cluster beschrieben. Weil die Simulation embarrassingly parallel ist, genügt
ein Replikat-Ansatz: Alle acht Knoten rechnen denselben Auftrag mit unabhängigen
Zufallsströmen, was statistisch einer achtfach größeren Simulation entspricht und
den Schätzfehler der Impulsantwort um rund $\sqrt{8}$ senkt -- ein Verhalten, das
das numerische Experiment in Abbildung~\ref{fig:mc-convergence} bestätigt. Ein
Shell-Skript startet die Knoten per SSH parallel, synchronisiert über \texttt{wait}
und sammelt die Ergebnisse per \texttt{scp} ein; \texttt{merge\_results.py}
verkettet die Trefferlisten zu einer gemeinsamen Datei im gewohnten Schema. Der
Prototyp erfüllt seine Aufgabe, weist aber in Konfigurierbarkeit, Fehlererkennung
und der Anbindung an die DSP-Pipeline die diskutierten Schwächen auf, die den
Rahmen für die Weiterentwicklung abstecken.
 im Master.
% Benoetigte Pakete: amsmath, listings, graphicx, booktabs.
% Fuer saubere Optik im Master: fontenc, babel[ngerman], microtype.
% Abbildungen: latexDocs/figures/cluster_arch.png, mc_convergence.png
%
% STANDALONE-VORSCHAU:
%   \documentclass[12pt]{article}
%   \usepackage[utf8]{inputenc}\usepackage[T1]{fontenc}
%   \usepackage[ngerman]{babel}\usepackage{microtype}
%   \usepackage{amsmath}\usepackage{listings}
%   \usepackage{graphicx}\usepackage{booktabs}
%   \begin{document}
%   % =====================================================================
% Kapitel: Verteilte Ausfuehrung auf dem Rechencluster
% Autor: Mier
%
% EINBINDEN: \input{Kapitel_Cluster_Mier} im Master.
% Benoetigte Pakete: amsmath, listings, graphicx, booktabs.
% Fuer saubere Optik im Master: fontenc, babel[ngerman], microtype.
% Abbildungen: latexDocs/figures/cluster_arch.png, mc_convergence.png
%
% STANDALONE-VORSCHAU:
%   \documentclass[12pt]{article}
%   \usepackage[utf8]{inputenc}\usepackage[T1]{fontenc}
%   \usepackage[ngerman]{babel}\usepackage{microtype}
%   \usepackage{amsmath}\usepackage{listings}
%   \usepackage{graphicx}\usepackage{booktabs}
%   \begin{document}
%   \input{Kapitel_Cluster_Mier}
%   \end{document}
% =====================================================================

\lstdefinestyle{clusterlst}{
    basicstyle=\ttfamily\small,
    showstringspaces=false,
    breaklines=true,
    numbers=left,
    numberstyle=\tiny,
    frame=single,
    tabsize=2
}

%======================================================================
\section{Verteilte Ausführung auf dem Rechencluster}
\label{sec:cluster}

\subsection{Einordnung und Aufgabenstellung}
\label{subsec:cluster-einordnung}

Die Strahlverfolgung ist der teuerste Schritt der gesamten Pipeline. Für eine
brauchbare Impulsantwort müssen Hunderttausende bis Millionen einzelner
Schallstrahlen durch die Raumgeometrie verfolgt werden, und jeder dieser Strahlen
verlangt mehrere Schnitt- und Reflexionsberechnungen. Der entscheidende Vorteil
dabei: Kein Strahl hängt vom Ergebnis eines anderen ab. Das Problem ist damit
\emph{embarrassingly parallel} und lässt sich ohne Kommunikation zwischen den
Teilrechnungen auf beliebig viele Rechner verteilen.

Dieses Kapitel dokumentiert die Verteilungsschicht des Systems. Sie startet die
Rust-Simulation gleichzeitig auf allen acht Knoten des Labor-Clusters, wartet auf
deren Abschluss, sammelt die Teilergebnisse ein und fügt sie zu einer einzigen,
statistisch dichteren Impulsantwort zusammen. Aus Sicht der übrigen Pipeline
ändert sich dabei nichts: Am Ende steht wieder eine Datei im gewohnten Schema,
die die Signalverarbeitung unverändert weiterverarbeiten kann.

Zur Abgrenzung von den Nachbarkapiteln gehören zwei Punkte. Erstens ist die
\emph{prozessinterne} Nebenläufigkeit der Rust-Engine, also die Verteilung auf
die Kerne eines einzelnen Rechners mittels \texttt{rayon}, an anderer Stelle
beschrieben und bleibt hier außen vor. Sie parallelisiert innerhalb eines
Prozesses; dieses Kapitel parallelisiert über Prozess- und Rechnergrenzen hinweg.
Zweitens geht es nicht um den einzelnen Subprozess-Aufruf aus der Python-Pipeline,
sondern um die Ebene darüber: die gleichzeitige Ausführung auf mehreren physischen
Knoten und die anschließende Aggregation ihrer Ergebnisse.

\subsection{Theoretischer Hintergrund}
\label{subsec:cluster-theorie}

\subsubsection{Monte-Carlo-Strahlverfolgung}
\label{subsubsec:mc}

Die Engine schätzt die Impulsantwort eines Raumes stochastisch. Sie verschickt
zufällig gerichtete Strahlen von der Schallquelle, verfolgt ihre Reflexionen und
registriert, wann und mit welchem Restdruck sie den Empfänger erreichen. Jeder
einzelne Strahl ist eine Zufallsstichprobe; erst die Summe vieler Strahlen ergibt
ein aussagekräftiges Bild der Raumakustik. Solche Monte-Carlo-Verfahren haben eine
charakteristische Eigenschaft: Ihr Schätzfehler sinkt mit der Wurzel der
Stichprobengröße. Verdoppelt man die Zahl der Strahlen, halbiert sich der Fehler
also nicht, sondern sinkt nur um den Faktor $\sqrt{2}$. Genau diese langsame
Konvergenz macht große Strahlzahlen nötig und motiviert die Verteilung auf mehrere
Rechner.

\subsubsection{Zwei Wege der Verteilung}
\label{subsubsec:strategien}

Ein Monte-Carlo-Verfahren lässt sich auf zwei grundsätzlich verschiedene Arten
verteilen. Tabelle~\ref{tab:strategien} stellt sie gegenüber.

\begin{table}[htbp]
    \centering
    \caption{Zwei Strategien zur Verteilung der Strahlverfolgung.}
    \label{tab:strategien}
    \begin{tabular}{p{3.2cm} p{5.2cm} p{4.2cm}}
        \toprule
        & \textbf{Work-Partitioning} & \textbf{Replikat-Parallelismus} \\
        \midrule
        Prinzip & Gesamtzahl der Strahlen wird in disjunkte Anteile zerlegt & Jeder Knoten rechnet denselben Auftrag mit eigenem Zufallsstrom \\
        Koordination & Zuteilung der Anteile nötig & keine, jeder Knoten ist autark \\
        Zufallszahlen & ein globaler Strom, aufgeteilt & unabhängige Ströme pro Knoten \\
        Robustheit & Ausfall trifft einen festen Anteil & Ausfall kostet nur Stichprobengröße \\
        \bottomrule
    \end{tabular}
\end{table}

Das umgesetzte System verwendet den zweiten Weg. Alle acht Knoten starten die
Simulation mit derselben Konfigurationsdatei \texttt{input\_config.json} und
werfen dennoch jeweils andere Strahlen, weil jeder Prozess seinen eigenen
Zufallsstrom besitzt. Der Ansatz ist bewusst schlicht gewählt: Er braucht keine
zentrale Zuteilung und keinen Austausch zwischen den Knoten, sondern nur, dass am
Ende alle Ergebnisse eingesammelt werden.

\subsubsection{Unabhängigkeit der Zufallsströme}
\label{subsubsec:rng}

Die Korrektheit des Replikat-Ansatzes steht und fällt mit einer Bedingung: Die
Zufallszahlen auf den Knoten müssen unabhängig sein. Würde jeder Knoten dieselbe
Folge zufälliger Richtungen erzeugen, so lieferten alle acht exakt dasselbe
Ergebnis, und das Zusammenführen brächte keinen Erkenntnisgewinn, sondern nur
achtfach dieselben Strahlen. Die Rust-Engine schließt das aus, indem sie ihren
Zufallsstrom pro Thread bzw. Prozess frisch aus der Systementropie zieht:

\begin{lstlisting}[style=clusterlst, language=Rust, caption={Initialisierung des Zufallsstroms pro Prozess (\texttt{simulation.rs}).}]
let mut rng = rand::rng();                       // Omnidirectional, random rays
let random_angle: f32 = rng.random_range(0.0..(2.0 * PI));
let direction = Vec2::new(random_angle.cos(), random_angle.sin());
\end{lstlisting}

Weil jeder der acht Prozesse auf seinem Knoten einen anderen Startzustand erhält,
erzeugt jeder eine andere Stichprobe. Acht Knoten mit je einer Million Strahlen
liefern damit acht Millionen \emph{unabhängige} Stichproben derselben Raumakustik.

\subsubsection{Statistische Konvergenz}
\label{subsubsec:konvergenz}

Statistisch verhält sich der verteilte Lauf wie eine einzige, achtfach größere
Simulation. Für einen Monte-Carlo-Schätzer aus $N$ unabhängigen Teilstichproben
gleicher Größe sinkt die Standardabweichung des Mittelwerts mit
\[
    \sigma_N \;=\; \frac{\sigma_1}{\sqrt{N}} .
\]
Acht Knoten reduzieren die Streuung der geschätzten Impulsantwort demnach um den
Faktor $\sqrt{8} \approx 2{,}83$ gegenüber einem einzelnen Knoten.
Abbildung~\ref{fig:mc-convergence} bestätigt dieses Verhalten in einem numerischen
Experiment: Aufgetragen ist der mittlere quadratische Fehler der geschätzten
Impulsantwort über der Knotenzahl. Die Messpunkte folgen sauber der
$1/\sqrt{N}$-Referenzkurve; von einem auf acht Knoten sinkt der Fehler um den
Faktor $2{,}78$, praktisch exakt der theoretisch erwartete Wert.

\begin{figure}[htbp]
    \centering
    \includegraphics[width=0.8\textwidth]{figures/mc_convergence}
    \caption{Mittlerer quadratischer Fehler der geschätzten Impulsantwort über der
    Knotenzahl $N$ (je $5\cdot10^4$ Strahlen, doppelt logarithmisch). Die
    Messpunkte folgen der $1/\sqrt{N}$-Referenz; acht Knoten senken den Fehler um
    den Faktor $2{,}78 \approx \sqrt{8}$.}
    \label{fig:mc-convergence}
\end{figure}

Bemerkenswert ist die Kehrseite dieser Kurve: Um den Fehler zu halbieren, braucht
es die vierfache Knotenzahl. Der Nutzen zusätzlicher Knoten nimmt also ab, was bei
der Wahl der Clustergröße zu berücksichtigen ist. Für die hier angestrebte
Genauigkeit sind acht Knoten ein sinnvoller Kompromiss zwischen Aufwand und
Qualität.

\subsection{Architektur des Clusters}
\label{subsec:cluster-arch}

Das System folgt einem \emph{Manager-Worker}-Muster
(Abbildung~\ref{fig:cluster-arch}). Ein Manager-Knoten koordiniert den Ablauf,
acht Worker-Knoten führen die eigentliche Simulation aus. Der Manager selbst
rechnet nicht mit, sondern startet die Worker, wartet auf sie und führt am Ende
die Ergebnisse zusammen.

\begin{figure}[htbp]
    \centering
    \includegraphics[width=0.82\textwidth]{figures/cluster_arch}
    \caption{Manager-Worker-Architektur. Der Manager startet die acht Knoten
    parallel per \texttt{ssh}, holt nach Abschluss jede Ergebnisdatei per
    \texttt{scp} zurück und führt sie mit \texttt{merge\_results.py} zusammen.}
    \label{fig:cluster-arch}
\end{figure}

Die Kommunikation läuft vollständig über Standard-Werkzeuge des Betriebssystems.
\texttt{ssh} startet die Prozesse aus der Ferne, \texttt{scp} überträgt die
Ergebnisdateien zurück. Ein eigenes Netzwerkprotokoll oder ein
Message-Passing-Framework kommt nicht zum Einsatz. Für einen Stapelbetrieb, bei
dem jeder Knoten genau einmal startet, eine Weile rechnet und ein einziges
Ergebnis zurückliefert, wäre solcher Aufwand nicht gerechtfertigt: Es gibt keine
laufende Zustandssynchronisation und keinen Nachrichtenaustausch während der
Rechnung, den ein Framework verwalten müsste.

Damit das funktioniert, müssen drei Voraussetzungen erfüllt sein. Auf jedem Knoten
liegt dieselbe vorkompilierte Binärdatei \texttt{rust\_service} samt der
Konfiguration \texttt{input\_config.json}. Der Manager kann sich per SSH ohne
interaktive Passworteingabe anmelden, was eine Schlüssel-basierte
Authentifizierung voraussetzt. Und die Knoten sind über feste Adressen im
Labornetz erreichbar.

\subsection{Ablaufsteuerung: \texttt{run\_cluster.sh}}
\label{subsec:cluster-run}

Das Skript \texttt{run\_cluster.sh} steuert den verteilten Lauf in drei Phasen:
paralleles Starten, Warten und Einsammeln.

\begin{lstlisting}[style=clusterlst, language=bash, caption={Verteilter Start und Einsammeln der Ergebnisse (gekürzt).}]
# Phase 1: alle Knoten gleichzeitig starten (Hintergrundprozesse)
ssh node1@139.6.65.29 "./rust_service --config-file input_config.json" &
ssh node2@139.6.65.18 "./rust_service --config-file input_config.json" &
# ... node3 bis node8 analog ...

wait                       # Phase 2: auf Abschluss ALLER Knoten warten

# Phase 3: Ergebnis jedes Knotens zum Manager kopieren
scp node1@139.6.65.29:~/ir_output.json ~/ir_node1.json
scp node2@139.6.65.18:~/ir_output.json ~/ir_node2.json
# ... node3 bis node8 analog ...
\end{lstlisting}

Zwei Eigenschaften der Shell nutzt der Ablauf gezielt. Das kaufmännische Und
(\texttt{\&}) am Zeilenende schickt jeden \texttt{ssh}-Aufruf in den Hintergrund.
Die acht Simulationen laufen dadurch nahezu gleichzeitig an, statt nacheinander
abgearbeitet zu werden. Das darauffolgende \texttt{wait} blockiert, bis
\emph{alle} Hintergrundprozesse beendet sind, und bildet so eine
Synchronisationsbarriere zwischen der Rechen- und der Einsammel-Phase. Erst wenn
der letzte Knoten fertig ist, beginnt \texttt{scp}, von jedem Knoten dessen
\texttt{ir\_output.json} zu holen und unter einem knotenspezifischen Namen auf dem
Manager abzulegen.

Diese Barriere hat eine wichtige Konsequenz für die Laufzeit: Der Gesamtlauf ist
so langsam wie der \emph{langsamste} Knoten. Solange die Hardware homogen ist,
fällt das kaum ins Gewicht; auf ungleichen Knoten würde ein einzelner langsamer
Rechner den ganzen Lauf ausbremsen, während die übrigen sieben bereits fertig
warten. Dieses Straggler-Problem ist der zentrale Grund, warum ein statischer
Replikat-Ansatz mit gleich großen Aufträgen homogene Hardware voraussetzt.

\subsection{Aggregation der Teilergebnisse: \texttt{merge\_results.py}}
\label{subsec:cluster-merge}

Sobald alle Teilergebnisse auf dem Manager liegen, führt \texttt{merge\_results.py}
sie zu einer einzigen Impulsantwort zusammen. Jeder Knoten liefert dasselbe
Ausgabeschema, nämlich Metadaten sowie zwei positionsgleiche Listen aus
Verzögerungen und Drücken. Die Aggregation besteht deshalb im Kern aus einer
schlichten Verkettung dieser Listen:

\begin{lstlisting}[style=clusterlst, language=Python, caption={Zusammenführen der Knoten-Ergebnisse, gekürzt (\texttt{merge\_results.py}).}]
all_delays, all_pressures = [], []
for f in sorted(glob.glob("/home/manager/ir_node*.json")):
    with open(f) as file:
        data = json.load(file)
        all_delays    += data["hits"]["delays_seconds"]
        all_pressures += data["hits"]["pressures"]

merged = {
    "metadata": {"sample_rate": 44100, "rays_cast": 1000000 * 8,
                 "rays_received": len(all_delays), "room_name": "MVP_Test_Room_1"},
    "hits": {"delays_seconds": all_delays, "pressures": all_pressures},
}
json.dump(merged, open("/home/manager/ir_merged.json", "w"), indent=2)
\end{lstlisting}

Dass die Verkettung hier die richtige Operation ist, folgt direkt aus dem
Replikat-Modell. Weil die einzelnen Läufe unabhängige Stichproben desselben Raums
sind, ist die Vereinigung aller Treffer statistisch identisch mit dem Ergebnis
eines einzigen Laufs mit acht Millionen Strahlen. Es findet keine Mittelung und
keine Gewichtung statt, sondern nur ein Aneinanderhängen. Das Feld
\texttt{rays\_cast} ist im Skript fest auf $8 \cdot 10^{6}$ gesetzt; solange der in
Abschnitt~\ref{subsec:cluster-grenzen} beschriebene Kopierfehler bei Knoten~6
besteht, fließen tatsächlich nur sieben Knoten in die Aggregation ein, und dieser
Metadatenwert überschätzt die reale Strahlzahl. \texttt{rays\_received} zählt
dagegen die tatsächlich empfangenen Treffer. Die
zusammengeführte Datei trägt exakt dasselbe Schema wie die Ausgabe eines einzelnen
Knotens und ist damit ohne weitere Anpassung für die DSP-Pipeline verwendbar.

\subsection{Datenfluss im Überblick}
\label{subsec:cluster-datenfluss}

Der gesamte verteilte Ablauf lässt sich in fünf Stufen zusammenfassen:

\begin{enumerate}
    \item Der Manager startet per SSH auf jedem der acht Knoten
    \texttt{rust\_service} mit identischer Konfiguration, alle parallel.
    \item Jeder Knoten simuliert unabhängig eine Million Strahlen und schreibt
    lokal seine \texttt{ir\_output.json}.
    \item \texttt{wait} synchronisiert: Die Einsammel-Phase beginnt erst, wenn der
    letzte Knoten fertig ist.
    \item \texttt{scp} überträgt jede \texttt{ir\_output.json} als
    \texttt{ir\_node$k$.json} zum Manager.
    \item \texttt{merge\_results.py} verkettet alle Trefferlisten zu einer
    \texttt{ir\_merged.json}, die anschließend die DSP-Pipeline verarbeitet.
\end{enumerate}

\subsection{Leistungsbetrachtung}
\label{subsec:cluster-leistung}

\subsubsection{Erwarteter Speedup}
\label{subsubsec:speedup}

Bei der Leistung sind zwei Sichtweisen zu trennen. In der umgesetzten
Replikat-Betriebsart rechnet jeder Knoten dieselbe Strahlzahl, sodass acht Knoten
in nahezu derselben Wanduhrzeit die achtfache Stichprobe liefern. Der Gewinn ist
hier keine kürzere Laufzeit, sondern die höhere Genauigkeit aus
Abbildung~\ref{fig:mc-convergence}; man spricht vom schwachen Skalieren im Sinne
von Gustafson, bei dem mit der Rechnerzahl die bearbeitete Problemgröße wächst,
nicht die Geschwindigkeit einer festen Aufgabe. Will man dagegen eine feste
Gesamtstrahlzahl in kürzerer Zeit erreichen, so verteilt man diese auf die Knoten,
sodass jeder nur $8\cdot10^{6}/k$ Strahlen rechnet. Erst in dieser
Partitionierungs-Betriebsart greift das Gesetz von Amdahl. Bezeichnet $p$ den parallelen
Anteil der Arbeit, hier die reine Strahlverfolgung, und $1-p$ den seriellen Rest
aus Prozessstart, Ergebnistransfer und Aggregation, so beträgt der theoretische
Speedup
\[
    S(k) \;=\; \frac{1}{(1-p) + \dfrac{p}{k}} .
\]
Da die Strahlverfolgung die Laufzeit klar dominiert, liegt $p$ nahe bei eins, und
$S(k)$ steigt bis $k=8$ nahezu linear. In der Praxis bleibt der reale Speedup
darunter, und zwar aus drei Gründen: dem bereits erwähnten Straggler-Effekt (der
Lauf ist so schnell wie der langsamste Knoten), dem Overhead für den
SSH-Verbindungsaufbau sowie der Übertragung und Zusammenführung der
Ergebnisdateien.

\subsubsection{Kommunikations- und I/O-Overhead}
\label{subsubsec:overhead}

Der Overhead dieses Ansatzes ist klein, aber nicht null. Pro Knoten fallen ein
SSH-Verbindungsaufbau und eine \texttt{scp}-Übertragung an. Übertragen wird jeweils
eine JSON-Datei, deren Größe mit der Zahl der empfangenen Treffer wächst; bei
einer Million geworfener Strahlen liegt sie im Bereich einiger Megabyte. Da diese
acht Übertragungen erst \emph{nach} der Rechenphase stattfinden und nicht mit ihr
überlappen, addieren sie sich voll zur Gesamtlaufzeit. Bei deutlich mehr Knoten
oder größeren Ergebnisdateien würde dieser sequenzielle Einsammel-Schritt zum
Flaschenhals; eine Überlappung von Rechnung und Transfer oder ein kompakteres
Ausgabeformat als das gut lesbare, aber platzintensive JSON wären dann die
naheliegenden Gegenmaßnahmen.

\subsubsection{Messung am Cluster}
\label{subsubsec:messung}

Die konkreten Laufzeiten hängen von der Laborhardware ab und sind am Cluster zu
messen. Um einen aussagekräftigen Speedup zu erhalten, wird in der
Partitionierungs-Betriebsart gemessen: In \texttt{input\_config.json} wird
\texttt{rays\_to\_cast} auf $8\cdot10^{6}/k$ gesetzt, sodass die Gesamtarbeit über
alle $k$ Knoten konstant bleibt. Dann wird \texttt{run\_cluster.sh} für
verschiedene Knotenzahlen mit der Shell-eigenen Zeitmessung ausgeführt:

\begin{lstlisting}[style=clusterlst, language=bash, caption={Messvorschrift für die Laufzeit am Cluster.}]
time ./run_cluster.sh        # Wanduhrzeit fuer den kompletten verteilten Lauf
\end{lstlisting}

Tabelle~\ref{tab:speedup} nimmt die Messwerte auf. Als Referenz dient der Lauf mit
einem einzelnen Knoten; der Speedup ergibt sich als Verhältnis der Wanduhrzeiten.

\begin{table}[htbp]
    \centering
    \caption{Wanduhrzeit und Speedup des verteilten Laufs (am Cluster zu messen).}
    \label{tab:speedup}
    \begin{tabular}{r r r}
        \toprule
        Knoten $k$ & Wanduhrzeit [s] & Speedup $S$ \\
        \midrule
        1 & \dots & $1{,}00$ \\
        2 & \dots & \dots \\
        4 & \dots & \dots \\
        8 & \dots & \dots \\
        \bottomrule
    \end{tabular}
\end{table}

\subsection{Fehlertoleranz und Sicherheit}
\label{subsec:cluster-fehler}

Zwei Querschnittsthemen verdienen eigene Betrachtung. Beim \textbf{Sicherheits}aspekt
ist die Schlüssel-basierte SSH-Authentifizierung zentral: Der Manager benötigt
passwortlosen Zugriff auf jeden Knoten, was über hinterlegte öffentliche Schlüssel
gelöst ist. Der private Schlüssel des Managers ist damit das entscheidende
Geheimnis; wer ihn besitzt, kann auf allen Knoten Code ausführen. In einem
abgeschlossenen Labornetz ist das vertretbar, für einen offeneren Betrieb wäre eine
engere Rechtevergabe angebracht.

Beim Thema \textbf{Fehlertoleranz} zeigt der aktuelle Stand seine deutlichste
Schwäche, die im folgenden Abschnitt zusammen mit den übrigen Grenzen diskutiert
wird. Der Kern des Problems: Ein ausgefallener Knoten wird nicht erkannt, sondern
schlicht übergangen, sodass die Stichprobengröße unbemerkt sinkt.

\subsection{Grenzen des aktuellen Standes und Ausblick}
\label{subsec:cluster-grenzen}

Die Verteilungsschicht ist als funktionsfähiger Prototyp umgesetzt. Für einen
robusten Betrieb bestehen mehrere klar benannte Schwachstellen, die zugleich die
nächsten Entwicklungsschritte vorgeben.

\begin{itemize}
    \item \textbf{Fest verdrahtete Konfiguration.} IP-Adressen, Knotennamen, die
    Knotenzahl und die Ablagepfade stehen direkt im Skript. Ein Umzug auf andere
    Hardware oder eine andere Knotenzahl erfordert Änderungen am Quelltext. Eine
    Auslagerung dieser Werte in eine Konfigurationsdatei wäre der naheliegende
    erste Schritt.

    \item \textbf{Kopierfehler bei Knoten~6.} Im Einsammel-Schritt wird das
    Ergebnis von Knoten~6 versehentlich nach \texttt{ir\_node5.json} kopiert und
    überschreibt damit das Ergebnis von Knoten~5. Effektiv fließen so nur sieben
    statt acht unabhängige Stichproben in die Aggregation ein. Der Fehler ist ein
    Tippfehler im Zielnamen und leicht zu beheben, führt aber vor Augen, wie
    fehleranfällig die von Hand ausgeschriebene Knotenliste ist -- ein weiteres
    Argument für die genannte Auslagerung in eine Schleife über eine
    Konfigurationsliste.

    \item \textbf{Keine Fehlererkennung.} Fällt ein Knoten aus oder liefert er kein
    Ergebnis, schlägt der zugehörige \texttt{scp}-Aufruf fehl, der Gesamtlauf setzt
    sich aber fort und aggregiert stillschweigend über die vorhandenen Dateien. Ein
    Knotenausfall bleibt damit unbemerkt und verfälscht die Stichprobengröße. Eine
    Prüfung der Rückgabecodes und eine Warnung bei fehlenden Ergebnissen wären hier
    angebracht.

    \item \textbf{Manuelle Kopplung an die DSP-Pipeline.} Die Aggregation erzeugt
    \texttt{ir\_merged.json}, die Python-DSP-Pipeline liest jedoch
    \texttt{ir\_output.json}. Der Übergang zwischen verteiltem Lauf und Faltung ist
    derzeit ein manueller Schritt. Eine saubere Integration würde den Cluster-Lauf
    direkt aus der Orchestrierung heraus anstoßen und die zusammengeführte Datei
    automatisch weiterreichen.

    \item \textbf{Statische Verteilung.} Der Replikat-Ansatz mit gleich großen
    Aufträgen ist einfach und robust, setzt aber homogene Hardware voraus und lastet
    unterschiedlich schnelle Knoten nicht unterschiedlich aus. Ein Modell mit
    dynamischer Arbeitsverteilung -- ein Manager vergibt Ray-Kontingente, sobald ein
    Knoten frei wird -- würde heterogene Hardware besser nutzen und die Knotenzahl
    flexibel machen.
\end{itemize}

Alle genannten Punkte betreffen Robustheit und Bedienbarkeit, nicht die
grundsätzliche Korrektheit des Verfahrens. Die verteilte Ausführung liefert eine
gültige, statistisch dichtere Impulsantwort, und das Manager-Worker-Muster über
SSH und \texttt{scp} erfüllt seinen Zweck im Laborbetrieb.

\subsection{Zusammenfassung}
\label{subsec:cluster-zusammenfassung}

Dieses Kapitel hat die verteilte Ausführung der Strahlverfolgung auf dem
Labor-Cluster beschrieben. Weil die Simulation embarrassingly parallel ist, genügt
ein Replikat-Ansatz: Alle acht Knoten rechnen denselben Auftrag mit unabhängigen
Zufallsströmen, was statistisch einer achtfach größeren Simulation entspricht und
den Schätzfehler der Impulsantwort um rund $\sqrt{8}$ senkt -- ein Verhalten, das
das numerische Experiment in Abbildung~\ref{fig:mc-convergence} bestätigt. Ein
Shell-Skript startet die Knoten per SSH parallel, synchronisiert über \texttt{wait}
und sammelt die Ergebnisse per \texttt{scp} ein; \texttt{merge\_results.py}
verkettet die Trefferlisten zu einer gemeinsamen Datei im gewohnten Schema. Der
Prototyp erfüllt seine Aufgabe, weist aber in Konfigurierbarkeit, Fehlererkennung
und der Anbindung an die DSP-Pipeline die diskutierten Schwächen auf, die den
Rahmen für die Weiterentwicklung abstecken.

%   \end{document}
% =====================================================================

\lstdefinestyle{clusterlst}{
    basicstyle=\ttfamily\small,
    showstringspaces=false,
    breaklines=true,
    numbers=left,
    numberstyle=\tiny,
    frame=single,
    tabsize=2
}

%======================================================================
\section{Verteilte Ausführung auf dem Rechencluster}
\label{sec:cluster}

\subsection{Einordnung und Aufgabenstellung}
\label{subsec:cluster-einordnung}

Die Strahlverfolgung ist der teuerste Schritt der gesamten Pipeline. Für eine
brauchbare Impulsantwort müssen Hunderttausende bis Millionen einzelner
Schallstrahlen durch die Raumgeometrie verfolgt werden, und jeder dieser Strahlen
verlangt mehrere Schnitt- und Reflexionsberechnungen. Der entscheidende Vorteil
dabei: Kein Strahl hängt vom Ergebnis eines anderen ab. Das Problem ist damit
\emph{embarrassingly parallel} und lässt sich ohne Kommunikation zwischen den
Teilrechnungen auf beliebig viele Rechner verteilen.

Dieses Kapitel dokumentiert die Verteilungsschicht des Systems. Sie startet die
Rust-Simulation gleichzeitig auf allen acht Knoten des Labor-Clusters, wartet auf
deren Abschluss, sammelt die Teilergebnisse ein und fügt sie zu einer einzigen,
statistisch dichteren Impulsantwort zusammen. Aus Sicht der übrigen Pipeline
ändert sich dabei nichts: Am Ende steht wieder eine Datei im gewohnten Schema,
die die Signalverarbeitung unverändert weiterverarbeiten kann.

Zur Abgrenzung von den Nachbarkapiteln gehören zwei Punkte. Erstens ist die
\emph{prozessinterne} Nebenläufigkeit der Rust-Engine, also die Verteilung auf
die Kerne eines einzelnen Rechners mittels \texttt{rayon}, an anderer Stelle
beschrieben und bleibt hier außen vor. Sie parallelisiert innerhalb eines
Prozesses; dieses Kapitel parallelisiert über Prozess- und Rechnergrenzen hinweg.
Zweitens geht es nicht um den einzelnen Subprozess-Aufruf aus der Python-Pipeline,
sondern um die Ebene darüber: die gleichzeitige Ausführung auf mehreren physischen
Knoten und die anschließende Aggregation ihrer Ergebnisse.

\subsection{Theoretischer Hintergrund}
\label{subsec:cluster-theorie}

\subsubsection{Monte-Carlo-Strahlverfolgung}
\label{subsubsec:mc}

Die Engine schätzt die Impulsantwort eines Raumes stochastisch. Sie verschickt
zufällig gerichtete Strahlen von der Schallquelle, verfolgt ihre Reflexionen und
registriert, wann und mit welchem Restdruck sie den Empfänger erreichen. Jeder
einzelne Strahl ist eine Zufallsstichprobe; erst die Summe vieler Strahlen ergibt
ein aussagekräftiges Bild der Raumakustik. Solche Monte-Carlo-Verfahren haben eine
charakteristische Eigenschaft: Ihr Schätzfehler sinkt mit der Wurzel der
Stichprobengröße. Verdoppelt man die Zahl der Strahlen, halbiert sich der Fehler
also nicht, sondern sinkt nur um den Faktor $\sqrt{2}$. Genau diese langsame
Konvergenz macht große Strahlzahlen nötig und motiviert die Verteilung auf mehrere
Rechner.

\subsubsection{Zwei Wege der Verteilung}
\label{subsubsec:strategien}

Ein Monte-Carlo-Verfahren lässt sich auf zwei grundsätzlich verschiedene Arten
verteilen. Tabelle~\ref{tab:strategien} stellt sie gegenüber.

\begin{table}[htbp]
    \centering
    \caption{Zwei Strategien zur Verteilung der Strahlverfolgung.}
    \label{tab:strategien}
    \begin{tabular}{p{3.2cm} p{5.2cm} p{4.2cm}}
        \toprule
        & \textbf{Work-Partitioning} & \textbf{Replikat-Parallelismus} \\
        \midrule
        Prinzip & Gesamtzahl der Strahlen wird in disjunkte Anteile zerlegt & Jeder Knoten rechnet denselben Auftrag mit eigenem Zufallsstrom \\
        Koordination & Zuteilung der Anteile nötig & keine, jeder Knoten ist autark \\
        Zufallszahlen & ein globaler Strom, aufgeteilt & unabhängige Ströme pro Knoten \\
        Robustheit & Ausfall trifft einen festen Anteil & Ausfall kostet nur Stichprobengröße \\
        \bottomrule
    \end{tabular}
\end{table}

Das umgesetzte System verwendet den zweiten Weg. Alle acht Knoten starten die
Simulation mit derselben Konfigurationsdatei \texttt{input\_config.json} und
werfen dennoch jeweils andere Strahlen, weil jeder Prozess seinen eigenen
Zufallsstrom besitzt. Der Ansatz ist bewusst schlicht gewählt: Er braucht keine
zentrale Zuteilung und keinen Austausch zwischen den Knoten, sondern nur, dass am
Ende alle Ergebnisse eingesammelt werden.

\subsubsection{Unabhängigkeit der Zufallsströme}
\label{subsubsec:rng}

Die Korrektheit des Replikat-Ansatzes steht und fällt mit einer Bedingung: Die
Zufallszahlen auf den Knoten müssen unabhängig sein. Würde jeder Knoten dieselbe
Folge zufälliger Richtungen erzeugen, so lieferten alle acht exakt dasselbe
Ergebnis, und das Zusammenführen brächte keinen Erkenntnisgewinn, sondern nur
achtfach dieselben Strahlen. Die Rust-Engine schließt das aus, indem sie ihren
Zufallsstrom pro Thread bzw. Prozess frisch aus der Systementropie zieht:

\begin{lstlisting}[style=clusterlst, language=Rust, caption={Initialisierung des Zufallsstroms pro Prozess (\texttt{simulation.rs}).}]
let mut rng = rand::rng();                       // Omnidirectional, random rays
let random_angle: f32 = rng.random_range(0.0..(2.0 * PI));
let direction = Vec2::new(random_angle.cos(), random_angle.sin());
\end{lstlisting}

Weil jeder der acht Prozesse auf seinem Knoten einen anderen Startzustand erhält,
erzeugt jeder eine andere Stichprobe. Acht Knoten mit je einer Million Strahlen
liefern damit acht Millionen \emph{unabhängige} Stichproben derselben Raumakustik.

\subsubsection{Statistische Konvergenz}
\label{subsubsec:konvergenz}

Statistisch verhält sich der verteilte Lauf wie eine einzige, achtfach größere
Simulation. Für einen Monte-Carlo-Schätzer aus $N$ unabhängigen Teilstichproben
gleicher Größe sinkt die Standardabweichung des Mittelwerts mit
\[
    \sigma_N \;=\; \frac{\sigma_1}{\sqrt{N}} .
\]
Acht Knoten reduzieren die Streuung der geschätzten Impulsantwort demnach um den
Faktor $\sqrt{8} \approx 2{,}83$ gegenüber einem einzelnen Knoten.
Abbildung~\ref{fig:mc-convergence} bestätigt dieses Verhalten in einem numerischen
Experiment: Aufgetragen ist der mittlere quadratische Fehler der geschätzten
Impulsantwort über der Knotenzahl. Die Messpunkte folgen sauber der
$1/\sqrt{N}$-Referenzkurve; von einem auf acht Knoten sinkt der Fehler um den
Faktor $2{,}78$, praktisch exakt der theoretisch erwartete Wert.

\begin{figure}[htbp]
    \centering
    \includegraphics[width=0.8\textwidth]{figures/mc_convergence}
    \caption{Mittlerer quadratischer Fehler der geschätzten Impulsantwort über der
    Knotenzahl $N$ (je $5\cdot10^4$ Strahlen, doppelt logarithmisch). Die
    Messpunkte folgen der $1/\sqrt{N}$-Referenz; acht Knoten senken den Fehler um
    den Faktor $2{,}78 \approx \sqrt{8}$.}
    \label{fig:mc-convergence}
\end{figure}

Bemerkenswert ist die Kehrseite dieser Kurve: Um den Fehler zu halbieren, braucht
es die vierfache Knotenzahl. Der Nutzen zusätzlicher Knoten nimmt also ab, was bei
der Wahl der Clustergröße zu berücksichtigen ist. Für die hier angestrebte
Genauigkeit sind acht Knoten ein sinnvoller Kompromiss zwischen Aufwand und
Qualität.

\subsection{Architektur des Clusters}
\label{subsec:cluster-arch}

Das System folgt einem \emph{Manager-Worker}-Muster
(Abbildung~\ref{fig:cluster-arch}). Ein Manager-Knoten koordiniert den Ablauf,
acht Worker-Knoten führen die eigentliche Simulation aus. Der Manager selbst
rechnet nicht mit, sondern startet die Worker, wartet auf sie und führt am Ende
die Ergebnisse zusammen.

\begin{figure}[htbp]
    \centering
    \includegraphics[width=0.82\textwidth]{figures/cluster_arch}
    \caption{Manager-Worker-Architektur. Der Manager startet die acht Knoten
    parallel per \texttt{ssh}, holt nach Abschluss jede Ergebnisdatei per
    \texttt{scp} zurück und führt sie mit \texttt{merge\_results.py} zusammen.}
    \label{fig:cluster-arch}
\end{figure}

Die Kommunikation läuft vollständig über Standard-Werkzeuge des Betriebssystems.
\texttt{ssh} startet die Prozesse aus der Ferne, \texttt{scp} überträgt die
Ergebnisdateien zurück. Ein eigenes Netzwerkprotokoll oder ein
Message-Passing-Framework kommt nicht zum Einsatz. Für einen Stapelbetrieb, bei
dem jeder Knoten genau einmal startet, eine Weile rechnet und ein einziges
Ergebnis zurückliefert, wäre solcher Aufwand nicht gerechtfertigt: Es gibt keine
laufende Zustandssynchronisation und keinen Nachrichtenaustausch während der
Rechnung, den ein Framework verwalten müsste.

Damit das funktioniert, müssen drei Voraussetzungen erfüllt sein. Auf jedem Knoten
liegt dieselbe vorkompilierte Binärdatei \texttt{rust\_service} samt der
Konfiguration \texttt{input\_config.json}. Der Manager kann sich per SSH ohne
interaktive Passworteingabe anmelden, was eine Schlüssel-basierte
Authentifizierung voraussetzt. Und die Knoten sind über feste Adressen im
Labornetz erreichbar.

\subsection{Ablaufsteuerung: \texttt{run\_cluster.sh}}
\label{subsec:cluster-run}

Das Skript \texttt{run\_cluster.sh} steuert den verteilten Lauf in drei Phasen:
paralleles Starten, Warten und Einsammeln.

\begin{lstlisting}[style=clusterlst, language=bash, caption={Verteilter Start und Einsammeln der Ergebnisse (gekürzt).}]
# Phase 1: alle Knoten gleichzeitig starten (Hintergrundprozesse)
ssh node1@139.6.65.29 "./rust_service --config-file input_config.json" &
ssh node2@139.6.65.18 "./rust_service --config-file input_config.json" &
# ... node3 bis node8 analog ...

wait                       # Phase 2: auf Abschluss ALLER Knoten warten

# Phase 3: Ergebnis jedes Knotens zum Manager kopieren
scp node1@139.6.65.29:~/ir_output.json ~/ir_node1.json
scp node2@139.6.65.18:~/ir_output.json ~/ir_node2.json
# ... node3 bis node8 analog ...
\end{lstlisting}

Zwei Eigenschaften der Shell nutzt der Ablauf gezielt. Das kaufmännische Und
(\texttt{\&}) am Zeilenende schickt jeden \texttt{ssh}-Aufruf in den Hintergrund.
Die acht Simulationen laufen dadurch nahezu gleichzeitig an, statt nacheinander
abgearbeitet zu werden. Das darauffolgende \texttt{wait} blockiert, bis
\emph{alle} Hintergrundprozesse beendet sind, und bildet so eine
Synchronisationsbarriere zwischen der Rechen- und der Einsammel-Phase. Erst wenn
der letzte Knoten fertig ist, beginnt \texttt{scp}, von jedem Knoten dessen
\texttt{ir\_output.json} zu holen und unter einem knotenspezifischen Namen auf dem
Manager abzulegen.

Diese Barriere hat eine wichtige Konsequenz für die Laufzeit: Der Gesamtlauf ist
so langsam wie der \emph{langsamste} Knoten. Solange die Hardware homogen ist,
fällt das kaum ins Gewicht; auf ungleichen Knoten würde ein einzelner langsamer
Rechner den ganzen Lauf ausbremsen, während die übrigen sieben bereits fertig
warten. Dieses Straggler-Problem ist der zentrale Grund, warum ein statischer
Replikat-Ansatz mit gleich großen Aufträgen homogene Hardware voraussetzt.

\subsection{Aggregation der Teilergebnisse: \texttt{merge\_results.py}}
\label{subsec:cluster-merge}

Sobald alle Teilergebnisse auf dem Manager liegen, führt \texttt{merge\_results.py}
sie zu einer einzigen Impulsantwort zusammen. Jeder Knoten liefert dasselbe
Ausgabeschema, nämlich Metadaten sowie zwei positionsgleiche Listen aus
Verzögerungen und Drücken. Die Aggregation besteht deshalb im Kern aus einer
schlichten Verkettung dieser Listen:

\begin{lstlisting}[style=clusterlst, language=Python, caption={Zusammenführen der Knoten-Ergebnisse, gekürzt (\texttt{merge\_results.py}).}]
all_delays, all_pressures = [], []
for f in sorted(glob.glob("/home/manager/ir_node*.json")):
    with open(f) as file:
        data = json.load(file)
        all_delays    += data["hits"]["delays_seconds"]
        all_pressures += data["hits"]["pressures"]

merged = {
    "metadata": {"sample_rate": 44100, "rays_cast": 1000000 * 8,
                 "rays_received": len(all_delays), "room_name": "MVP_Test_Room_1"},
    "hits": {"delays_seconds": all_delays, "pressures": all_pressures},
}
json.dump(merged, open("/home/manager/ir_merged.json", "w"), indent=2)
\end{lstlisting}

Dass die Verkettung hier die richtige Operation ist, folgt direkt aus dem
Replikat-Modell. Weil die einzelnen Läufe unabhängige Stichproben desselben Raums
sind, ist die Vereinigung aller Treffer statistisch identisch mit dem Ergebnis
eines einzigen Laufs mit acht Millionen Strahlen. Es findet keine Mittelung und
keine Gewichtung statt, sondern nur ein Aneinanderhängen. Das Feld
\texttt{rays\_cast} ist im Skript fest auf $8 \cdot 10^{6}$ gesetzt; solange der in
Abschnitt~\ref{subsec:cluster-grenzen} beschriebene Kopierfehler bei Knoten~6
besteht, fließen tatsächlich nur sieben Knoten in die Aggregation ein, und dieser
Metadatenwert überschätzt die reale Strahlzahl. \texttt{rays\_received} zählt
dagegen die tatsächlich empfangenen Treffer. Die
zusammengeführte Datei trägt exakt dasselbe Schema wie die Ausgabe eines einzelnen
Knotens und ist damit ohne weitere Anpassung für die DSP-Pipeline verwendbar.

\subsection{Datenfluss im Überblick}
\label{subsec:cluster-datenfluss}

Der gesamte verteilte Ablauf lässt sich in fünf Stufen zusammenfassen:

\begin{enumerate}
    \item Der Manager startet per SSH auf jedem der acht Knoten
    \texttt{rust\_service} mit identischer Konfiguration, alle parallel.
    \item Jeder Knoten simuliert unabhängig eine Million Strahlen und schreibt
    lokal seine \texttt{ir\_output.json}.
    \item \texttt{wait} synchronisiert: Die Einsammel-Phase beginnt erst, wenn der
    letzte Knoten fertig ist.
    \item \texttt{scp} überträgt jede \texttt{ir\_output.json} als
    \texttt{ir\_node$k$.json} zum Manager.
    \item \texttt{merge\_results.py} verkettet alle Trefferlisten zu einer
    \texttt{ir\_merged.json}, die anschließend die DSP-Pipeline verarbeitet.
\end{enumerate}

\subsection{Leistungsbetrachtung}
\label{subsec:cluster-leistung}

\subsubsection{Erwarteter Speedup}
\label{subsubsec:speedup}

Bei der Leistung sind zwei Sichtweisen zu trennen. In der umgesetzten
Replikat-Betriebsart rechnet jeder Knoten dieselbe Strahlzahl, sodass acht Knoten
in nahezu derselben Wanduhrzeit die achtfache Stichprobe liefern. Der Gewinn ist
hier keine kürzere Laufzeit, sondern die höhere Genauigkeit aus
Abbildung~\ref{fig:mc-convergence}; man spricht vom schwachen Skalieren im Sinne
von Gustafson, bei dem mit der Rechnerzahl die bearbeitete Problemgröße wächst,
nicht die Geschwindigkeit einer festen Aufgabe. Will man dagegen eine feste
Gesamtstrahlzahl in kürzerer Zeit erreichen, so verteilt man diese auf die Knoten,
sodass jeder nur $8\cdot10^{6}/k$ Strahlen rechnet. Erst in dieser
Partitionierungs-Betriebsart greift das Gesetz von Amdahl. Bezeichnet $p$ den parallelen
Anteil der Arbeit, hier die reine Strahlverfolgung, und $1-p$ den seriellen Rest
aus Prozessstart, Ergebnistransfer und Aggregation, so beträgt der theoretische
Speedup
\[
    S(k) \;=\; \frac{1}{(1-p) + \dfrac{p}{k}} .
\]
Da die Strahlverfolgung die Laufzeit klar dominiert, liegt $p$ nahe bei eins, und
$S(k)$ steigt bis $k=8$ nahezu linear. In der Praxis bleibt der reale Speedup
darunter, und zwar aus drei Gründen: dem bereits erwähnten Straggler-Effekt (der
Lauf ist so schnell wie der langsamste Knoten), dem Overhead für den
SSH-Verbindungsaufbau sowie der Übertragung und Zusammenführung der
Ergebnisdateien.

\subsubsection{Kommunikations- und I/O-Overhead}
\label{subsubsec:overhead}

Der Overhead dieses Ansatzes ist klein, aber nicht null. Pro Knoten fallen ein
SSH-Verbindungsaufbau und eine \texttt{scp}-Übertragung an. Übertragen wird jeweils
eine JSON-Datei, deren Größe mit der Zahl der empfangenen Treffer wächst; bei
einer Million geworfener Strahlen liegt sie im Bereich einiger Megabyte. Da diese
acht Übertragungen erst \emph{nach} der Rechenphase stattfinden und nicht mit ihr
überlappen, addieren sie sich voll zur Gesamtlaufzeit. Bei deutlich mehr Knoten
oder größeren Ergebnisdateien würde dieser sequenzielle Einsammel-Schritt zum
Flaschenhals; eine Überlappung von Rechnung und Transfer oder ein kompakteres
Ausgabeformat als das gut lesbare, aber platzintensive JSON wären dann die
naheliegenden Gegenmaßnahmen.

\subsubsection{Messung am Cluster}
\label{subsubsec:messung}

Die konkreten Laufzeiten hängen von der Laborhardware ab und sind am Cluster zu
messen. Um einen aussagekräftigen Speedup zu erhalten, wird in der
Partitionierungs-Betriebsart gemessen: In \texttt{input\_config.json} wird
\texttt{rays\_to\_cast} auf $8\cdot10^{6}/k$ gesetzt, sodass die Gesamtarbeit über
alle $k$ Knoten konstant bleibt. Dann wird \texttt{run\_cluster.sh} für
verschiedene Knotenzahlen mit der Shell-eigenen Zeitmessung ausgeführt:

\begin{lstlisting}[style=clusterlst, language=bash, caption={Messvorschrift für die Laufzeit am Cluster.}]
time ./run_cluster.sh        # Wanduhrzeit fuer den kompletten verteilten Lauf
\end{lstlisting}

Tabelle~\ref{tab:speedup} nimmt die Messwerte auf. Als Referenz dient der Lauf mit
einem einzelnen Knoten; der Speedup ergibt sich als Verhältnis der Wanduhrzeiten.

\begin{table}[htbp]
    \centering
    \caption{Wanduhrzeit und Speedup des verteilten Laufs (am Cluster zu messen).}
    \label{tab:speedup}
    \begin{tabular}{r r r}
        \toprule
        Knoten $k$ & Wanduhrzeit [s] & Speedup $S$ \\
        \midrule
        1 & \dots & $1{,}00$ \\
        2 & \dots & \dots \\
        4 & \dots & \dots \\
        8 & \dots & \dots \\
        \bottomrule
    \end{tabular}
\end{table}

\subsection{Fehlertoleranz und Sicherheit}
\label{subsec:cluster-fehler}

Zwei Querschnittsthemen verdienen eigene Betrachtung. Beim \textbf{Sicherheits}aspekt
ist die Schlüssel-basierte SSH-Authentifizierung zentral: Der Manager benötigt
passwortlosen Zugriff auf jeden Knoten, was über hinterlegte öffentliche Schlüssel
gelöst ist. Der private Schlüssel des Managers ist damit das entscheidende
Geheimnis; wer ihn besitzt, kann auf allen Knoten Code ausführen. In einem
abgeschlossenen Labornetz ist das vertretbar, für einen offeneren Betrieb wäre eine
engere Rechtevergabe angebracht.

Beim Thema \textbf{Fehlertoleranz} zeigt der aktuelle Stand seine deutlichste
Schwäche, die im folgenden Abschnitt zusammen mit den übrigen Grenzen diskutiert
wird. Der Kern des Problems: Ein ausgefallener Knoten wird nicht erkannt, sondern
schlicht übergangen, sodass die Stichprobengröße unbemerkt sinkt.

\subsection{Grenzen des aktuellen Standes und Ausblick}
\label{subsec:cluster-grenzen}

Die Verteilungsschicht ist als funktionsfähiger Prototyp umgesetzt. Für einen
robusten Betrieb bestehen mehrere klar benannte Schwachstellen, die zugleich die
nächsten Entwicklungsschritte vorgeben.

\begin{itemize}
    \item \textbf{Fest verdrahtete Konfiguration.} IP-Adressen, Knotennamen, die
    Knotenzahl und die Ablagepfade stehen direkt im Skript. Ein Umzug auf andere
    Hardware oder eine andere Knotenzahl erfordert Änderungen am Quelltext. Eine
    Auslagerung dieser Werte in eine Konfigurationsdatei wäre der naheliegende
    erste Schritt.

    \item \textbf{Kopierfehler bei Knoten~6.} Im Einsammel-Schritt wird das
    Ergebnis von Knoten~6 versehentlich nach \texttt{ir\_node5.json} kopiert und
    überschreibt damit das Ergebnis von Knoten~5. Effektiv fließen so nur sieben
    statt acht unabhängige Stichproben in die Aggregation ein. Der Fehler ist ein
    Tippfehler im Zielnamen und leicht zu beheben, führt aber vor Augen, wie
    fehleranfällig die von Hand ausgeschriebene Knotenliste ist -- ein weiteres
    Argument für die genannte Auslagerung in eine Schleife über eine
    Konfigurationsliste.

    \item \textbf{Keine Fehlererkennung.} Fällt ein Knoten aus oder liefert er kein
    Ergebnis, schlägt der zugehörige \texttt{scp}-Aufruf fehl, der Gesamtlauf setzt
    sich aber fort und aggregiert stillschweigend über die vorhandenen Dateien. Ein
    Knotenausfall bleibt damit unbemerkt und verfälscht die Stichprobengröße. Eine
    Prüfung der Rückgabecodes und eine Warnung bei fehlenden Ergebnissen wären hier
    angebracht.

    \item \textbf{Manuelle Kopplung an die DSP-Pipeline.} Die Aggregation erzeugt
    \texttt{ir\_merged.json}, die Python-DSP-Pipeline liest jedoch
    \texttt{ir\_output.json}. Der Übergang zwischen verteiltem Lauf und Faltung ist
    derzeit ein manueller Schritt. Eine saubere Integration würde den Cluster-Lauf
    direkt aus der Orchestrierung heraus anstoßen und die zusammengeführte Datei
    automatisch weiterreichen.

    \item \textbf{Statische Verteilung.} Der Replikat-Ansatz mit gleich großen
    Aufträgen ist einfach und robust, setzt aber homogene Hardware voraus und lastet
    unterschiedlich schnelle Knoten nicht unterschiedlich aus. Ein Modell mit
    dynamischer Arbeitsverteilung -- ein Manager vergibt Ray-Kontingente, sobald ein
    Knoten frei wird -- würde heterogene Hardware besser nutzen und die Knotenzahl
    flexibel machen.
\end{itemize}

Alle genannten Punkte betreffen Robustheit und Bedienbarkeit, nicht die
grundsätzliche Korrektheit des Verfahrens. Die verteilte Ausführung liefert eine
gültige, statistisch dichtere Impulsantwort, und das Manager-Worker-Muster über
SSH und \texttt{scp} erfüllt seinen Zweck im Laborbetrieb.

\subsection{Zusammenfassung}
\label{subsec:cluster-zusammenfassung}

Dieses Kapitel hat die verteilte Ausführung der Strahlverfolgung auf dem
Labor-Cluster beschrieben. Weil die Simulation embarrassingly parallel ist, genügt
ein Replikat-Ansatz: Alle acht Knoten rechnen denselben Auftrag mit unabhängigen
Zufallsströmen, was statistisch einer achtfach größeren Simulation entspricht und
den Schätzfehler der Impulsantwort um rund $\sqrt{8}$ senkt -- ein Verhalten, das
das numerische Experiment in Abbildung~\ref{fig:mc-convergence} bestätigt. Ein
Shell-Skript startet die Knoten per SSH parallel, synchronisiert über \texttt{wait}
und sammelt die Ergebnisse per \texttt{scp} ein; \texttt{merge\_results.py}
verkettet die Trefferlisten zu einer gemeinsamen Datei im gewohnten Schema. Der
Prototyp erfüllt seine Aufgabe, weist aber in Konfigurierbarkeit, Fehlererkennung
und der Anbindung an die DSP-Pipeline die diskutierten Schwächen auf, die den
Rahmen für die Weiterentwicklung abstecken.

%   \end{document}
% =====================================================================

\lstdefinestyle{clusterlst}{
    basicstyle=\ttfamily\small,
    showstringspaces=false,
    breaklines=true,
    numbers=left,
    numberstyle=\tiny,
    frame=single,
    tabsize=2
}

%======================================================================
\section{Verteilte Ausführung auf dem Rechencluster}
\label{sec:cluster}

\subsection{Einordnung und Aufgabenstellung}
\label{subsec:cluster-einordnung}

Die Strahlverfolgung ist der teuerste Schritt der gesamten Pipeline. Für eine
brauchbare Impulsantwort müssen Hunderttausende bis Millionen einzelner
Schallstrahlen durch die Raumgeometrie verfolgt werden, und jeder dieser Strahlen
verlangt mehrere Schnitt- und Reflexionsberechnungen. Der entscheidende Vorteil
dabei: Kein Strahl hängt vom Ergebnis eines anderen ab. Das Problem ist damit
\emph{embarrassingly parallel} und lässt sich ohne Kommunikation zwischen den
Teilrechnungen auf beliebig viele Rechner verteilen.

Dieses Kapitel dokumentiert die Verteilungsschicht des Systems. Sie startet die
Rust-Simulation gleichzeitig auf allen acht Knoten des Labor-Clusters, wartet auf
deren Abschluss, sammelt die Teilergebnisse ein und fügt sie zu einer einzigen,
statistisch dichteren Impulsantwort zusammen. Aus Sicht der übrigen Pipeline
ändert sich dabei nichts: Am Ende steht wieder eine Datei im gewohnten Schema,
die die Signalverarbeitung unverändert weiterverarbeiten kann.

Zur Abgrenzung von den Nachbarkapiteln gehören zwei Punkte. Erstens ist die
\emph{prozessinterne} Nebenläufigkeit der Rust-Engine, also die Verteilung auf
die Kerne eines einzelnen Rechners mittels \texttt{rayon}, an anderer Stelle
beschrieben und bleibt hier außen vor. Sie parallelisiert innerhalb eines
Prozesses; dieses Kapitel parallelisiert über Prozess- und Rechnergrenzen hinweg.
Zweitens geht es nicht um den einzelnen Subprozess-Aufruf aus der Python-Pipeline,
sondern um die Ebene darüber: die gleichzeitige Ausführung auf mehreren physischen
Knoten und die anschließende Aggregation ihrer Ergebnisse.

\subsection{Theoretischer Hintergrund}
\label{subsec:cluster-theorie}

\subsubsection{Monte-Carlo-Strahlverfolgung}
\label{subsubsec:mc}

Die Engine schätzt die Impulsantwort eines Raumes stochastisch. Sie verschickt
zufällig gerichtete Strahlen von der Schallquelle, verfolgt ihre Reflexionen und
registriert, wann und mit welchem Restdruck sie den Empfänger erreichen. Jeder
einzelne Strahl ist eine Zufallsstichprobe; erst die Summe vieler Strahlen ergibt
ein aussagekräftiges Bild der Raumakustik. Solche Monte-Carlo-Verfahren haben eine
charakteristische Eigenschaft: Ihr Schätzfehler sinkt mit der Wurzel der
Stichprobengröße. Verdoppelt man die Zahl der Strahlen, halbiert sich der Fehler
also nicht, sondern sinkt nur um den Faktor $\sqrt{2}$. Genau diese langsame
Konvergenz macht große Strahlzahlen nötig und motiviert die Verteilung auf mehrere
Rechner.

\subsubsection{Zwei Wege der Verteilung}
\label{subsubsec:strategien}

Ein Monte-Carlo-Verfahren lässt sich auf zwei grundsätzlich verschiedene Arten
verteilen. Tabelle~\ref{tab:strategien} stellt sie gegenüber.

\begin{table}[htbp]
    \centering
    \caption{Zwei Strategien zur Verteilung der Strahlverfolgung.}
    \label{tab:strategien}
    \begin{tabular}{p{3.2cm} p{5.2cm} p{4.2cm}}
        \toprule
        & \textbf{Work-Partitioning} & \textbf{Replikat-Parallelismus} \\
        \midrule
        Prinzip & Gesamtzahl der Strahlen wird in disjunkte Anteile zerlegt & Jeder Knoten rechnet denselben Auftrag mit eigenem Zufallsstrom \\
        Koordination & Zuteilung der Anteile nötig & keine, jeder Knoten ist autark \\
        Zufallszahlen & ein globaler Strom, aufgeteilt & unabhängige Ströme pro Knoten \\
        Robustheit & Ausfall trifft einen festen Anteil & Ausfall kostet nur Stichprobengröße \\
        \bottomrule
    \end{tabular}
\end{table}

Das umgesetzte System verwendet den zweiten Weg. Alle acht Knoten starten die
Simulation mit derselben Konfigurationsdatei \texttt{input\_config.json} und
werfen dennoch jeweils andere Strahlen, weil jeder Prozess seinen eigenen
Zufallsstrom besitzt. Der Ansatz ist bewusst schlicht gewählt: Er braucht keine
zentrale Zuteilung und keinen Austausch zwischen den Knoten, sondern nur, dass am
Ende alle Ergebnisse eingesammelt werden.

\subsubsection{Unabhängigkeit der Zufallsströme}
\label{subsubsec:rng}

Die Korrektheit des Replikat-Ansatzes steht und fällt mit einer Bedingung: Die
Zufallszahlen auf den Knoten müssen unabhängig sein. Würde jeder Knoten dieselbe
Folge zufälliger Richtungen erzeugen, so lieferten alle acht exakt dasselbe
Ergebnis, und das Zusammenführen brächte keinen Erkenntnisgewinn, sondern nur
achtfach dieselben Strahlen. Die Rust-Engine schließt das aus, indem sie ihren
Zufallsstrom pro Thread bzw. Prozess frisch aus der Systementropie zieht:

\begin{lstlisting}[style=clusterlst, language=Rust, caption={Initialisierung des Zufallsstroms pro Prozess (\texttt{simulation.rs}).}]
let mut rng = rand::rng();                       // Omnidirectional, random rays
let random_angle: f32 = rng.random_range(0.0..(2.0 * PI));
let direction = Vec2::new(random_angle.cos(), random_angle.sin());
\end{lstlisting}

Weil jeder der acht Prozesse auf seinem Knoten einen anderen Startzustand erhält,
erzeugt jeder eine andere Stichprobe. Acht Knoten mit je einer Million Strahlen
liefern damit acht Millionen \emph{unabhängige} Stichproben derselben Raumakustik.

\subsubsection{Statistische Konvergenz}
\label{subsubsec:konvergenz}

Statistisch verhält sich der verteilte Lauf wie eine einzige, achtfach größere
Simulation. Für einen Monte-Carlo-Schätzer aus $N$ unabhängigen Teilstichproben
gleicher Größe sinkt die Standardabweichung des Mittelwerts mit
\[
    \sigma_N \;=\; \frac{\sigma_1}{\sqrt{N}} .
\]
Acht Knoten reduzieren die Streuung der geschätzten Impulsantwort demnach um den
Faktor $\sqrt{8} \approx 2{,}83$ gegenüber einem einzelnen Knoten.
Abbildung~\ref{fig:mc-convergence} bestätigt dieses Verhalten in einem numerischen
Experiment: Aufgetragen ist der mittlere quadratische Fehler der geschätzten
Impulsantwort über der Knotenzahl. Die Messpunkte folgen sauber der
$1/\sqrt{N}$-Referenzkurve; von einem auf acht Knoten sinkt der Fehler um den
Faktor $2{,}78$, praktisch exakt der theoretisch erwartete Wert.

\begin{figure}[htbp]
    \centering
    \includegraphics[width=0.8\textwidth]{figures/mc_convergence}
    \caption{Mittlerer quadratischer Fehler der geschätzten Impulsantwort über der
    Knotenzahl $N$ (je $5\cdot10^4$ Strahlen, doppelt logarithmisch). Die
    Messpunkte folgen der $1/\sqrt{N}$-Referenz; acht Knoten senken den Fehler um
    den Faktor $2{,}78 \approx \sqrt{8}$.}
    \label{fig:mc-convergence}
\end{figure}

Bemerkenswert ist die Kehrseite dieser Kurve: Um den Fehler zu halbieren, braucht
es die vierfache Knotenzahl. Der Nutzen zusätzlicher Knoten nimmt also ab, was bei
der Wahl der Clustergröße zu berücksichtigen ist. Für die hier angestrebte
Genauigkeit sind acht Knoten ein sinnvoller Kompromiss zwischen Aufwand und
Qualität.

\subsection{Architektur des Clusters}
\label{subsec:cluster-arch}

Das System folgt einem \emph{Manager-Worker}-Muster
(Abbildung~\ref{fig:cluster-arch}). Ein Manager-Knoten koordiniert den Ablauf,
acht Worker-Knoten führen die eigentliche Simulation aus. Der Manager selbst
rechnet nicht mit, sondern startet die Worker, wartet auf sie und führt am Ende
die Ergebnisse zusammen.

\begin{figure}[htbp]
    \centering
    \includegraphics[width=0.82\textwidth]{figures/cluster_arch}
    \caption{Manager-Worker-Architektur. Der Manager startet die acht Knoten
    parallel per \texttt{ssh}, holt nach Abschluss jede Ergebnisdatei per
    \texttt{scp} zurück und führt sie mit \texttt{merge\_results.py} zusammen.}
    \label{fig:cluster-arch}
\end{figure}

Die Kommunikation läuft vollständig über Standard-Werkzeuge des Betriebssystems.
\texttt{ssh} startet die Prozesse aus der Ferne, \texttt{scp} überträgt die
Ergebnisdateien zurück. Ein eigenes Netzwerkprotokoll oder ein
Message-Passing-Framework kommt nicht zum Einsatz. Für einen Stapelbetrieb, bei
dem jeder Knoten genau einmal startet, eine Weile rechnet und ein einziges
Ergebnis zurückliefert, wäre solcher Aufwand nicht gerechtfertigt: Es gibt keine
laufende Zustandssynchronisation und keinen Nachrichtenaustausch während der
Rechnung, den ein Framework verwalten müsste.

Damit das funktioniert, müssen drei Voraussetzungen erfüllt sein. Auf jedem Knoten
liegt dieselbe vorkompilierte Binärdatei \texttt{rust\_service} samt der
Konfiguration \texttt{input\_config.json}. Der Manager kann sich per SSH ohne
interaktive Passworteingabe anmelden, was eine Schlüssel-basierte
Authentifizierung voraussetzt. Und die Knoten sind über feste Adressen im
Labornetz erreichbar.

\subsection{Ablaufsteuerung: \texttt{run\_cluster.sh}}
\label{subsec:cluster-run}

Das Skript \texttt{run\_cluster.sh} steuert den verteilten Lauf in drei Phasen:
paralleles Starten, Warten und Einsammeln.

\begin{lstlisting}[style=clusterlst, language=bash, caption={Verteilter Start und Einsammeln der Ergebnisse (gekürzt).}]
# Phase 1: alle Knoten gleichzeitig starten (Hintergrundprozesse)
ssh node1@139.6.65.29 "./rust_service --config-file input_config.json" &
ssh node2@139.6.65.18 "./rust_service --config-file input_config.json" &
# ... node3 bis node8 analog ...

wait                       # Phase 2: auf Abschluss ALLER Knoten warten

# Phase 3: Ergebnis jedes Knotens zum Manager kopieren
scp node1@139.6.65.29:~/ir_output.json ~/ir_node1.json
scp node2@139.6.65.18:~/ir_output.json ~/ir_node2.json
# ... node3 bis node8 analog ...
\end{lstlisting}

Zwei Eigenschaften der Shell nutzt der Ablauf gezielt. Das kaufmännische Und
(\texttt{\&}) am Zeilenende schickt jeden \texttt{ssh}-Aufruf in den Hintergrund.
Die acht Simulationen laufen dadurch nahezu gleichzeitig an, statt nacheinander
abgearbeitet zu werden. Das darauffolgende \texttt{wait} blockiert, bis
\emph{alle} Hintergrundprozesse beendet sind, und bildet so eine
Synchronisationsbarriere zwischen der Rechen- und der Einsammel-Phase. Erst wenn
der letzte Knoten fertig ist, beginnt \texttt{scp}, von jedem Knoten dessen
\texttt{ir\_output.json} zu holen und unter einem knotenspezifischen Namen auf dem
Manager abzulegen.

Diese Barriere hat eine wichtige Konsequenz für die Laufzeit: Der Gesamtlauf ist
so langsam wie der \emph{langsamste} Knoten. Solange die Hardware homogen ist,
fällt das kaum ins Gewicht; auf ungleichen Knoten würde ein einzelner langsamer
Rechner den ganzen Lauf ausbremsen, während die übrigen sieben bereits fertig
warten. Dieses Straggler-Problem ist der zentrale Grund, warum ein statischer
Replikat-Ansatz mit gleich großen Aufträgen homogene Hardware voraussetzt.

\subsection{Aggregation der Teilergebnisse: \texttt{merge\_results.py}}
\label{subsec:cluster-merge}

Sobald alle Teilergebnisse auf dem Manager liegen, führt \texttt{merge\_results.py}
sie zu einer einzigen Impulsantwort zusammen. Jeder Knoten liefert dasselbe
Ausgabeschema, nämlich Metadaten sowie zwei positionsgleiche Listen aus
Verzögerungen und Drücken. Die Aggregation besteht deshalb im Kern aus einer
schlichten Verkettung dieser Listen:

\begin{lstlisting}[style=clusterlst, language=Python, caption={Zusammenführen der Knoten-Ergebnisse, gekürzt (\texttt{merge\_results.py}).}]
all_delays, all_pressures = [], []
for f in sorted(glob.glob("/home/manager/ir_node*.json")):
    with open(f) as file:
        data = json.load(file)
        all_delays    += data["hits"]["delays_seconds"]
        all_pressures += data["hits"]["pressures"]

merged = {
    "metadata": {"sample_rate": 44100, "rays_cast": 1000000 * 8,
                 "rays_received": len(all_delays), "room_name": "MVP_Test_Room_1"},
    "hits": {"delays_seconds": all_delays, "pressures": all_pressures},
}
json.dump(merged, open("/home/manager/ir_merged.json", "w"), indent=2)
\end{lstlisting}

Dass die Verkettung hier die richtige Operation ist, folgt direkt aus dem
Replikat-Modell. Weil die einzelnen Läufe unabhängige Stichproben desselben Raums
sind, ist die Vereinigung aller Treffer statistisch identisch mit dem Ergebnis
eines einzigen Laufs mit acht Millionen Strahlen. Es findet keine Mittelung und
keine Gewichtung statt, sondern nur ein Aneinanderhängen. Das Feld
\texttt{rays\_cast} ist im Skript fest auf $8 \cdot 10^{6}$ gesetzt; solange der in
Abschnitt~\ref{subsec:cluster-grenzen} beschriebene Kopierfehler bei Knoten~6
besteht, fließen tatsächlich nur sieben Knoten in die Aggregation ein, und dieser
Metadatenwert überschätzt die reale Strahlzahl. \texttt{rays\_received} zählt
dagegen die tatsächlich empfangenen Treffer. Die
zusammengeführte Datei trägt exakt dasselbe Schema wie die Ausgabe eines einzelnen
Knotens und ist damit ohne weitere Anpassung für die DSP-Pipeline verwendbar.

\subsection{Datenfluss im Überblick}
\label{subsec:cluster-datenfluss}

Der gesamte verteilte Ablauf lässt sich in fünf Stufen zusammenfassen:

\begin{enumerate}
    \item Der Manager startet per SSH auf jedem der acht Knoten
    \texttt{rust\_service} mit identischer Konfiguration, alle parallel.
    \item Jeder Knoten simuliert unabhängig eine Million Strahlen und schreibt
    lokal seine \texttt{ir\_output.json}.
    \item \texttt{wait} synchronisiert: Die Einsammel-Phase beginnt erst, wenn der
    letzte Knoten fertig ist.
    \item \texttt{scp} überträgt jede \texttt{ir\_output.json} als
    \texttt{ir\_node$k$.json} zum Manager.
    \item \texttt{merge\_results.py} verkettet alle Trefferlisten zu einer
    \texttt{ir\_merged.json}, die anschließend die DSP-Pipeline verarbeitet.
\end{enumerate}

\subsection{Leistungsbetrachtung}
\label{subsec:cluster-leistung}

\subsubsection{Erwarteter Speedup}
\label{subsubsec:speedup}

Bei der Leistung sind zwei Sichtweisen zu trennen. In der umgesetzten
Replikat-Betriebsart rechnet jeder Knoten dieselbe Strahlzahl, sodass acht Knoten
in nahezu derselben Wanduhrzeit die achtfache Stichprobe liefern. Der Gewinn ist
hier keine kürzere Laufzeit, sondern die höhere Genauigkeit aus
Abbildung~\ref{fig:mc-convergence}; man spricht vom schwachen Skalieren im Sinne
von Gustafson, bei dem mit der Rechnerzahl die bearbeitete Problemgröße wächst,
nicht die Geschwindigkeit einer festen Aufgabe. Will man dagegen eine feste
Gesamtstrahlzahl in kürzerer Zeit erreichen, so verteilt man diese auf die Knoten,
sodass jeder nur $8\cdot10^{6}/k$ Strahlen rechnet. Erst in dieser
Partitionierungs-Betriebsart greift das Gesetz von Amdahl. Bezeichnet $p$ den parallelen
Anteil der Arbeit, hier die reine Strahlverfolgung, und $1-p$ den seriellen Rest
aus Prozessstart, Ergebnistransfer und Aggregation, so beträgt der theoretische
Speedup
\[
    S(k) \;=\; \frac{1}{(1-p) + \dfrac{p}{k}} .
\]
Da die Strahlverfolgung die Laufzeit klar dominiert, liegt $p$ nahe bei eins, und
$S(k)$ steigt bis $k=8$ nahezu linear. In der Praxis bleibt der reale Speedup
darunter, und zwar aus drei Gründen: dem bereits erwähnten Straggler-Effekt (der
Lauf ist so schnell wie der langsamste Knoten), dem Overhead für den
SSH-Verbindungsaufbau sowie der Übertragung und Zusammenführung der
Ergebnisdateien.

\subsubsection{Kommunikations- und I/O-Overhead}
\label{subsubsec:overhead}

Der Overhead dieses Ansatzes ist klein, aber nicht null. Pro Knoten fallen ein
SSH-Verbindungsaufbau und eine \texttt{scp}-Übertragung an. Übertragen wird jeweils
eine JSON-Datei, deren Größe mit der Zahl der empfangenen Treffer wächst; bei
einer Million geworfener Strahlen liegt sie im Bereich einiger Megabyte. Da diese
acht Übertragungen erst \emph{nach} der Rechenphase stattfinden und nicht mit ihr
überlappen, addieren sie sich voll zur Gesamtlaufzeit. Bei deutlich mehr Knoten
oder größeren Ergebnisdateien würde dieser sequenzielle Einsammel-Schritt zum
Flaschenhals; eine Überlappung von Rechnung und Transfer oder ein kompakteres
Ausgabeformat als das gut lesbare, aber platzintensive JSON wären dann die
naheliegenden Gegenmaßnahmen.

\subsubsection{Messung am Cluster}
\label{subsubsec:messung}

Die konkreten Laufzeiten hängen von der Laborhardware ab und sind am Cluster zu
messen. Um einen aussagekräftigen Speedup zu erhalten, wird in der
Partitionierungs-Betriebsart gemessen: In \texttt{input\_config.json} wird
\texttt{rays\_to\_cast} auf $8\cdot10^{6}/k$ gesetzt, sodass die Gesamtarbeit über
alle $k$ Knoten konstant bleibt. Dann wird \texttt{run\_cluster.sh} für
verschiedene Knotenzahlen mit der Shell-eigenen Zeitmessung ausgeführt:

\begin{lstlisting}[style=clusterlst, language=bash, caption={Messvorschrift für die Laufzeit am Cluster.}]
time ./run_cluster.sh        # Wanduhrzeit fuer den kompletten verteilten Lauf
\end{lstlisting}

Tabelle~\ref{tab:speedup} nimmt die Messwerte auf. Als Referenz dient der Lauf mit
einem einzelnen Knoten; der Speedup ergibt sich als Verhältnis der Wanduhrzeiten.

\begin{table}[htbp]
    \centering
    \caption{Wanduhrzeit und Speedup des verteilten Laufs (am Cluster zu messen).}
    \label{tab:speedup}
    \begin{tabular}{r r r}
        \toprule
        Knoten $k$ & Wanduhrzeit [s] & Speedup $S$ \\
        \midrule
        1 & \dots & $1{,}00$ \\
        2 & \dots & \dots \\
        4 & \dots & \dots \\
        8 & \dots & \dots \\
        \bottomrule
    \end{tabular}
\end{table}

\subsection{Fehlertoleranz und Sicherheit}
\label{subsec:cluster-fehler}

Zwei Querschnittsthemen verdienen eigene Betrachtung. Beim \textbf{Sicherheits}aspekt
ist die Schlüssel-basierte SSH-Authentifizierung zentral: Der Manager benötigt
passwortlosen Zugriff auf jeden Knoten, was über hinterlegte öffentliche Schlüssel
gelöst ist. Der private Schlüssel des Managers ist damit das entscheidende
Geheimnis; wer ihn besitzt, kann auf allen Knoten Code ausführen. In einem
abgeschlossenen Labornetz ist das vertretbar, für einen offeneren Betrieb wäre eine
engere Rechtevergabe angebracht.

Beim Thema \textbf{Fehlertoleranz} zeigt der aktuelle Stand seine deutlichste
Schwäche, die im folgenden Abschnitt zusammen mit den übrigen Grenzen diskutiert
wird. Der Kern des Problems: Ein ausgefallener Knoten wird nicht erkannt, sondern
schlicht übergangen, sodass die Stichprobengröße unbemerkt sinkt.

\subsection{Grenzen des aktuellen Standes und Ausblick}
\label{subsec:cluster-grenzen}

Die Verteilungsschicht ist als funktionsfähiger Prototyp umgesetzt. Für einen
robusten Betrieb bestehen mehrere klar benannte Schwachstellen, die zugleich die
nächsten Entwicklungsschritte vorgeben.

\begin{itemize}
    \item \textbf{Fest verdrahtete Konfiguration.} IP-Adressen, Knotennamen, die
    Knotenzahl und die Ablagepfade stehen direkt im Skript. Ein Umzug auf andere
    Hardware oder eine andere Knotenzahl erfordert Änderungen am Quelltext. Eine
    Auslagerung dieser Werte in eine Konfigurationsdatei wäre der naheliegende
    erste Schritt.

    \item \textbf{Kopierfehler bei Knoten~6.} Im Einsammel-Schritt wird das
    Ergebnis von Knoten~6 versehentlich nach \texttt{ir\_node5.json} kopiert und
    überschreibt damit das Ergebnis von Knoten~5. Effektiv fließen so nur sieben
    statt acht unabhängige Stichproben in die Aggregation ein. Der Fehler ist ein
    Tippfehler im Zielnamen und leicht zu beheben, führt aber vor Augen, wie
    fehleranfällig die von Hand ausgeschriebene Knotenliste ist -- ein weiteres
    Argument für die genannte Auslagerung in eine Schleife über eine
    Konfigurationsliste.

    \item \textbf{Keine Fehlererkennung.} Fällt ein Knoten aus oder liefert er kein
    Ergebnis, schlägt der zugehörige \texttt{scp}-Aufruf fehl, der Gesamtlauf setzt
    sich aber fort und aggregiert stillschweigend über die vorhandenen Dateien. Ein
    Knotenausfall bleibt damit unbemerkt und verfälscht die Stichprobengröße. Eine
    Prüfung der Rückgabecodes und eine Warnung bei fehlenden Ergebnissen wären hier
    angebracht.

    \item \textbf{Manuelle Kopplung an die DSP-Pipeline.} Die Aggregation erzeugt
    \texttt{ir\_merged.json}, die Python-DSP-Pipeline liest jedoch
    \texttt{ir\_output.json}. Der Übergang zwischen verteiltem Lauf und Faltung ist
    derzeit ein manueller Schritt. Eine saubere Integration würde den Cluster-Lauf
    direkt aus der Orchestrierung heraus anstoßen und die zusammengeführte Datei
    automatisch weiterreichen.

    \item \textbf{Statische Verteilung.} Der Replikat-Ansatz mit gleich großen
    Aufträgen ist einfach und robust, setzt aber homogene Hardware voraus und lastet
    unterschiedlich schnelle Knoten nicht unterschiedlich aus. Ein Modell mit
    dynamischer Arbeitsverteilung -- ein Manager vergibt Ray-Kontingente, sobald ein
    Knoten frei wird -- würde heterogene Hardware besser nutzen und die Knotenzahl
    flexibel machen.
\end{itemize}

Alle genannten Punkte betreffen Robustheit und Bedienbarkeit, nicht die
grundsätzliche Korrektheit des Verfahrens. Die verteilte Ausführung liefert eine
gültige, statistisch dichtere Impulsantwort, und das Manager-Worker-Muster über
SSH und \texttt{scp} erfüllt seinen Zweck im Laborbetrieb.

\subsection{Zusammenfassung}
\label{subsec:cluster-zusammenfassung}

Dieses Kapitel hat die verteilte Ausführung der Strahlverfolgung auf dem
Labor-Cluster beschrieben. Weil die Simulation embarrassingly parallel ist, genügt
ein Replikat-Ansatz: Alle acht Knoten rechnen denselben Auftrag mit unabhängigen
Zufallsströmen, was statistisch einer achtfach größeren Simulation entspricht und
den Schätzfehler der Impulsantwort um rund $\sqrt{8}$ senkt -- ein Verhalten, das
das numerische Experiment in Abbildung~\ref{fig:mc-convergence} bestätigt. Ein
Shell-Skript startet die Knoten per SSH parallel, synchronisiert über \texttt{wait}
und sammelt die Ergebnisse per \texttt{scp} ein; \texttt{merge\_results.py}
verkettet die Trefferlisten zu einer gemeinsamen Datei im gewohnten Schema. Der
Prototyp erfüllt seine Aufgabe, weist aber in Konfigurierbarkeit, Fehlererkennung
und der Anbindung an die DSP-Pipeline die diskutierten Schwächen auf, die den
Rahmen für die Weiterentwicklung abstecken.

%   \end{document}
% =====================================================================

\lstdefinestyle{clusterlst}{
    basicstyle=\ttfamily\small,
    showstringspaces=false,
    breaklines=true,
    numbers=left,
    numberstyle=\tiny,
    frame=single,
    tabsize=2
}

%======================================================================
\section{Verteilte Ausführung auf dem Rechencluster}
\label{sec:cluster}

\subsection{Einordnung und Aufgabenstellung}
\label{subsec:cluster-einordnung}

Die Strahlverfolgung ist der teuerste Schritt der gesamten Pipeline. Für eine
brauchbare Impulsantwort müssen Hunderttausende bis Millionen einzelner
Schallstrahlen durch die Raumgeometrie verfolgt werden, und jeder dieser Strahlen
verlangt mehrere Schnitt- und Reflexionsberechnungen. Der entscheidende Vorteil
dabei: Kein Strahl hängt vom Ergebnis eines anderen ab. Das Problem ist damit
\emph{embarrassingly parallel} und lässt sich ohne Kommunikation zwischen den
Teilrechnungen auf beliebig viele Rechner verteilen.

Dieses Kapitel dokumentiert die Verteilungsschicht des Systems. Sie startet die
Rust-Simulation gleichzeitig auf allen acht Knoten des Labor-Clusters, wartet auf
deren Abschluss, sammelt die Teilergebnisse ein und fügt sie zu einer einzigen,
statistisch dichteren Impulsantwort zusammen. Aus Sicht der übrigen Pipeline
ändert sich dabei nichts: Am Ende steht wieder eine Datei im gewohnten Schema,
die die Signalverarbeitung unverändert weiterverarbeiten kann.

Zur Abgrenzung von den Nachbarkapiteln gehören zwei Punkte. Erstens ist die
\emph{prozessinterne} Nebenläufigkeit der Rust-Engine, also die Verteilung auf
die Kerne eines einzelnen Rechners mittels \texttt{rayon}, an anderer Stelle
beschrieben und bleibt hier außen vor. Sie parallelisiert innerhalb eines
Prozesses; dieses Kapitel parallelisiert über Prozess- und Rechnergrenzen hinweg.
Zweitens geht es nicht um den einzelnen Subprozess-Aufruf aus der Python-Pipeline,
sondern um die Ebene darüber: die gleichzeitige Ausführung auf mehreren physischen
Knoten und die anschließende Aggregation ihrer Ergebnisse.

\subsection{Theoretischer Hintergrund}
\label{subsec:cluster-theorie}

\subsubsection{Monte-Carlo-Strahlverfolgung}
\label{subsubsec:mc}

Die Engine schätzt die Impulsantwort eines Raumes stochastisch. Sie verschickt
zufällig gerichtete Strahlen von der Schallquelle, verfolgt ihre Reflexionen und
registriert, wann und mit welchem Restdruck sie den Empfänger erreichen. Jeder
einzelne Strahl ist eine Zufallsstichprobe; erst die Summe vieler Strahlen ergibt
ein aussagekräftiges Bild der Raumakustik. Solche Monte-Carlo-Verfahren haben eine
charakteristische Eigenschaft: Ihr Schätzfehler sinkt mit der Wurzel der
Stichprobengröße. Verdoppelt man die Zahl der Strahlen, halbiert sich der Fehler
also nicht, sondern sinkt nur um den Faktor $\sqrt{2}$. Genau diese langsame
Konvergenz macht große Strahlzahlen nötig und motiviert die Verteilung auf mehrere
Rechner.

\subsubsection{Zwei Wege der Verteilung}
\label{subsubsec:strategien}

Ein Monte-Carlo-Verfahren lässt sich auf zwei grundsätzlich verschiedene Arten
verteilen. Tabelle~\ref{tab:strategien} stellt sie gegenüber.

\begin{table}[htbp]
    \centering
    \caption{Zwei Strategien zur Verteilung der Strahlverfolgung.}
    \label{tab:strategien}
    \begin{tabular}{p{3.2cm} p{5.2cm} p{4.2cm}}
        \toprule
        & \textbf{Work-Partitioning} & \textbf{Replikat-Parallelismus} \\
        \midrule
        Prinzip & Gesamtzahl der Strahlen wird in disjunkte Anteile zerlegt & Jeder Knoten rechnet denselben Auftrag mit eigenem Zufallsstrom \\
        Koordination & Zuteilung der Anteile nötig & keine, jeder Knoten ist autark \\
        Zufallszahlen & ein globaler Strom, aufgeteilt & unabhängige Ströme pro Knoten \\
        Robustheit & Ausfall trifft einen festen Anteil & Ausfall kostet nur Stichprobengröße \\
        \bottomrule
    \end{tabular}
\end{table}

Das umgesetzte System verwendet den zweiten Weg. Alle acht Knoten starten die
Simulation mit derselben Konfigurationsdatei \texttt{input\_config.json} und
werfen dennoch jeweils andere Strahlen, weil jeder Prozess seinen eigenen
Zufallsstrom besitzt. Der Ansatz ist bewusst schlicht gewählt: Er braucht keine
zentrale Zuteilung und keinen Austausch zwischen den Knoten, sondern nur, dass am
Ende alle Ergebnisse eingesammelt werden.

\subsubsection{Unabhängigkeit der Zufallsströme}
\label{subsubsec:rng}

Die Korrektheit des Replikat-Ansatzes steht und fällt mit einer Bedingung: Die
Zufallszahlen auf den Knoten müssen unabhängig sein. Würde jeder Knoten dieselbe
Folge zufälliger Richtungen erzeugen, so lieferten alle acht exakt dasselbe
Ergebnis, und das Zusammenführen brächte keinen Erkenntnisgewinn, sondern nur
achtfach dieselben Strahlen. Die Rust-Engine schließt das aus, indem sie ihren
Zufallsstrom pro Thread bzw. Prozess frisch aus der Systementropie zieht:

\begin{lstlisting}[style=clusterlst, language=Rust, caption={Initialisierung des Zufallsstroms pro Prozess (\texttt{simulation.rs}).}]
let mut rng = rand::rng();                       // Omnidirectional, random rays
let random_angle: f32 = rng.random_range(0.0..(2.0 * PI));
let direction = Vec2::new(random_angle.cos(), random_angle.sin());
\end{lstlisting}

Weil jeder der acht Prozesse auf seinem Knoten einen anderen Startzustand erhält,
erzeugt jeder eine andere Stichprobe. Acht Knoten mit je einer Million Strahlen
liefern damit acht Millionen \emph{unabhängige} Stichproben derselben Raumakustik.

\subsubsection{Statistische Konvergenz}
\label{subsubsec:konvergenz}

Statistisch verhält sich der verteilte Lauf wie eine einzige, achtfach größere
Simulation. Für einen Monte-Carlo-Schätzer aus $N$ unabhängigen Teilstichproben
gleicher Größe sinkt die Standardabweichung des Mittelwerts mit
\[
    \sigma_N \;=\; \frac{\sigma_1}{\sqrt{N}} .
\]
Acht Knoten reduzieren die Streuung der geschätzten Impulsantwort demnach um den
Faktor $\sqrt{8} \approx 2{,}83$ gegenüber einem einzelnen Knoten.
Abbildung~\ref{fig:mc-convergence} bestätigt dieses Verhalten in einem numerischen
Experiment: Aufgetragen ist der mittlere quadratische Fehler der geschätzten
Impulsantwort über der Knotenzahl. Die Messpunkte folgen sauber der
$1/\sqrt{N}$-Referenzkurve; von einem auf acht Knoten sinkt der Fehler um den
Faktor $2{,}78$, praktisch exakt der theoretisch erwartete Wert.

\begin{figure}[htbp]
    \centering
    \includegraphics[width=0.8\textwidth]{figures/mc_convergence}
    \caption{Mittlerer quadratischer Fehler der geschätzten Impulsantwort über der
    Knotenzahl $N$ (je $5\cdot10^4$ Strahlen, doppelt logarithmisch). Die
    Messpunkte folgen der $1/\sqrt{N}$-Referenz; acht Knoten senken den Fehler um
    den Faktor $2{,}78 \approx \sqrt{8}$.}
    \label{fig:mc-convergence}
\end{figure}

Bemerkenswert ist die Kehrseite dieser Kurve: Um den Fehler zu halbieren, braucht
es die vierfache Knotenzahl. Der Nutzen zusätzlicher Knoten nimmt also ab, was bei
der Wahl der Clustergröße zu berücksichtigen ist. Für die hier angestrebte
Genauigkeit sind acht Knoten ein sinnvoller Kompromiss zwischen Aufwand und
Qualität.

\subsection{Architektur des Clusters}
\label{subsec:cluster-arch}

Das System folgt einem \emph{Manager-Worker}-Muster
(Abbildung~\ref{fig:cluster-arch}). Ein Manager-Knoten koordiniert den Ablauf,
acht Worker-Knoten führen die eigentliche Simulation aus. Der Manager selbst
rechnet nicht mit, sondern startet die Worker, wartet auf sie und führt am Ende
die Ergebnisse zusammen.

\begin{figure}[htbp]
    \centering
    \includegraphics[width=0.82\textwidth]{figures/cluster_arch}
    \caption{Manager-Worker-Architektur. Der Manager startet die acht Knoten
    parallel per \texttt{ssh}, holt nach Abschluss jede Ergebnisdatei per
    \texttt{scp} zurück und führt sie mit \texttt{merge\_results.py} zusammen.}
    \label{fig:cluster-arch}
\end{figure}

Die Kommunikation läuft vollständig über Standard-Werkzeuge des Betriebssystems.
\texttt{ssh} startet die Prozesse aus der Ferne, \texttt{scp} überträgt die
Ergebnisdateien zurück. Ein eigenes Netzwerkprotokoll oder ein
Message-Passing-Framework kommt nicht zum Einsatz. Für einen Stapelbetrieb, bei
dem jeder Knoten genau einmal startet, eine Weile rechnet und ein einziges
Ergebnis zurückliefert, wäre solcher Aufwand nicht gerechtfertigt: Es gibt keine
laufende Zustandssynchronisation und keinen Nachrichtenaustausch während der
Rechnung, den ein Framework verwalten müsste.

Damit das funktioniert, müssen drei Voraussetzungen erfüllt sein. Auf jedem Knoten
liegt dieselbe vorkompilierte Binärdatei \texttt{rust\_service} samt der
Konfiguration \texttt{input\_config.json}. Der Manager kann sich per SSH ohne
interaktive Passworteingabe anmelden, was eine Schlüssel-basierte
Authentifizierung voraussetzt. Und die Knoten sind über feste Adressen im
Labornetz erreichbar.

\subsection{Ablaufsteuerung: \texttt{run\_cluster.sh}}
\label{subsec:cluster-run}

Das Skript \texttt{run\_cluster.sh} steuert den verteilten Lauf in drei Phasen:
paralleles Starten, Warten und Einsammeln.

\begin{lstlisting}[style=clusterlst, language=bash, caption={Verteilter Start und Einsammeln der Ergebnisse (gekürzt).}]
# Phase 1: alle Knoten gleichzeitig starten (Hintergrundprozesse)
ssh node1@139.6.65.29 "./rust_service --config-file input_config.json" &
ssh node2@139.6.65.18 "./rust_service --config-file input_config.json" &
# ... node3 bis node8 analog ...

wait                       # Phase 2: auf Abschluss ALLER Knoten warten

# Phase 3: Ergebnis jedes Knotens zum Manager kopieren
scp node1@139.6.65.29:~/ir_output.json ~/ir_node1.json
scp node2@139.6.65.18:~/ir_output.json ~/ir_node2.json
# ... node3 bis node8 analog ...
\end{lstlisting}

Zwei Eigenschaften der Shell nutzt der Ablauf gezielt. Das kaufmännische Und
(\texttt{\&}) am Zeilenende schickt jeden \texttt{ssh}-Aufruf in den Hintergrund.
Die acht Simulationen laufen dadurch nahezu gleichzeitig an, statt nacheinander
abgearbeitet zu werden. Das darauffolgende \texttt{wait} blockiert, bis
\emph{alle} Hintergrundprozesse beendet sind, und bildet so eine
Synchronisationsbarriere zwischen der Rechen- und der Einsammel-Phase. Erst wenn
der letzte Knoten fertig ist, beginnt \texttt{scp}, von jedem Knoten dessen
\texttt{ir\_output.json} zu holen und unter einem knotenspezifischen Namen auf dem
Manager abzulegen.

Diese Barriere hat eine wichtige Konsequenz für die Laufzeit: Der Gesamtlauf ist
so langsam wie der \emph{langsamste} Knoten. Solange die Hardware homogen ist,
fällt das kaum ins Gewicht; auf ungleichen Knoten würde ein einzelner langsamer
Rechner den ganzen Lauf ausbremsen, während die übrigen sieben bereits fertig
warten. Dieses Straggler-Problem ist der zentrale Grund, warum ein statischer
Replikat-Ansatz mit gleich großen Aufträgen homogene Hardware voraussetzt.

\subsection{Aggregation der Teilergebnisse: \texttt{merge\_results.py}}
\label{subsec:cluster-merge}

Sobald alle Teilergebnisse auf dem Manager liegen, führt \texttt{merge\_results.py}
sie zu einer einzigen Impulsantwort zusammen. Jeder Knoten liefert dasselbe
Ausgabeschema, nämlich Metadaten sowie zwei positionsgleiche Listen aus
Verzögerungen und Drücken. Die Aggregation besteht deshalb im Kern aus einer
schlichten Verkettung dieser Listen:

\begin{lstlisting}[style=clusterlst, language=Python, caption={Zusammenführen der Knoten-Ergebnisse, gekürzt (\texttt{merge\_results.py}).}]
all_delays, all_pressures = [], []
for f in sorted(glob.glob("/home/manager/ir_node*.json")):
    with open(f) as file:
        data = json.load(file)
        all_delays    += data["hits"]["delays_seconds"]
        all_pressures += data["hits"]["pressures"]

merged = {
    "metadata": {"sample_rate": 44100, "rays_cast": 1000000 * 8,
                 "rays_received": len(all_delays), "room_name": "MVP_Test_Room_1"},
    "hits": {"delays_seconds": all_delays, "pressures": all_pressures},
}
json.dump(merged, open("/home/manager/ir_merged.json", "w"), indent=2)
\end{lstlisting}

Dass die Verkettung hier die richtige Operation ist, folgt direkt aus dem
Replikat-Modell. Weil die einzelnen Läufe unabhängige Stichproben desselben Raums
sind, ist die Vereinigung aller Treffer statistisch identisch mit dem Ergebnis
eines einzigen Laufs mit acht Millionen Strahlen. Es findet keine Mittelung und
keine Gewichtung statt, sondern nur ein Aneinanderhängen. Das Feld
\texttt{rays\_cast} ist im Skript fest auf $8 \cdot 10^{6}$ gesetzt; solange der in
Abschnitt~\ref{subsec:cluster-grenzen} beschriebene Kopierfehler bei Knoten~6
besteht, fließen tatsächlich nur sieben Knoten in die Aggregation ein, und dieser
Metadatenwert überschätzt die reale Strahlzahl. \texttt{rays\_received} zählt
dagegen die tatsächlich empfangenen Treffer. Die
zusammengeführte Datei trägt exakt dasselbe Schema wie die Ausgabe eines einzelnen
Knotens und ist damit ohne weitere Anpassung für die DSP-Pipeline verwendbar.

\subsection{Datenfluss im Überblick}
\label{subsec:cluster-datenfluss}

Der gesamte verteilte Ablauf lässt sich in fünf Stufen zusammenfassen:

\begin{enumerate}
    \item Der Manager startet per SSH auf jedem der acht Knoten
    \texttt{rust\_service} mit identischer Konfiguration, alle parallel.
    \item Jeder Knoten simuliert unabhängig eine Million Strahlen und schreibt
    lokal seine \texttt{ir\_output.json}.
    \item \texttt{wait} synchronisiert: Die Einsammel-Phase beginnt erst, wenn der
    letzte Knoten fertig ist.
    \item \texttt{scp} überträgt jede \texttt{ir\_output.json} als
    \texttt{ir\_node$k$.json} zum Manager.
    \item \texttt{merge\_results.py} verkettet alle Trefferlisten zu einer
    \texttt{ir\_merged.json}, die anschließend die DSP-Pipeline verarbeitet.
\end{enumerate}

\subsection{Leistungsbetrachtung}
\label{subsec:cluster-leistung}

\subsubsection{Erwarteter Speedup}
\label{subsubsec:speedup}

Bei der Leistung sind zwei Sichtweisen zu trennen. In der umgesetzten
Replikat-Betriebsart rechnet jeder Knoten dieselbe Strahlzahl, sodass acht Knoten
in nahezu derselben Wanduhrzeit die achtfache Stichprobe liefern. Der Gewinn ist
hier keine kürzere Laufzeit, sondern die höhere Genauigkeit aus
Abbildung~\ref{fig:mc-convergence}; man spricht vom schwachen Skalieren im Sinne
von Gustafson, bei dem mit der Rechnerzahl die bearbeitete Problemgröße wächst,
nicht die Geschwindigkeit einer festen Aufgabe. Will man dagegen eine feste
Gesamtstrahlzahl in kürzerer Zeit erreichen, so verteilt man diese auf die Knoten,
sodass jeder nur $8\cdot10^{6}/k$ Strahlen rechnet. Erst in dieser
Partitionierungs-Betriebsart greift das Gesetz von Amdahl. Bezeichnet $p$ den parallelen
Anteil der Arbeit, hier die reine Strahlverfolgung, und $1-p$ den seriellen Rest
aus Prozessstart, Ergebnistransfer und Aggregation, so beträgt der theoretische
Speedup
\[
    S(k) \;=\; \frac{1}{(1-p) + \dfrac{p}{k}} .
\]
Da die Strahlverfolgung die Laufzeit klar dominiert, liegt $p$ nahe bei eins, und
$S(k)$ steigt bis $k=8$ nahezu linear. In der Praxis bleibt der reale Speedup
darunter, und zwar aus drei Gründen: dem bereits erwähnten Straggler-Effekt (der
Lauf ist so schnell wie der langsamste Knoten), dem Overhead für den
SSH-Verbindungsaufbau sowie der Übertragung und Zusammenführung der
Ergebnisdateien.

\subsubsection{Kommunikations- und I/O-Overhead}
\label{subsubsec:overhead}

Der Overhead dieses Ansatzes ist klein, aber nicht null. Pro Knoten fallen ein
SSH-Verbindungsaufbau und eine \texttt{scp}-Übertragung an. Übertragen wird jeweils
eine JSON-Datei, deren Größe mit der Zahl der empfangenen Treffer wächst; bei
einer Million geworfener Strahlen liegt sie im Bereich einiger Megabyte. Da diese
acht Übertragungen erst \emph{nach} der Rechenphase stattfinden und nicht mit ihr
überlappen, addieren sie sich voll zur Gesamtlaufzeit. Bei deutlich mehr Knoten
oder größeren Ergebnisdateien würde dieser sequenzielle Einsammel-Schritt zum
Flaschenhals; eine Überlappung von Rechnung und Transfer oder ein kompakteres
Ausgabeformat als das gut lesbare, aber platzintensive JSON wären dann die
naheliegenden Gegenmaßnahmen.

\subsubsection{Messung am Cluster}
\label{subsubsec:messung}

Die konkreten Laufzeiten hängen von der Laborhardware ab und sind am Cluster zu
messen. Um einen aussagekräftigen Speedup zu erhalten, wird in der
Partitionierungs-Betriebsart gemessen: In \texttt{input\_config.json} wird
\texttt{rays\_to\_cast} auf $8\cdot10^{6}/k$ gesetzt, sodass die Gesamtarbeit über
alle $k$ Knoten konstant bleibt. Dann wird \texttt{run\_cluster.sh} für
verschiedene Knotenzahlen mit der Shell-eigenen Zeitmessung ausgeführt:

\begin{lstlisting}[style=clusterlst, language=bash, caption={Messvorschrift für die Laufzeit am Cluster.}]
time ./run_cluster.sh        # Wanduhrzeit fuer den kompletten verteilten Lauf
\end{lstlisting}

Tabelle~\ref{tab:speedup} nimmt die Messwerte auf. Als Referenz dient der Lauf mit
einem einzelnen Knoten; der Speedup ergibt sich als Verhältnis der Wanduhrzeiten.

\begin{table}[htbp]
    \centering
    \caption{Wanduhrzeit und Speedup des verteilten Laufs (am Cluster zu messen).}
    \label{tab:speedup}
    \begin{tabular}{r r r}
        \toprule
        Knoten $k$ & Wanduhrzeit [s] & Speedup $S$ \\
        \midrule
        1 & \dots & $1{,}00$ \\
        2 & \dots & \dots \\
        4 & \dots & \dots \\
        8 & \dots & \dots \\
        \bottomrule
    \end{tabular}
\end{table}

\subsection{Fehlertoleranz und Sicherheit}
\label{subsec:cluster-fehler}

Zwei Querschnittsthemen verdienen eigene Betrachtung. Beim \textbf{Sicherheits}aspekt
ist die Schlüssel-basierte SSH-Authentifizierung zentral: Der Manager benötigt
passwortlosen Zugriff auf jeden Knoten, was über hinterlegte öffentliche Schlüssel
gelöst ist. Der private Schlüssel des Managers ist damit das entscheidende
Geheimnis; wer ihn besitzt, kann auf allen Knoten Code ausführen. In einem
abgeschlossenen Labornetz ist das vertretbar, für einen offeneren Betrieb wäre eine
engere Rechtevergabe angebracht.

Beim Thema \textbf{Fehlertoleranz} zeigt der aktuelle Stand seine deutlichste
Schwäche, die im folgenden Abschnitt zusammen mit den übrigen Grenzen diskutiert
wird. Der Kern des Problems: Ein ausgefallener Knoten wird nicht erkannt, sondern
schlicht übergangen, sodass die Stichprobengröße unbemerkt sinkt.

\subsection{Grenzen des aktuellen Standes und Ausblick}
\label{subsec:cluster-grenzen}

Die Verteilungsschicht ist als funktionsfähiger Prototyp umgesetzt. Für einen
robusten Betrieb bestehen mehrere klar benannte Schwachstellen, die zugleich die
nächsten Entwicklungsschritte vorgeben.

\begin{itemize}
    \item \textbf{Fest verdrahtete Konfiguration.} IP-Adressen, Knotennamen, die
    Knotenzahl und die Ablagepfade stehen direkt im Skript. Ein Umzug auf andere
    Hardware oder eine andere Knotenzahl erfordert Änderungen am Quelltext. Eine
    Auslagerung dieser Werte in eine Konfigurationsdatei wäre der naheliegende
    erste Schritt.

    \item \textbf{Kopierfehler bei Knoten~6.} Im Einsammel-Schritt wird das
    Ergebnis von Knoten~6 versehentlich nach \texttt{ir\_node5.json} kopiert und
    überschreibt damit das Ergebnis von Knoten~5. Effektiv fließen so nur sieben
    statt acht unabhängige Stichproben in die Aggregation ein. Der Fehler ist ein
    Tippfehler im Zielnamen und leicht zu beheben, führt aber vor Augen, wie
    fehleranfällig die von Hand ausgeschriebene Knotenliste ist -- ein weiteres
    Argument für die genannte Auslagerung in eine Schleife über eine
    Konfigurationsliste.

    \item \textbf{Keine Fehlererkennung.} Fällt ein Knoten aus oder liefert er kein
    Ergebnis, schlägt der zugehörige \texttt{scp}-Aufruf fehl, der Gesamtlauf setzt
    sich aber fort und aggregiert stillschweigend über die vorhandenen Dateien. Ein
    Knotenausfall bleibt damit unbemerkt und verfälscht die Stichprobengröße. Eine
    Prüfung der Rückgabecodes und eine Warnung bei fehlenden Ergebnissen wären hier
    angebracht.

    \item \textbf{Manuelle Kopplung an die DSP-Pipeline.} Die Aggregation erzeugt
    \texttt{ir\_merged.json}, die Python-DSP-Pipeline liest jedoch
    \texttt{ir\_output.json}. Der Übergang zwischen verteiltem Lauf und Faltung ist
    derzeit ein manueller Schritt. Eine saubere Integration würde den Cluster-Lauf
    direkt aus der Orchestrierung heraus anstoßen und die zusammengeführte Datei
    automatisch weiterreichen.

    \item \textbf{Statische Verteilung.} Der Replikat-Ansatz mit gleich großen
    Aufträgen ist einfach und robust, setzt aber homogene Hardware voraus und lastet
    unterschiedlich schnelle Knoten nicht unterschiedlich aus. Ein Modell mit
    dynamischer Arbeitsverteilung -- ein Manager vergibt Ray-Kontingente, sobald ein
    Knoten frei wird -- würde heterogene Hardware besser nutzen und die Knotenzahl
    flexibel machen.
\end{itemize}

Alle genannten Punkte betreffen Robustheit und Bedienbarkeit, nicht die
grundsätzliche Korrektheit des Verfahrens. Die verteilte Ausführung liefert eine
gültige, statistisch dichtere Impulsantwort, und das Manager-Worker-Muster über
SSH und \texttt{scp} erfüllt seinen Zweck im Laborbetrieb.

\subsection{Zusammenfassung}
\label{subsec:cluster-zusammenfassung}

Dieses Kapitel hat die verteilte Ausführung der Strahlverfolgung auf dem
Labor-Cluster beschrieben. Weil die Simulation embarrassingly parallel ist, genügt
ein Replikat-Ansatz: Alle acht Knoten rechnen denselben Auftrag mit unabhängigen
Zufallsströmen, was statistisch einer achtfach größeren Simulation entspricht und
den Schätzfehler der Impulsantwort um rund $\sqrt{8}$ senkt -- ein Verhalten, das
das numerische Experiment in Abbildung~\ref{fig:mc-convergence} bestätigt. Ein
Shell-Skript startet die Knoten per SSH parallel, synchronisiert über \texttt{wait}
und sammelt die Ergebnisse per \texttt{scp} ein; \texttt{merge\_results.py}
verkettet die Trefferlisten zu einer gemeinsamen Datei im gewohnten Schema. Der
Prototyp erfüllt seine Aufgabe, weist aber in Konfigurierbarkeit, Fehlererkennung
und der Anbindung an die DSP-Pipeline die diskutierten Schwächen auf, die den
Rahmen für die Weiterentwicklung abstecken.
